% !TeX program = lualatex
% Standalone LuaLaTeX manuscript. All vector graphics are encoded after \end{document}.
% Run latexmk -lualatex A_Prescribed_Gap_and_Two_Arithmetic_Selections.tex (or run lualatex twice).
% The companion ZIP supplies the expanded evidence, finite tests and scoped Lean replays.
\documentclass[12pt,letterpaper]{article}
\usepackage{iftex}\RequireLuaTeX
\usepackage[margin=.94in,headheight=15pt,headsep=14pt,footskip=27pt]{geometry}
\usepackage{fontspec}\setmainfont{Latin Modern Roman}\setsansfont{Latin Modern Sans}
\usepackage{amsmath,amssymb,amsthm,mathtools,booktabs,array,tabularx,longtable,graphicx,microtype,xcolor,fancyhdr,enumitem,hyperref,bookmark}
\definecolor{ink}{HTML}{183831}\definecolor{accent}{HTML}{095D58}\definecolor{mute}{HTML}{455E55}
\hypersetup{colorlinks=true,urlcolor=accent,linkcolor=accent,citecolor=accent,pdftitle={A prescribed gap and two arithmetic selections},pdfauthor={Uriel Research}}
\newcolumntype{Y}{>{\raggedright\arraybackslash}X}
\setlength{\parindent}{0pt}\setlength{\parskip}{4pt plus 0.7pt}\setlength{\emergencystretch}{3em}\raggedbottom
\pagestyle{fancy}\fancyhf{}\fancyhead[L]{\small\sffamily\color{mute}Uriel Research}\fancyhead[R]{\small\sffamily\color{mute}Two arithmetic selections}\fancyfoot[C]{\small\sffamily\thepage}
\renewcommand{\headrulewidth}{.3pt}\setcounter{tocdepth}{1}\setlist{itemsep=2pt,topsep=5pt}



\directlua{
  local source = assert(io.open(status.filename, "rb"))
  local figures, retained = {}, {}
  local name
  for line in source:lines() do
    local start = line:match("^BEGIN FIGURE ([a-z]+)")
    if start then name = start; figures[name] = {}
    elseif line:match("^END FIGURE") then name = nil
    elseif name then table.insert(figures[name], line:match("^[0-9a-f]+"))
    end
  end
  source:close()
  function bandFigure(key)
    local hex = table.concat((assert(figures[key], "Missing figure " .. key)))
    local data = hex:gsub("..", function(pair) return string.char(tonumber(pair,16)) end)
    local document, fileid = pdfe.new(data, string.len(data), key)
    retained[key] = document
    local figure = img.scan{filename = fileid or document, pagebox = "crop",
                            attr = "/Group<</S/Transparency/K false/I false>>"}
    node.write(img.node(figure))
  end
}
% Use the same graphicx size calculation as an ordinary vector PDF.
\makeatletter
\newsavebox{\bandplotbox}
\newcommand{\bandfigure}[2]{\begingroup
  \sbox{\bandplotbox}{\directlua{bandFigure("#2")}}%
  \Gin@bboxtrue\def\Gin@llx{0}\def\Gin@lly{0}%
  \Gin@defaultbp\Gin@urx{\wd\bandplotbox}%
  \Gin@defaultbp\Gin@ury{\ht\bandplotbox}%
  \def\Ginclude@band##1{\GPT@setup@scale\Ginclude@@pdftex{\usebox\bandplotbox}}%
  \def\Ginclude@graphics##1{\Gin@setfile{band}{}{##1}}%
  \includegraphics[width=#1]{#2}\endgroup}
\makeatother





\begin{document}
\thispagestyle{plain}
{\sffamily\small\bfseries\color{accent}URIEL RESEARCH\hfill 23 NOVEMBER 2026}\par\vspace{15pt}
{\Huge\bfseries\color{ink}A prescribed gap and two arithmetic selections\par}\vspace{8pt}
{\Large\color{mute}What survives a change of cutoff in the localized prime-sum argument?\par}\vspace{14pt}
\textbf{Abstract.} Fix an interval before asking the arithmetic construction for either of two windows. Can those selections certify every point between their observations? We obtain a conditional affirmative answer on the original local predicates. The auxiliary primes can be held fixed at distinct nearby outer queries. Once the odd base fields and the moving field of this particular prime pair have been fixed, the lists of exceptional weights admit a simultaneous choice at all cutoffs. A weight in the previously inspected range can disappear at a finer cutoff. Its first appearance there is impossible. This one-sided relation suffices to combine an earlier test at the left boundary and a later test at the right, provided the floor at the left exceeds the allowance at the right and the later arithmetic estimates hold throughout.

The endpoint tests themselves are substantial arithmetic inputs. Even an admissible rank-one center for the actual conductor-one slot at 3 can defeat the screening rule on an infinite family of intervals fixed before the selectors, including against every permissible enlargement of the cutoff at those intervals. On that very field a separate argument proves the required sum of cohomological lengths equals zero at all scales. We remove this slot from the original odd-field obligation without changing its rank constants, its pair, or any other slot. A direct enclosure independently certifies the actual corrected zeta sum on the finite target \([8,12]\). The moving field and the remaining odd slots prevent us from claiming an unconditional arithmetic certificate for an arbitrary prescribed interval.

\medskip\textbf{Evidence.} Four primary protocols, exact continuous-interval controls, source and object fingerprints, supplementary Lean files, kernel receipts and an independently replayable numerical certificate accompany this manuscript. The source proof is pinned at \texttt{40e56ba}. The scope of each verification appears alongside the claim.
\clearpage\tableofcontents

\section{Choose the gap before the witnesses}
Suppose we place two marks on a path before anyone inspects it. We want a guarantee about every step between the marks. The construction already knows how to find stretches of good ground arbitrarily far along the path. It can choose the stretch after seeing the obstacles. We have reversed this order. Even if we request two inspections whose permitted ranges cover our marks, the stretches they actually certify may lie elsewhere. To reach the ground between the marks, we must recover the information each inspection used in choosing its stretch.

The relevant information in this problem has an unexpectedly precise shape. At a given cutoff, a finite list singles out integer weights where a routine estimate of a cohomological quantity may be too costly. An ordinate \(x\) tolerates some of these weights. It only uses integers below its own cutoff, and the size of an allowable exception increases with \(x\). On passing to a larger outer query the routine estimate can mark a different list. The question is whether a newly marked weight can overturn a previously successful inspection. One can answer this at the level of the definitions even though the proof gives no numerical values for most of the centers that generate the lists.

\subsection{Four orders of choice}
The variable \(x\) runs through the prescribed closed interval \(I=[a,b]\). Write \(Y\) for the lower scale in a prime-pair request and \(X=Y^J\) for its outer ceiling. The inner start returned by the published selector is \(W\). Its selected window \([W,W^2]\) lies in the ambient range. An inclusion \(I\subseteq[Y,X]\) says where a datum \emph{could} be assembled; an inclusion in the selected window says where a datum is already supplied. They are different promises. We work in the original arithmetic coordinate. For the conductor-one principal observable its previously derived rescaling factor equals one. \href{https://github.com/nousentllc/riemann-hypothesis/pull/18}{[Previous endpoint and rescaling study]}

\begin{enumerate}
\item The existing supply theorem first accepts a sufficiently large outer query and then selects primes and a successful window. Two calls give two existential records. A list of common fields in their types does not provide an identity between their chosen centers.
\item Our prospective construction fixes \([a,b]\), brackets it inside two admissible envelopes, and arranges one pair for both. For fixed character data, the odd tables can be chosen before the pair. The moving table is chosen for that pair. Both choices precede all subsequent truncations of their tables.
\item The positive bridge then requires two endpoint tests on these very tables. It can use a later estimate with its own constants on all of \(I\). It does not require information about each integer independently at every intermediate point.
\item An independent evaluation of the scalar accepts \([a,b]\) first and computes the finite prime-power terms needed there. It requires no choice of primes \((q,d)\), field model, layer or center. Its output belongs to that one finite interval. The proof's localization problem concerns intervals arbitrarily far out.
\end{enumerate}

\subsection{What would falsify the proposed bridge?}
A successful left and right check on coherent records followed by a failed interior check, under all the stated guards, would falsify the projection argument. A pair of legal queries with incompatible prime size bands would falsify unrestricted pair reuse. A fixed center and an infinite preassigned family of failing intervals would falsify eventual screenability from center reuse alone. Finally, a nonzero layer invariant at any positive weight in the claimed conductor-one slot would invalidate the proposed direct bypass. Each assertion has a distinct failure mode and a different verifier.

We chose \([8,12]\) for the explicit gap before evaluation, \([2,4]\) to expose the equality boundary, and a fixed sparse schedule for the asymptotic obstruction. An additional strict-crossover comparison uses \([2,7/2]\). These are arithmetic ordinates. Integers indexing exceptional weights can be much larger. Where we choose a small exponent to make the mechanism visible by hand, we separately repeat the relevant comparison with the source exponent \(401\). None of the lower thresholds returned by classical choice has been evaluated at 8. A genuine field statement at every scale is explicitly identified when available.

\subsection{A hand test for the main inference}
Let \(Z_1\le Z_2\) be outer cutoffs. Denote by \(E_i\) the exceptions recorded at \(Z_i\). Let \(t(x)=\lfloor x^{C_t}\rfloor\) and write \(\eta\) for one slot's allowance exponent. Suppose the earlier cutoff satisfies \(t(a)\le Z_1+1\), so that it inspected every weight strictly below \(t(a)\). Suppose the finer list marks \(m<t(x)\) at an interior point. Since \(t(x)\le t(b)\), the right-hand check gives \(m\le b^\eta\). If \(b^\eta<t(a)\), then \(m<t(a)\). The left cutoff inspected this integer, and the finer list can retain it there only if the earlier list had marked it. The earlier success now improves the bound to \(m\le a^\eta\le x^\eta\). These three lines explain both the power and the limitation of the bridge. It uses membership at two resolutions, the strict floor, and monotonicity. A count of how many integers can be marked cannot replace the first step relating the lists. The later field bound still has to be affordable on the earlier part of the gap.


\section{Can the two selections use the same prime pair?}
\label{sec:gap-pairs}

A fixed target and two independent calls to an existence theorem give little control over identity. The source proves that at every sufficiently large base there is a good pair of auxiliary primes. Even if the same pair would work twice, two calls need not return it. There is, however, a stronger operation available to someone choosing the schedule. One can reserve enough arithmetic information in a single invocation to produce valid witnesses for a continuum of distinct, nearby outer queries. The interval of admissible queries can be determined independently of the pair. This gives a positive answer about sharing an actual arithmetic witness, with two important limitations. The certified motion can be very small, and each query can still select its cohomological subwindow somewhere other than the target interval.

\subsection{Three coordinates and four pieces of prime information}

Let \(M\) be the character modulus, take \(0<\varepsilon<1/52\), and fix an outer exponent \(J\ge1\). Write \(Y=X^{1/J}\) for the canonical base used by the prime selection at outer scale \(X\). We shall often run the parameterization in the other direction and query \(X=Y^J\). The sought scalar or arithmetic ordinate will be denoted \(x\). A fourth number, \(W\), is the start of a subwindow chosen by the cohomological selector after it has seen \(q,d\). Thus \(Y\) and \(X\) constrain which primes and test points are admissible. Neither is an observation of where \(W\) will fall. The source explicitly chooses the pair at \(Y=X^{1/J}\) and then passes its prime-row property into the \(W\) selector. \href{https://github.com/nousentllc/riemann-hypothesis/blob/40e56bab026552e089d7afcfeff08d2defb9c005/src/Development/AutomorphicArithmetic/CharacterBoxes/WinPairs.lean#L26-L35}{[Outer query and order of choice]}

Define the two scale functions
\[
 z(Y)=Y^{\varepsilon/5},\qquad H(Y)=z(Y)(\log Y)^2.
\]
The original \texttt{WinPair} requires an auxiliary prime \(q\equiv3\bmod8\) and a prime \(d\equiv7\bmod16\) satisfying
\[
 \tfrac12 z(Y)<\log q\le z(Y),\qquad
 \tfrac12 H(Y)<\log d\le H(Y).
\]
In addition, \(q\mid d+1\), and every odd prime \(p\le z(Y)\) divides \(d+1\). The upper limit of the small-prime test moves with \(Y\). At a later base, the two upper bounds become easier. Their strict lower bounds become harder. A new prime passing under the small-prime threshold adds a fresh divisibility requirement. A later row therefore need not follow from an earlier one, even when its two continuous bands still overlap. \href{https://github.com/nousentllc/riemann-hypothesis/blob/40e56bab026552e089d7afcfeff08d2defb9c005/src/Development/AutomorphicArithmetic/CharacterBoxes/ExistsPageN.lean#L27-L47}{[Prime row and cohomological selector]}

The stronger \texttt{PairGood} consists of the row, coprimality \( (M,qd)=1\), a collection of twist estimates, and zero-free boxes for the real-character twists. The twist estimates hold at every \(y\in[Y,X]\). They include both \(\lVert g_\chi(y)\rVert\le y^{1/2+\varepsilon/2}\) and a bound for the logarithmic derivative at one. For \(j=0,1\) the box heights are
\[
 B(M,Y,j)=8M\exp(z(Y))\exp(H(Y))^j.
\]
The box excludes a zero when its real part exceeds \(1-1/400\) and its imaginary part has absolute value at most this height. The source's monotonicity result reduces the upper endpoint \(X\) while keeping \(Y\) fixed. It changes neither the row nor either box. Using it to move the base would overlook several requirements at once. \href{https://github.com/nousentllc/riemann-hypothesis/blob/40e56bab026552e089d7afcfeff08d2defb9c005/src/Development/AutomorphicArithmetic/CharacterBoxes/PairGood.lean#L9-L42}{[Twists, heights, existence and upper-end monotonicity]}

What about the many primes in the proof of existence? The source first finds a large good subset of a candidate set for \(q\), then a large subset of a second candidate set for \(d\), conditional on that \(q\). The first ambient set changes its log-size band with \(Y\). The second changes the modulus in \(d+1\), its upper and lower log bounds, and potentially its chosen \(q\). Two subsets of density greater than one half must intersect when measured against the \emph{same} finite ambient set; there is no such conclusion for the two families of sets just described. Their explicit dependence on the scale and on the first prime is present before the proof passes to an existential \texttt{PairGood}. \href{https://github.com/nousentllc/riemann-hypothesis/blob/40e56bab026552e089d7afcfeff08d2defb9c005/src/Development/AutomorphicArithmetic/MeanSquareBounds/PrimeSetTwoBounds.lean#L35-L65}{[Candidate prime sets]}\quad\href{https://github.com/nousentllc/riemann-hypothesis/blob/40e56bab026552e089d7afcfeff08d2defb9c005/src/Development/AutomorphicArithmetic/MeanSquareBounds/PairExists.lean#L27-L43}{[Taking the two successive prime witnesses]}

\subsection{An absolute obstruction to reusing a row}

Suppose \(1<Y\le V\). If \(z(V)\ge2z(Y)\), every admissible first prime at \(V\) differs from every admissible first prime at \(Y\). Otherwise a common \(q\) would obey \(z(V)/2<\log q\le z(Y)\). Likewise if \(H(V)\ge2H(Y)\), the second primes must differ. Equality suffices for exclusion because the lower band is open. These are supplementary theorems about the actual source predicate. They require neither a model of which primes occur nor access to the center of a cohomology module. For \(Y>1\) the ratio in the second test is
\[
 \frac{H(V)}{H(Y)}
  =\left(\frac VY\right)^{\varepsilon/5}
    \left(\frac{\log V}{\log Y}\right)^2.
\]
The extra square of a logarithm can separate the \(d\) bands earlier than the \(q\) bands.

A useful consequence concerns any fixed power spacing \(V=Y^s\), \(s>1\). Then \(z(V)/z(Y)=Y^{(s-1)\varepsilon/5}\) tends to infinity. Beyond a threshold depending on \(s,\varepsilon\), both primes must change. Expressed in the outer variable, \(X_2=X_1^s\) gives the same separation at fixed \(J\). The two ambient ranges can nevertheless have a substantial overlap. For instance, \(1<s<J\) leaves \([V,Y^J]\) available to both. That overlap supplies common \emph{coordinates}; it says nothing about a common prime row there. This is a rigorous obstruction to a particular proof strategy. It does not say the number fields associated with different rows cannot share some abstract exceptional information or satisfy their own bounds separately.

\subsection{Reserving room for two nearby queries in advance}

We now reverse the conclusion for appropriately close bases. For \(u\ge0\) let \(N(u)=\lfloor u\rfloor+1\), the least integer strictly greater than \(u\). At base \(Y>1\) and a prospective second base \(V\ge Y\), perform the following input-only tests:
\begin{align*}
 z(V)&<N(z(Y)),\\
 \exp(z(V)/2)&<N(\exp(z(Y)/2)),\\
 \exp(H(V)/2)&<N(\exp(H(Y)/2)),\\
 B(M,V,j)&<B(2M,Y,j) \qquad (j=0,1),\\
 V^J&<Y^{J+1}.
\end{align*}
Every quantity is determined before \(q,d\) are chosen. An equality \(u=\lfloor u\rfloor\) still leaves a full unit to \(N(u)\); a value just below the next integer can leave very little. The guard deliberately protects against \emph{all} integers, including composites. This conservatism makes it possible to leave both primes unobserved.

At \(V=Y\) the first three strict inequalities hold by the floor formula. The next two hold because \(B(2M,Y,j)=2B(M,Y,j)>B(M,Y,j)\). The last follows from \(Y>1\). Each left side is continuous at the fixed positive \(Y\); the right sides do not move in this continuity argument. Intersect the finitely many open neighborhoods on which the comparisons hold. One obtains \(r>0\) such that they all hold for \(Y\le V<Y+r\). The fact that the floors change discontinuously as \(Y\) varies does not impair this proof. We froze the initial base before increasing \(V\). The supplementary \texttt{guardOpen} verifies the full claim, and the explicit guard is also usable directly at a proposed \(V\) without calculating the best radius.

Once \(Y\) is above the threshold for \texttt{PairEx} at inputs \(2M,\varepsilon,J+1\), choose a pair at that modulus with twist estimates on \([Y,Y^{J+1}]\). The larger modulus is only an auxiliary input to this stronger witness. The desired output for both queries is a pair at the target modulus \(M\). This is a prospective coordinated selection. It does not identify primes returned by earlier independent calls.

First consider the row. Since \(q\) is an integer strictly greater than \(\exp(z(Y)/2)\), it is at least \(N(\exp(z(Y)/2))\). The second test therefore keeps \(q\) strictly above the new lower threshold. The third test does the same for \(d\). Both old upper inequalities imply the new ones by monotonicity for \(Y,V\ge1\). An integer \(p\le z(V)\) cannot exceed \(z(Y)\), because it would then be at least \(N(z(Y))>z(V)\). Thus the original \(d+1\) divisibility conditions suffice without adding a single new prime. Congruences and primality are unchanged. Coprimality with \(2M\) implies coprimality with \(M\). This proves the requisite arithmetic row at the new base for the same pair.

It remains to pay for the analytic estimates. A Dirichlet character modulo \(M\) lifts along \(M\mid 2M\). At integers coprime to \(2M\) the lifted character agrees with the original. The source induction definition tests equality on integers coprime to a specified product. Passing from \(2Mqd\) to \(2(2M)qd\) only restricts which integers are tested, since the old product divides the new one. Hence an induced twist at modulus \(M\) qualifies for the old stronger estimates after lifting the character to modulus \(2M\). An observed \(y\in[V,V^J]\) lies inside \([Y,Y^{J+1}]\) by the first base order and the reserved horizon. The same estimates therefore hold at every such \(y\). For the real twists in the box, \(\rho.N\mid8M\) implies \(\rho.N\mid8(2M)\); the induction test again weakens by divisibility. The height reservation then puts the new box inside the known old box. No assertion that zeros move continuously with the choice of prime is necessary. \href{https://github.com/nousentllc/riemann-hypothesis/blob/40e56bab026552e089d7afcfeff08d2defb9c005/src/Formalism/SharpRestriction/RamificationFrobenius/TotientGaloisClasses/Analytic.lean#L26-L30}{[Meaning of induced twist]}\quad\href{https://github.com/nousentllc/riemann-hypothesis/blob/40e56bab026552e089d7afcfeff08d2defb9c005/src/Development/AutomorphicArithmetic/MeanSquareBounds/pairEx.lean#L112-L130}{[Original auxiliary-modulus prime construction]}

Together these arguments prove the following new implication in the original Lean types. Given \texttt{PairEx}, for every sufficiently large \(Y\) there are \(r>0,q,d\) such that
\[
   \forall V\in[Y,Y+r),\qquad
       \operatorname{PairGood}(M,\varepsilon,V,V^J,q,d).
\]
The proof of \texttt{commonNeighborhood} obtains \(r\) by continuity of the explicit guard \emph{before} invoking the prime existence theorem. Strictly increasing positive bases give strictly increasing outer scales \(V^J\). Thus these are genuinely different canonical outer queries, although they have been designed to fit one stronger choice. The separate statement \texttt{successiveSelections} takes the original cohomological window input at the same character and obtains its actual \(J\). For every \(V\) in a sufficiently late neighborhood we now have both \texttt{PairGood} and some
\[
 V\le W_V,\quad W_V^2\le V^J,\qquad
 \forall u\in[W_V,W_V^2]\quad
      \operatorname{VAt}(\psi,\varepsilon,401,q,d,e(u),u).
\]
Here \(e\) is one fixed page schedule, chosen before either outer query. The existential \(W_V\) occurs inside the quantifier over \(V\). It need settle neither the same subwindow as \(W_Y\) nor membership of a previously named \(x\) in either selected subwindow.

\subsection{Choose the target interval before choosing a pair}

Take an independently prescribed interval \([a,b]\) and choose a sufficiently late \(Y\) satisfying \(Y<a\le b\le Y^J\). Replace the guard radius by \(\min(r,a-Y)>0\). Every allowed \(V\) now satisfies \(V<a\), and \(b\le Y^J\le V^J\). Consequently the entire prescribed interval lies in the ambient domain of every \texttt{PairGood} conclusion. The supplementary theorem \texttt{prescribedAmbient} retains both features simultaneously, the containment and the prime identity. In particular the twist estimates contained in that record hold at every target ordinate for the twists that meet its induction conditions. They are estimates for those auxiliary twists. The \texttt{VAt} and final target-character scalar comparison involve further hypotheses.

This result accommodates a useful class of gaps chosen parametrically. Fix \(1<\lambda<J\), and consider every sufficiently late \([a,a^\lambda]\). Choose any \(s\) strictly between \(\lambda/J\) and \(1\), then choose \(Y=a^s\). As \(a\) tends to infinity, \(Y\) passes the unknown source threshold. Moreover \(Y<a\) and \(a^\lambda\le Y^J\). All choices belong to the input side of the prime existence step. The argument does not return a universal positive lower bound for \(r\), and the second canonical query may have to be extremely close to the first. If instead two specific outer bases were fixed independently in advance, one must actually check the guard or supply a different common-pair argument for them.

This clarifies what is purchased by stabilizing the pair. One may present one fixed target to two legitimate arithmetic calls using the same character, tolerance, page schedule and primes. To conclude field control on that target one must also manage the cutoff-dependent sets and verify their activity tests or some sufficient alternative. To use two already-computed endpoint certificates rather than a joint construction, one must identify the relevant witnesses in those certificates. Neither task is solved by calling the prime-existence theorem a second time. Conversely, failure of the sufficient guard does not prove that retaining a particular known pair is impossible. Its use of every integer in the two lower-band buffers illustrates exactly where it can discard useful information.

\subsection{A finite prime-row experiment with the gap fixed first}

We fixed \([a,b]=[\exp(72),\exp(76)]\), \(\varepsilon=1/72\) and \(J=4\) before searching for any witness. The four canonical bases have \(\log Y\) equal to \(36,\ 36+10^{-320},\ 36+10^{-4},\ 50\). In each case \(Y<a<b<Y^4\). The modest numerical \(J\) is a control parameter and has not been extracted from the cohomological selector. With \(u=\log Y\), both band functions are computable as \(z=e^{u/360}\) and \(H=u^2 e^{u/360}\).

At \(u=36\) the bands are set by \(z=1.1051709180\ldots\) and \(H=1432.3015098260\ldots\). The program chose \(q=3\) and constructed a 1,400-bit prime \(d\) satisfying \(d\equiv7\bmod16\), \(3\mid d+1\), and \(\log d\in[969.7129056033,969.7129056034]\). There are no odd primes below \(z\) on this finite panel, so the small-prime condition holds there vacuously. The remaining \texttt{WinPair} fields, including both strict lower bands, hold for this same pair at the first three bases. The initial prime \(3\) is allowed by the source row type; downstream large-scale host assembly excludes it. These tests verify a concrete instance of the row arithmetic. They do not claim \texttt{PairGood} twist estimates, field estimates, or occupancy of this target by an original host.

\begin{center}
\begin{tabular}{lcccc}
\toprule
Second log base & Target fits & Guard & \(q\) in band & \(d\) in band\\
\midrule
\(36+10^{-320}\)&yes&pass&yes&yes\\
\(36+10^{-4}\)&yes&fail&yes&yes\\
\(50\)&yes&fail&yes&no\\
\bottomrule
\end{tabular}
\end{center}

The middle row isolates the cost of protecting unobserved integers. Raising \(u\) by \(10^{-4}\) moves \(H\) only to \(1432.3098649314\ldots\), far from endangering the observed \(\log d\). Yet \(\exp(H/2)\) already spans more than one integer because the old value is of order \(10^{311}\). The input-only integer buffer fails. At \(u=50\), in contrast, \(H=2872.4910666028\ldots>2H(36)\). The open lower boundary means \emph{every} admissible second prime there must differ from any admissible second prime at 36. The first-prime size parameter is still only \(1.1489964266\ldots\) and continues to admit \(q=3\). The panel thus separates a conservative failed test from a mathematically unavoidable change.

The primality test is reproducible without assuming that a probable prime is prime. Starting from the trial-divided prime \(1031\), each of eight steps produces an integer \(n=kr+1\) with known prime \(r\) and \(r^2>n\). A supplied witness \(t\) obeys \(t^{n-1}\equiv1\bmod n\) and \(\gcd(t^k-1,n)=1\). For any prime factor \(p\mid n\), the order of \(t\) modulo \(p\) divides \(n-1\) but fails to divide \(k=(n-1)/r\), so \(r\) divides this order. Hence \(p\ge r+1>\sqrt n\); a composite \(n\) would have a smaller prime divisor. The full chain consists of integers in the receipt. This arithmetic argument has a different failure mode from the interval calculations for logarithms and exponentials. The latter bracket the inputs to monotone elementary functions, widen each correctly rounded decimal evaluation by one unit in its final position, and make comparisons by exact rational interval arithmetic at 815 significant digits. When a required comparison cannot be decided, the script raises an error. A separate exact sweep at \(\log_2Y=360u,\ 1\le u\le12\) yields 11 first-band equality exclusions among 78 ordered base pairs. These equality cases would have been misclassified by treating both ends of the source band as closed.

\subsection{How small can the new neighborhood be?}

The input guard permits a quantitative diagnosis of the small radius. This estimate is an additional calculation on paper. Let \(t=\log Y>0\), write \(\alpha=\varepsilon/5\) and put \(F(t)=t^2e^{\alpha t}\). If \(\delta=\log(V/Y)>0\) passes the guard for the second prime, then
\[
 e^{F(t+\delta)/2}<\lfloor e^{F(t)/2}\rfloor+1
 \le e^{F(t)/2}+1 .
\]
Taking logarithms and using \(\log(1+s)\le s\) gives
\[
 F(t+\delta)-F(t)<2e^{-F(t)/2}.
\]
Since \(F'(t)=e^{\alpha t}(2t+\alpha t^2)\) is positive and increasing for \(t>0\), integration yields
\[
 0<\delta<\frac{2 e^{-F(t)/2}}{e^{\alpha t}(2t+\alpha t^2)}.
\]
At \(t=36\), this necessary upper bound is below \(2.285\times10^{-313}\). The successful \(10^{-320}\) displacement was measured in the log base. Even at this moderate log base, the guard forbids displacements on the order of \(10^{-300}\). The bound applies to the sufficient input guard. We have already exhibited a valid prime row at a much larger displacement. The gap between the two statements has a concrete origin. The uniform argument supposes that the unknown selected prime might be the very first integer above the lower edge of its band. Measured as an integer, one unit is small compared with a lower edge of order \(10^{311}\).

If an actual row has been supplied, the underlying calculation looks different. Define the two observed slacks by \(s_q=2\log q-z(Y)\) and \(s_d=2\log d-H(Y)\), both positive. On the right of \(Y\), the same primes meet the size tests precisely while \(z(V)-z(Y)<s_q\) and \(H(V)-H(Y)<s_d\). Coprimality and congruences at a fixed modulus remain true. For the remaining row predicate one can admit every newly crossed odd prime if it divides the already known \(d+1\). This is a useful exact test when the pair is numerical. At our test point, the second-prime slack exceeds \(500\), whereas the protective enlargement permitted by the integer guard in \(H\) is less than \(2e^{-716}\). From knowledge of the row alone one still cannot extend a zero-free box to a higher ordinate. The proof using enlarged modulus and extended upper scale supplies exactly that missing analytic reserve in the local persistence theorem.

One might hope to enlarge the guaranteed radius by using the fact that both selected integers are prime and satisfy congruences. For instance, the first possible integer \(d\) above \(e^{H(Y)/2}\) in the residue class \(7\bmod16\) can lie several integer steps beyond the floor used in our guard. Likewise \(q\mid d+1\) and the small-prime product thin the candidate set further. If one knows the actual selected \(q\), the admissible residue classes for \(d\) become much sparser. To turn this observation into a conclusion uniform in \(Y\) and independently of \(q,d\), one needs a proved lower bound on the first available candidate above the moving edge, or a fresh joint selection argument that prefers rows with slack. Neither is part of the present calculation.

This distinction matters for interpreting the count of compatible selections. A continuum of query points can fit within a single conservative neighborhood, even while their physical displacement in log scale is tiny. Thus the cardinality of the supported family is a poor proxy for the span of a useful bridge. A prescribed gap can nevertheless have appreciable width inside the overlapping ambient intervals because their right endpoints are \(Y^J,V^J\), independent of how small \(V-Y\) is. What remains to exploit that width is a field argument or a separate choice of valid hosts. A claim about the length of a cohomologically controlled gap should record the target, the common ambient interval, and the actual field certificate on the target as three distinct pieces of data. The guard and its successful controls settle the first two roles; the field analysis addresses the third.

There is a practical asymmetry in trying to invert the construction. From a later unmodified \texttt{PairGood} one has twist estimates only for ordinates above its own later base. Moving it backward introduces new ordinates below that base; reducing its upper endpoint cannot recover them. Starting from an unmodified earlier certificate has a different defect, since moving to the right raises the box heights and may extend the upper horizon. The local construction starts earlier, anticipates the future horizon and buys extra box height before selecting the primes. This direction of proof explains why testing equality of the final numerical rows, even if those rows were recoverable, would still leave an analytic obligation.


\section{What actually persists when the cutoff changes}

Suppose we keep the character, its page schedule, and the field fixed. An upper bound on the number of exceptional integers says little about which integers a later query will return. A set with one element can change from one query to the next. The original construction, however, contains more information than its cardinality estimate. It obtains the finite sets from a fixed family of points in a normed extension of the $\ell$-adic numbers. Changing the cutoff lengthens the search over integers and simultaneously shrinks the neighborhoods of those points. These two changes have opposite effects. Their interaction gives a transfer principle precise enough for a prescribed interval.

\subsection{Reopening an existential witness}

Let $L_0$ be the base number field of one slot, let $\ell$ be its prime, and write $R$ for the rank budget. The first component of the original predicate \texttt{ABound} gives a normed field $F$ with an isometric normed algebra structure over $\mathbb Q_\ell$, an ultrametric norm, integers $c,r$ with $r\le R$, and a family $s_i\in F$, $i\in\{1,\ldots,r\}$. A primitive $\ell^c$-th root exists in $L_0$. For $n\ge1$, a layer of level $m$ over $L_0$, and any integer $h\ge0$, the tuple supplies the implication
\[
 \bigl(\forall i,\ \ell^{-h}\le\|s_i-n\|\bigr)
 \quad\Longrightarrow\quad
 \operatorname{shaInv}(L_1,\ell,n)
 \le r\bigl(\max(m,j_0)+c+h\bigr),
 \qquad j_0=\lfloor\log_\ell R\rfloor+3.
\]
At $R=0$ the source uses Lean's total natural-number logarithm. The displayed logarithms of positive cutoffs elsewhere in this section have their usual integer-part meaning. The second component of \texttt{ABound} supplies the distinct estimate \texttt{SecondBound}, which compares cohomology on different layers. The new result keeps that second component wherever the original field estimate calls for it. The first component is the relevant place to seek information about repeated selections. The arithmetic source makes the universal quantifiers over the weight, the layer and $h$ come \emph{after} the choice of centers. A new center for each $h$ would make the argument below unavailable. \href{https://github.com/nousentllc/riemann-hypothesis/blob/40e56bab026552e089d7afcfeff08d2defb9c005/src/Development/AutomorphicArithmetic/AutomorphicVarieties/LayeredIsomorphisms/SecondBound.lean#L11-L44}{Pinned definitions of the first, second, even and odd field inputs.}

Write $h(Z)=\lfloor\log_\ell Z\rfloor$ for $Z\ge1$. Fix a complete admissible tuple once and set
\[
 E(Z)=\left\{n\in\{1,\ldots,Z\}:
                \ \exists i,\ \|s_i-n\|<\ell^{-h(Z)}\right\}.
\]
At $Z=0$ this definition gives the empty set. Both conventions at the boundary matter. The inequality defining the ball is strict. An integer whose distance is exactly the radius can use the away estimate. The new declaration \texttt{coherentFirst} identifies the filtered set with the union used in the original proof and proves the original \texttt{FirstBound} for it at every cutoff. For a weight outside $E(Z)$, substitute $h(Z)$ in the away estimate and use $r\le R$. For the cardinality, if $n$ and $n'$ both occupy one ball, ultrametricity implies $\|n-n'\|_\ell<\ell^{-h(Z)}$. When $n\ne n'$, the valuation of their difference is at least $h(Z)+1$. The difference is divisible by $\ell^{h(Z)+1}$, which is strictly larger than $Z$. Two positive integers in $[1,Z]$ cannot have such a nonzero difference. Summing one possible occupant per center bounds $|E(Z)|$ by $r$; duplicate occupants can only improve the bound. The center may lie in a field larger than $\mathbb Q_\ell$. The difference of the two integer labels lies in $\mathbb Q_\ell$ itself, which is why its valuation takes integer values. \href{https://github.com/nousentllc/riemann-hypothesis/blob/40e56bab026552e089d7afcfeff08d2defb9c005/src/Development/AutomorphicArithmetic/AutomorphicVarieties/LayeredIsomorphisms/Near1.lean#L6-L43}{Integer separation in the original field.} \href{https://github.com/nousentllc/riemann-hypothesis/blob/40e56bab026552e089d7afcfeff08d2defb9c005/src/Development/AutomorphicArithmetic/AutomorphicVarieties/LayeredIsomorphisms/Exc.lean#L7-L42}{The actual finite union and cohomological conclusion.}

The difference between an exceptional label and an observed violation is important. Membership grants the argument permission to use its exceptional-case estimate. Even with an explicit center and an exact membership test, it would take a separate computation or theorem to conclude that the ordinary cohomology bound fails at that weight. Conversely, it is possible to reason exactly about inclusion between filters while remaining unable to classify an individual label for the actual number field. The supplement performs that second kind of deduction.

\subsection{One-sided inclusion on the old range}

Choose integer cutoffs $u\le v$. The inequality $h(u)\le h(v)$ implies $\ell^{-h(v)}\le\ell^{-h(u)}$. A label $n\le u$ surviving at $v$ therefore lies inside the earlier ball around the very same center. This proves
\[
             E(v)\cap[1,u]\subseteq E(u).
\]
If $h(u)=h(v)$, the radii are equal and the displayed inclusion is equality. If $h(u)<h(v)$, a label near the boundary can disappear. At either depth the larger search interval can introduce a new label above $u$. In particular there can be an exchange rather than an increase or decrease of the finite set. At $\ell=3$, take the integer center $s=10$. For $u=3$ the strict inequality reduces to congruence to $10$ modulo $9$, and the only candidate is $1$. For $v=10$ it reduces to congruence modulo $27$, and the candidate is $10$. The later set has empty intersection with the old domain; both sets have size one. This example also exhibits what goes wrong with a proposed universal subset relation in either direction.

We can apply the late field estimate throughout the gap while using the early table to verify its active weights. There is a numerical price to this choice. The original \texttt{FirstBound} at cutoff $v$ bounds the invariant by
\[
          R\bigl(\max(m,j_0)+c+h(v)\bigr).
\]
Restriction of the range of $n$ to $n\le u$ preserves this right-hand side. It gives the stipulated earlier right-hand side immediately if $h(u)=h(v)$, as verified in \texttt{restrictAtEqualDepth}. Without equality the late estimate has the larger allowance $R(h(v)-h(u))$. No argument based solely on restricting its universal quantifier subtracts this allowance. Our field proof always pays for the late cutoff in the single ambient estimate from which the interior conclusion is taken. The earlier list contributes the one-sided inclusion, with no attempted transfer of a smaller logarithmic cost.

A stronger eventual assertion has narrower quantitative meaning. For any fixed integer $B$ and the same fixed finite family of centers, \texttt{stabilizes} establishes
\[
   \exists Z_0\ \forall Z\ge Z_0\ \forall n\in[1,B],\qquad
   n\in E(Z)\quad\Longleftrightarrow\quad \exists i,\ s_i=n.
\]
For each of the finitely many pairs $(i,n)$ with $s_i\ne n$, the distance is positive. Choose a power of $\ell^{-1}$ at most this distance. Take an exponent at least as large as all these exponents and then take $Z$ with integer logarithm at least that exponent and $Z\ge B$. Every unequal pair has left its shrinking open ball. Exact equality continues to give distance zero. This proof covers the empty family and families with repeated centers. The quantifier on $B$ precedes $Z_0$. It gives neither an effective precision from $B,R$ alone nor stabilization for all labels below a moving fraction of $Z$.

A one-center example measures the cost of confusing these two statements. Set $s=3+\ell^m$ with $3<\ell^m$. Then the old label $3$ is present for every $3\le Z<\ell^m$, and disappears at $Z=\ell^m$, precisely when strictness at the new radius first matters. Increasing $m$ leaves $R=1$ and the label unchanged. The run checked $m=8,32,128,512$ separately at four primes, using exact integer arithmetic even at the largest cutoffs. There is a second obstruction to increasing precision at will. The field theorem covers $x\ge X^{1/J}$, so a prescribed left endpoint $a$ imposes $X\le a^J$ for positive $J$ and $a$. The existential stabilization cutoff on labels may lie beyond the permitted envelope. Even if usable stabilization were established, testing which centers are exactly equal to integer weights would still require information about those centers. Neither the theorem nor the finite experiment supplies these equalities for a genuine slot.

\section{Transferring two field tests to an entire real interval}

\subsection{The two branches of an active weight}

Fix a physical target $[a,b]$ before querying the arithmetic source. Write $t(x)=\lfloor x^{C_t}\rfloor$. An active weight is an integer $n$ with $1\le n<t(x)$. The source's truncated sum makes this a strict cutoff, so a label equal to $t(a)$ can become active immediately after the left endpoint. For one slot let $\eta\ge0$ be the exponent in the condition $n\le x^\eta$. A left check at cutoff $u$ certifies this inequality for active $n\in E(u)$ at $x=a$. A right check at cutoff $v$ certifies it for active $n\in E(v)$ at $x=b$. Assume
\[
              u\le v,\qquad t(a)\le u+1,\qquad b^\eta<t(a).
\]
Then every active late-table weight obeys the allowance at \emph{every real} $x\in[a,b]$, provided $0\le a$ and $t(x)\le t(b)$. To see this, pick a putative obstruction $n\in E(v)$ with $n<t(x)$. If $n<t(a)$, integrality and $t(a)\le u+1$ imply $n\le u$. The inclusion on the old range puts $n$ in $E(u)$, and the left check gives $n\le a^\eta\le x^\eta$. If $n\ge t(a)$, then $n<t(x)\le t(b)$ makes the right check applicable. It implies $n\le b^\eta<t(a)\le n$, a contradiction. This gives the entire interval without dividing it into a mesh or assuming that exceptional activity itself is monotone.

There are two distinct monotone quantities in this short proof. The cutoff $t(x)$ can introduce new active weights as $x$ increases. The allowance $x^\eta$ can make those same weights harmless farther to the right. An undetected intermediate obstruction is exactly what becomes possible if the time interval between entry and harmlessness lies between the endpoint tests. The crossover excludes this event for newly active late weights. The old-range inclusion excludes a different mechanism in which an already active integer was absent from the independently supplied left list. Each condition has its own negative control below. The terminal theorems use the exponents $\eta_A=w_A\epsilon/C_K$ and $\eta_B=w_B\epsilon/C_K$ separately. Applying the argument at different cutoffs to each of the finitely many slots preserves the original rank accounting. The actual outer scales $U\le X$ give $u=\lfloor U^{C_t}\rfloor$ and $v=\lfloor X^{C_t}\rfloor$. In the application $a\le U$ implies the required early truncation coverage.

\subsection{Which fields can share the table?}

The even and odd sides have different quantifier orders in the original arithmetic. For the even side, fix an abelian-field datum $M$. The source supplies $R_0(M)$ before the primes $q,d$. Once those primes satisfy \texttt{PrimePair}, it chooses a base field $L_0$, a proof of \texttt{TwoBase}, and a proof of \texttt{ABound} over that field. The concrete construction adjoins $i$ and uses the two quadratic data indexed by $q$ and $d$. The source computes the modulus of that datum as $4 M.m\,q d$. A rank uniform in all pairs therefore does not, by itself, identify their even coefficient fields, obstruction centers or lists. Keeping one pair makes it legal to take one even base and retain one tuple of centers before varying the cutoff. Two independent existential invocations at the same parameters also carry no theorem of pointwise equality between their possible choices. We construct a simultaneous witness by making the choice once. \href{https://github.com/nousentllc/riemann-hypothesis/blob/40e56bab026552e089d7afcfeff08d2defb9c005/src/Development/AutomorphicArithmetic/AutomorphicVarieties/QuadraticExtensions/ModulusAdjoinISqrtNegPrimes.lean#L16-L32}{Even base and its modulus.} \href{https://github.com/nousentllc/riemann-hypothesis/blob/40e56bab026552e089d7afcfeff08d2defb9c005/src/Development/AutomorphicArithmetic/AutomorphicVarieties/QuadraticExtensions/IsAlgClosureRat.lean#L87-L141}{Where the uniform rank enters both field constructors.}

On the odd side a fixed character $\psi$ determines a finite set $B_\chi$ and a field datum at each $\ell\in B_\chi$, namely $L_\chi(\ell)=\operatorname{cyc}(\ell\psi.N)$. Neither depends on $q,d$. For the conductor-one specialization, the datum at the slot $3$ is therefore $\operatorname{cyc}(3)$, represented by $\mathbb Q(\zeta_3)$. The source uses a model of each datum to obtain a bound. Other number-field realizations enter through a base-field isomorphism and an induced isomorphism of the chosen layers, which transport the cohomology estimate. This operation permits a single reference table to certify bounds for all such target realizations. It does not produce explicit images of centers in arbitrary coefficient fields, or an equality between tables selected independently there. The new theorem chooses one table per odd slot before either the pair or the outer cutoff. This is stronger than merely keeping the odd objects fixed after the primes have been chosen. \href{https://github.com/nousentllc/riemann-hypothesis/blob/40e56bab026552e089d7afcfeff08d2defb9c005/src/Development/AutomorphicArithmetic/AutomorphicVarieties/PrimeSetBoundDataOfPrimitiveChar.lean#L69-L77}{The odd datum and the common ambient exponent.} \href{https://github.com/nousentllc/riemann-hypothesis/blob/40e56bab026552e089d7afcfeff08d2defb9c005/src/Development/AutomorphicArithmetic/AutomorphicVarieties/LayeredIsomorphisms/PartB.lean#L50-L94}{Odd-slot isomorphisms and their actual use.}

The construction in \texttt{coherentJoint} keeps the same arithmetic inputs as the corresponding field branch, namely a constant satisfying \texttt{LemmaK}, \texttt{FixedIn} and \texttt{TwoIn}. For $w_A<11/20$ and $w_B<4/5$ it first chooses a common constant package. It obtains an eventual lower threshold before the pair is revealed. Given an admitted page and truncation schedule, it fixes the odd tables, then receives the pair, then fixes the even table, then accepts any sufficiently large outer query satisfying the logarithmic pair bound. This order prevents the late outer scale from silently changing the instruction used for the left test. Uniformity of the threshold comes from the original discriminant, height and layer bounds: the even branch has bounds uniform in the eventual pair under the logarithmic condition; the odd branch has finitely many fixed slots and takes a common majorant of their thresholds. The original proof uses these ingredients to find a good subwindow after making the finite union. The strengthened construction keeps enough of their dependence to test an interval specified in advance. \href{https://github.com/nousentllc/riemann-hypothesis/blob/40e56bab026552e089d7afcfeff08d2defb9c005/src/Development/AutomorphicArithmetic/AutomorphicVarieties/LayeredIsomorphisms/ThmBabOf.lean#L11-L97}{Original construction of the joint field bounds.}

\subsection{What finite calculations distinguish}

For integral centers in $\mathbb Q_\ell$, membership in the ball has a direct algorithm. It is the congruence $s\equiv n\pmod{\ell^{h(Z)+1}}$. Because $\ell^{h(Z)+1}>Z$, the least nonnegative representative is the only possible occupant; it is admitted exactly when between $1$ and $Z$. Enumerating all residues at the largest precision enumerates every one-center profile on the finite cutoff range for integral centers in $\mathbb Q_\ell$. The field proof allows larger normed extensions and arbitrary finite families of centers. Exact coverage of the finite model and the field-independent Lean argument are distinct results.

The first census used $1\le u\le v\le32$. At $\ell=2,3,5,7$ respectively the largest residue moduli were $64,81,125,49$. Each residue was tested at all $528$ ordered cutoff pairs. There were $168{,}432$ comparisons, zero failures of old-range inclusion, and zero differences on the old domain at equal depth. In $13{,}656$ comparisons an old occupant disappeared and a new occupant appeared. The repeated transitions count cutoff pairs and residue profiles, not independent arithmetic observations. The many successful replacements provide a check that the implementation actually changed both the search range and the precision. A code path that always returned the empty list would satisfy both positive inclusion checks and fail this control.

The second census used $t(x)=\lfloor x^2\rfloor$, $\eta=1$, and outer scales $U=4$, $X=8$, hence integer cutoffs $16$ and $64$. These numbers make the membership predicates and the interior failure check exact in rational or integer arithmetic. They are deliberately small model parameters. For a candidate $n$, failure at a real $x\in[a,b]$ requires $x<n$ and $n+1\le x^2$. With $a\ge0$ this is possible exactly when $a<n$, $n+1\le b^2$ and $n+1<n^2$, for $a\le b$. To verify sufficiency take $x=\max(a,\sqrt{n+1})$. Thus zero reported holes means zero in the real continuum of this particular model, as opposed to zero on a chosen numerical mesh.

On $[2,7/2]$, $b^\eta<t(a)=4$. Of the $677$ shared residue profiles across the four primes, $626$ passed both endpoint screens, and none of those $626$ had a late-table failure in the interior. Replacing the left and right centers independently while leaving the gap, endpoint tests and individual rank caps intact yielded $658$ hidden failures. A hand example uses $\ell=2$, the early center $7$, and the late center $3$. At $x=2$, the early center contributes no active violation. At $x=7/2$, the late label $3$ is allowed. But at $x=5/2$ it is active and $3>5/2$. The early table misses a genuinely old label of the late table. Independent selection has destroyed the premise that would have transferred it.

The preselected boundary case $[2,4]$ instead has $b^\eta=t(a)$. One shared-center profile at each prime passes both tests and fails in the interior. Set $s=4$. Its label sits exactly at the inactive left boundary, becomes active at $x=5/2$ when $\lfloor x^2\rfloor=6$, and is already affordable at $x=4$. Thus a weak crossover would fail even with the very same integer center and the same exceptional set $\{4\}$ at both cutoffs. These two controls break different steps in the proof. Independent-center failure is a test of identity across choices. Equality failure is a test of the strict endpoint geometry. Neither supplies a counterexample to a cohomology bound of the pinned formalization.

\subsection{The field conclusion and its unresolved input}

The intermediate terminal statement \texttt{gapVAt} gives the source predicate \texttt{VAt} at each $x\in[a,b]$ with the same pair and the original page function evaluated at $x$. This conclusion quantifies over every eligible number field realizing that character, page and pair. The construction calls the original field and height estimates. It assumes \texttt{PointHV} and \texttt{ThmBc} as well as the shared field interface established above. The numerical guards include the ambient interval $X^{1/J}\le a\le b\le X$, the lower bounds on both primes and auxiliary indices, and a bound on $\log(qd)$. They refer to the later query, whose list is actually used throughout. The early outer scale's role is to place the left truncation inside its valid integer range; a separate early prime estimate cannot substitute for the later guards. On a sufficiently late canonical query, the original lower-growth and logarithmic estimates for \texttt{WinPair} supply those guards. The preassigned two endpoints still have to satisfy the field screens. \href{https://github.com/nousentllc/riemann-hypothesis/blob/40e56bab026552e089d7afcfeff08d2defb9c005/src/Development/AutomorphicArithmetic/AutomorphicVarieties/CMFieldExponentialBounds/AdmissibleOfPageN.lean#L38-L94}{Canonical prime guards.}

The outstanding arithmetic information can be expressed in finite terms. For a fixed choice of tuple and finite cutoffs it suffices to certify that the balls of its centers miss the forbidden active integers at the two endpoints. A certified containing set for the possible residues would also suffice if that set missed those integers. To refute a proposed endpoint screen one would need a genuine center and an actual forbidden weight meeting the strict distance inequality, or a theorem forcing such proximity. The original cardinality estimate establishes the existence of some favorable subwindow by avoiding finitely many weights. It gives no decision procedure for the two screens named by the user. The new result locates precisely what such a procedure would purchase: one simultaneous bound on a whole prescribed continuum, expressed using the original arithmetic objects.

\subsection{Publishing the local arithmetic conclusion in the source's physical variables}

For the scale multiplier $k=1$, hence $c=1$, there is no change of physical coordinate between $x$ and the argument of the field estimate. Nevertheless a field estimate by itself leaves five further fields of the physical witness to supply. The terminal declaration \texttt{gapWinDatum} fixes $C_t=401$, uses the later canonical mark $Y=X^{1/J}$, and keeps the very same schedule $e(x)$ as the input \texttt{PageN}. If $x\in[a,b]$, the ambient assumptions give $Y\le a\le x$ and $x\le b\le X$, so $x^{1/J}\le Y$. A \texttt{PairGood} at $(Y,X)$ restricts to $(Y,x)$ by shrinking the upper end of its twist-bound interval. Its analytic zero boxes continue to refer to the same mark. The two page inequalities follow directly from \texttt{PageN} with multiplier one. For the last field, require $x\ge1$. Since $x^{400}\le x^{401}$, monotonicity of the truncated volume estimate lowers the cutoff from $\lfloor x^{401}\rfloor$ to $\lfloor x^{400}\rfloor$ at exactly the allowed square-root exponential bound. These six obligations reconstruct the source's \texttt{WinDatum} for every $x$ in the prescribed interval. In particular, the conductor-one specialization has a common mark and pair and can still change its page when $x$ crosses a page boundary. \href{https://github.com/nousentllc/riemann-hypothesis/blob/40e56bab026552e089d7afcfeff08d2defb9c005/src/Development/AutomorphicArithmetic/CharacterBoxes/TwistForm.lean#L19-L26}{The six fields of the physical witness.}

The same end-to-end argument accepts a direct certificate on selected odd slots. Fix a finite set $D$ of indices and require an actual pointwise \texttt{BbAt} estimate for $D$ at the same constants, truncation schedule and character datum. In forming the full odd bound on $B_\chi$, each required index is inspected individually. A slot in $D$ receives its estimate from the direct certificate. For each other slot, the two endpoint screens are transferred through its coherent table. This disjunction is implemented where the individual layer estimates are assembled, before the field and height conclusions consume the full \texttt{BbAt}. Consequently the terminal statement only requests the two odd screens for $\ell\in B_\chi\setminus D$. The original rank function and $J$ can still include the budget for an index whose estimate is supplied directly. Improving the exponent by erasing that rank from its formula would be a separate construction. At conductor one the direct zero-invariant estimate for the $3$-slot gives precisely the input for $D=\{3\}$. A proof that the unknown centers there happened to avoid both endpoints would be superfluous to this route.


\section{A fixed gap with two independently labeled witnesses}
\label{sec:twoscreens}
Fix the closed interval \(I=[a,b]\) before either arithmetic selection is made. Suppose the first selection supplies data usable at the left end, and the second supplies data usable at the right end. Grant, provisionally, that both data sets control the same target whenever their own hypotheses hold. The remaining question is one of coverage. If each point can use either witness, how much information has to cross from one selection to the other? Testing this weaker question avoids making equality of the two exception lists a prerequisite for every conceivable proof. It also identifies precisely which additional information would make a positive result applicable to an actual pair of selections.

In the original field-window argument the moving-field list is combined with a finite union of lists indexed by auxiliary primes. The bound on the cardinality of this union is used to select an interval. On that selected interval the proof restricts the global good predicate back to the component lists. This is a particular arrangement of quantifiers. Different outer selections can carry different lists, and the estimates promised on one family of lists also have an outer validity range. In the following calculation we expose the lists and their ranges as independent inputs. The previous experiment isolated the smaller test needed for the actual weights in a sum, which are positive integers \emph{strictly} below its integer cutoff. We use that test throughout. \href{https://github.com/nousentllc/riemann-hypothesis/blob/40e56bab026552e089d7afcfeff08d2defb9c005/src/Development/AutomorphicArithmetic/AutomorphicVarieties/LayeredIsomorphisms/ThmBabOf.lean#L77-L103}{[Original choice of union, window and component estimates]}\quad\href{https://github.com/nousentllc/riemann-hypothesis/blob/d79bd956c6fb2b8a3de51b606765be81d24fb1fc/research/uriel/experiments/cofinal-endpoint-certificates/evidence/collar/FieldEndpoints.lean#L363-L391}{[Prior one-witness active-weight implication]}

\subsection{How the direction of a failed test is determined}
For cutoff exponent \(C>0\), allowance exponent \(h>0\), and one finite set \(E\) of positive integers, write \(t(x)=\lfloor x^C\rfloor\) and take \(x\ge1\). Let \(T_E(x)\) mean that every \(m\in E\) with \(1\le m<t(x)\) satisfies \(m\le x^h\). For one integer \(m\), the failure set in the positive real axis is exactly
\begin{equation}
 B_m=[u_m,v_m),\qquad
 u_m=(m+1)^{1/C},\qquad v_m=m^{1/h},
 \label{eq:stitch-band}
\end{equation}
interpreted as empty if \(u_m\ge v_m\). Indeed \(m<t(x)\) is equivalent to \(m+1\le x^C\), while \(m>x^h\) requires \(x<m^{1/h}\). The left end is included and the right end excluded. Discarding either convention would change some of the exact certificates below. The failed part of a whole list is the union of its individual bands intersected with the prescribed \(I\).

Assume first that the source-scale crossover holds,
\begin{equation}
                        b^h<t(a).
 \label{eq:stitch-cross}
\end{equation}
Suppose \(T_E(a)\). If a label first obstructs \(E\) at a later \(x\), it must satisfy \(m\ge t(a)\). An earlier active label was already small at \(a\), and its allowance can only increase. A label entering at or above \(t(a)\) exceeds even \(b^h\), so it persists as an obstruction through the right end, once it enters at all. Thus a left-passing list has a prefix of usable points. If it ever fails, its first failure is an included point. Conversely, if a second list \(F\) satisfies \(T_F(b)\), every label actually harmful earlier is already active at \(a\). To see this, take a label active at \(x\) between \(a\) and \(b\). It is active at \(b\) as well, and its right-end test gives \(m\le b^h<t(a)\). Hence the second list can only clear obstructions as \(x\) increases. It has a suffix of usable points whose first element belongs to it. These two monotonicity statements rely on successful tests at their designated ends and on the crossover.

The supplementary Lean declarations \texttt{leftFailure} and \texttt{rightFailure} express the two equivalences for an integer cutoff and an already evaluated real allowance. Substituting \(s=x^h\) gives precisely the real inequality used by the earlier active-weight interface. The derivation uses monotonicity of the appropriate cutoffs and allowances. No primality, lattice construction or upper bound on the number of labels enters these equivalences. For a finite family of tagged lists with different allowance exponents, apply the argument to each tag with its own \(h\). If every left tag passes at \(a\) and satisfies its crossover, the left family first fails at the earliest activation among all its tags. A right family becomes usable at the latest clearance among its tags. Tags preserve essential provenance when a numerically equal integer has different allowances on the two branches.

\subsection{The two numbers needed for an exact handoff}
Let \(E\) be the left-passing list and \(F\) the right-passing list. If \(E\) has a harmful label on \(I\), let \(L\) be its earliest activation. Otherwise set \(L=+\infty\). If \(F\) starts with a harmful label let \(R\) be the last clearance of those initially harmful labels, and otherwise set \(R=a\). The usable sets are
\[
             I\cap(-\infty,L)\quad\text{and}\quad I\cap[R,\infty).
\]
The weight screen with free pointwise choice covers all of \(I\) exactly when \(R\le L\). If \(L<R\), its entire missing component is \([L,R)\). The left witness can supply strictly earlier points, and the right witness supplies \(R\) and later points. Two passing exterior tests guarantee only that \(L>a\) and \(R\le b\).

There is a short certificate even when the lists are large. In the nontrivial case \(s=\min\{m\in E:t(a)\le m<t(b)\}\) is the first persistent label to activate. On the right, \(r\) is the greatest label in \(F\) that is harmful at \(a\). The criterion takes the form
\begin{equation}
                 r^{C/h}\le s+1 .
                 \label{eq:stitch-extrema}
\end{equation}
A negative answer has an especially small witness. An entering left label \(s\) and an initially harmful right label \(r\) with \(r^{C/h}>s+1\) exhibit the uncovered point \(x=(s+1)^{1/C}\). It is unnecessary to prove that this pair is extremal to refute coverage. To prove a positive answer from two reported numbers, one needs bounds that hold for \emph{all} relevant labels. Knowing some common integer, equal list cardinalities, or bounds on just two arbitrarily selected labels leaves open the position of a worse pair.

One can ask for less switching but more transfer of information. If all labels of \(E\) that could cause failure at \(b\) belong to \(F\), then \(T_F(b)\) certifies \(T_E(b)\). The earlier one-witness result now fills \(I\) with \(E\) alone, subject to its validity range. Precisely, the membership check is needed only on \(E\cap\{m:1\le m<t(b),\ m>b^h\}\). Since \(F\) passes at \(b\), this intersection must be empty whenever it is contained in \(F\). The symmetric check for \(F\) uses left-end dangers. This criterion is strictly weaker than comparing complete lists. The two-witness criterion can be weaker still. Every list can have a harmful label and the two harmful regions can miss one another. Requiring that all exceptional labels agree across two outer fields can therefore overstate the arithmetic task, even though arbitrary different lists do not resolve it.

A host or field may be usable only over a limited range. Suppose the first is available on \([a,c_L]\) and the second on \([c_R,b]\), and include all other pointwise hypotheses in that availability assumption. In the two-arm model the effective left stop is \(D=\min(L,c_L)\) and the effective right start is \(H=\max(R,c_R)\). Provided the right start is included, the test becomes \(H\le D\). A left end truncated by its closed host range can still be available at \(D\). If it is truncated by a harmful activation it fails there. Equality nevertheless suffices in either case because the right witness covers \(H\). The supplementary \texttt{clippedIff} checks the finite-frontier statement for real cut points. An ineffective left boundary can be placed beyond the compact target. A single point can be sufficient for a patch even when the scalar argument and the underlying field argument provide no continuity of the actual witnesses.

\subsection{A closed interval fixed before enumeration}
We froze \([8,12]\), \(C=2\), \(h=1\) and four pairs of outer availability ranges before enumerating any list pair. Then \(t(8)=64>12\). The chosen menu of fourteen possible integers is
\[
  4,\ 8,\ 9,\ 10,\ 11,\ 12,\ 32,\ 63,\ 64,\ 80,\ 99,\ 120,\ 143,\ 144 .
\]
The two extremes 143 and 144 distinguish activation at the right endpoint from activation outside the interval. The middle labels 32 and 63 already obstruct the full interval and cannot occur in a list passing its designated end. The labels 4 and 8 are harmless throughout. We enumerated every subset with at most two elements, including the empty list, giving 106 possible lists and 11,236 ordered pairs before endpoint screening. The restriction to two entries makes hand certificates transparent. The general proof of \eqref{eq:stitch-extrema} applies to any finite size.

The arithmetic is exact over the continuum. Transform to the coordinate \(z=x^2\). A weight obstructs exactly when \(m+1\le z<m^2\). The program sorts the two boundary events per label, the domain cut points and both endpoints, discarding events outside \([64,144]\). Each event is tested separately from neighboring open cells. One exact rational value between consecutive events suffices for that open cell because all Boolean inequalities there have constant truth value. Coincidences leave only eleven event or open-cell representatives across this deliberately structured menu. A dense grid with eleven ordinary probes would have no corresponding guarantee. The second calculation predicts the answer by the two extreme integers without evaluating the interior. An assertion compares these independent decisions at every applicable ordered pair.

Of 37 lists passing on the left and 29 passing on the right, the product contains 1,073 endpoint-eligible ordered pairs. There are 658 complete patches and 415 holes. In 413 pairs a single list suffices on the full interval. In another 245 one must use both. Of the successful patches, 116 have no point at which both lists pass simultaneously. In 305 failing pairs the two lists have the same cardinality; in 30 they even have a common label. Conversely 215 successful patches require both lists and their two lists are disjoint. The 49 pairs with equal restrictions to the potentially harmful menu labels include just seven pairs with identical complete lists. These are counts of exact cases under a fixed enumeration rule. They supply no sampling distribution for arithmetic fields or experimental estimate of the probability of success at a future scale.

\begin{center}\begin{tabular}{rrrrr}
\toprule
\(c_L\)&\(c_R\)&Covered&Lost to range&Covered without shared point\\
\midrule
12&8 &658&0&116\\
10&10&352&306&96\\
11&9 &525&133&137\\
9&11 &0&658&0\\
\bottomrule
\end{tabular}\end{center}
The column on loss to range compares each row with both witnesses available throughout \([8,12]\). In the second row, a single handoff location has been prescribed in addition to the gap itself. Of the 658 unconstrained patches, 306 cannot hand off at 10. In the third row the two outer domains overlap, but 133 formerly valid patches still lose coverage once the domains are imposed. Geometric overlap of the outer domains says nothing about which witness is usable within their overlap. Conversely, the number of covered pairs with no common good point increases from 116 to 137 in that row. Limiting an availability domain can remove a common point while leaving a one-sided patch intact. The last row contains an open segment supported by neither witness, which defeats even two empty exception lists.

\subsection{Controls that can be checked without software}
For a failure take \(E=\{80\}\) and \(F=\{11\}\). The left obstruction begins at \(9\) because \(9^2=81\). The right obstruction ends at 11. Both exterior tests pass, but \([9,11)\) is uncovered. For instance at \(9\) the cutoff is 81, and both 80 and 11 are active and exceed the allowance 9. Appending 4 to both lists makes them intersect, preserves equal cardinalities and changes no conclusion. If an argument used the intersection instead of the union of the available witness tests, the harmless label 4 would produce a spurious certificate. Testing the union of all labels against one common allowance can also be unduly strong. With full domains and the same exponent, that test requires both witnesses to pass at every point, while the desired scalar conclusion can use either.

For a positive patch with a nonempty overlap take \(E=\{120\}\) and \(F=\{10\}\). The left test succeeds up to 11, the right from 10 onwards, and both succeed on \([10,11)\). Neither is adequate on the entire prescribed interval. The inequality in \eqref{eq:stitch-extrema} reads \(100\le121\). For a tight patch take \(E=\{99\}\), \(F=\{10\}\). At \(x=10\) the left test has already failed and the right test succeeds for the first time. Just to the left only the first test works, and from 10 to 12 only the second does. The condition reads \(100=100\). Demanding a point with two simultaneously valid tests would reject this complete patch. At the far boundary the same phenomenon occurs for \(E=\{143\}\), \(F=\{12\}\). The first test supplies the entire half-open interval \([8,12)\), while the second is usable only at the endpoint 12. A roundoff error that moved the final probe into the interior could lose a real obligation.

The examples also explain why the actual provenance of each endpoint matters. The first positive patch admits a whole interval of handoff values. The tight patch gives no freedom at all. A purported certificate could quote small errors for the two endpoint sums yet omit the fact that the two underlying families have different thresholds or different scalar observables. At a handoff one needs each usable witness to furnish the same desired pointwise conclusion on its own portion. There is no mathematical reason to identify two short-lived fields merely because they prove bounds for one scalar. There is equally no reason to transport an estimate from one scalar to another merely because the lists overlap. In particular these explicit finite lists describe the screen geometry and are not measured exceptional weights for either original field.

\begin{figure}[htbp]\centering\bandfigure{.98\linewidth}{handoff}\caption{Weights on \([8,12]\) with \(C_t=2\) and allowance exponent one. Top: the two integral 3-adic centers 8 and 169 give \(E(64)=\{7,8\}\), \(E(144)=\{8\}\). Middle: label 99 disables the left screen at 10 just as the right screen for 10 starts passing. Bottom: labels 80 and 11 leave the interval \([9,11)\) uncovered. Open and closed symbols distinguish the boundary tests. These finite choices illustrate the screening rule.}\end{figure}

\subsection{The exact scope of the crossover}
The strict bound \eqref{eq:stitch-cross} is easy to apply uniformly. There is a sharper diagnosis if one knows the candidate label universe \(U\). Call \(m\in U\) transient on the target when its failure band begins strictly after \(a\), is nonempty, and ends at or before \(b\). The equivalent arithmetic conditions are
\[
       \lfloor a^C\rfloor\le m\le\lfloor b^h\rfloor,
       \qquad m+1<m^{C/h}.
\]
If \(U\) has no such label then every sublist passing at the left has only persistent failures, and every sublist passing at the right has only initial failures. This follows immediately by intersecting \eqref{eq:stitch-band} with the target. If \(U\) contains a transient label, its singleton passes at both endpoints and fails in the interior. Hence absence of transients is necessary and sufficient for this orientation property uniformly over all sublists of \(U\). A particular list can mask a transient behind other failures. Existence of a transient in an ambient universe therefore gives a counterexample to a uniform claim and does not by itself classify any specified pair.

This refinement was derived after the baseline run and recorded before its own executable tests. A second program used three exponent pairs \((C,h)=(2,1),(3,2),(4,3)\), every two distinct endpoints of a ten-element rational menu, and the integers 1 through 128. In the power coordinate \(z=x^{Ch}\) the two changes for a single label occur at \((m+1)^h\) and \(m^C\). All comparisons, including equality, use rational arithmetic. Across 135 windows the strict crossover holds in 22, while 33 have no transient. Eleven additional windows therefore enjoy orientation for this fixed universe. All 34,560 window/subset cases on the eight-label subuniverse agree with the singleton prediction. For \(C=3,h=2\) on \([5/4,3/2]\), the integer 2 lies between the two cutoff numbers, but its potential band is empty since \(3^2>2^3\). A bare comparison between the cutoffs would overlook this second way for an obstruction to be excluded.

An adversarial long interval shows the price of dropping orientation. Fix \([2,256]\) with \(C=2,h=1\) and take \(E=\{9,6888\}\), \(F=\{3,82\}\). The failures on the left are \([\sqrt{10},9)\) and \([83,256]\), while the failures on the right are \([2,3)\) and \([\sqrt{83},82)\), all relative to the target. These parts are disjoint across the two witnesses. The complete interval is therefore coverable, but at \(2,4,10,100\) the only possible witnesses are respectively left, right, left, right. Any pointwise policy requires at least three switches. Three suffice. Use left until \(\sqrt{10}\), right until \(9\), left until \(83\) and right thereafter. A minimum-switch dynamic program, evaluated on the ordered events and open cells, independently returns three. This example identifies what the source-scale crossover buys: one can represent the entire two-list coverage problem by one handoff and two extremal numbers. Without it even full pointwise coverage need not admit a one-handoff certificate.

For the original fixed-gap question these calculations identify a finite certificate whose implications are exact. Obtain the two witnesses in the right order relative to the target, transport their actual availability domains into a common coordinate, and bound the earliest persistent activation of the first and the latest clearance of the second. A constructive upper bound on the second and lower bound on the first can suffice even when listing all weights is infeasible. A uniform cardinality estimate on each list alone supplies neither. The affirmative and negative examples establish these information requirements at the interface. Whether two specific outer selections satisfy them depends on the original number-field witnesses and the estimates that survive their reindexing.

\subsection{What an unmatched pair still contributes to the scalar problem}
Suppose the lists can actually be partially identified, but the activation and clearance numbers have interval-valued bounds. If \(L\in[L_-,L_+]\) and \(R\in[R_-,R_+]\), complete coverage is already certified when \(R_+\le L_-\). A positive missing interval is forced when \(L_+<R_-\). Between these two decisions the data leave the comparison unresolved. In the positive case only an upper bound for the clearing label and a lower bound for the entering label are needed. This has a practical implication for tracing arithmetic witnesses. Excluding harmful labels below a given cutoff on the entering side and above a complementary cutoff on the clearing side can suffice without presenting a full equality of exceptional sets. If both ranges are provided by the original theorem, one should first compare these one-sided bounds. Computing every center would then be unnecessary.

Suppose instead that \(D<H\), and the claimed fixed observable \(g\) really is the same on both sides. If the successful regions supply \(g(x)\le x^\sigma\), continuity also bounds the two scalar values \(g(D)\), \(g(H)\) whenever certified points approach their respective edges. This remains true at a left activation where no arithmetic witness exists exactly at \(D\). Thus an uncovered arithmetic endpoint and an unknown scalar endpoint are different failures. The earlier experiment derived a global slope bound for each fixed admissible character family. If \(K\) is such a bound, then on the remaining gap
\[
 g(x)\le \min\big\{D^\sigma+K(x-D),\ H^\sigma+K(H-x)\big\},
 \qquad D\le x\le H .
\]
In the information class of all \(K\)-Lipschitz curves subject only to the two upper endpoint readings, the displayed minimum on the interval is itself an admissible curve and gives the exact upper envelope. Each branch is \(K\)-Lipschitz and their minimum is \(K\)-Lipschitz. At each endpoint at least the branch based there satisfies the corresponding upper bound. An immediate sufficient criterion for exponent \(\theta>\sigma>0\) is \(H^\sigma+K(H-D)/2\le D^\theta\), since the nearest edge is at distance at most half the gap length. The formula uses ordinates of the scalar sum, after any required rescaling. It offers a way to use a nearly successful handoff quantitatively, while leaving intact the exact arithmetic answer that neither witness covers its missing subinterval. No finite list experiment here supplies those scalar endpoint readings or the constants of the originally selected family. \href{https://github.com/nousentllc/riemann-hypothesis/pull/18}{[Fixed-family slope proof and scope of the scalar continuation]}


\section{At conductor one the slot at 3 can be settled directly}
\label{sec:slot}
Our strongest test case has two unusual properties. We know its base field independently of the auxiliary primes, and we can compute the cohomological invariant which the exceptional list was introduced to bound. Thus we can ask two separate mathematical questions on the \emph{same} field. Which weights would a chosen center mark? How costly is each of those weights in the actual arithmetic sum? The answers can differ even if the generator, the matrix implementing its action, and the tower have all been fixed.

For the conductor-one character \(\psi=\mathbf 1\), the auxiliary set explicitly contains 3 and \(L_\psi(3)=\mathrm{AbF.cyc}(3)\). Consequently every permitted model for this slot represents \(\mathbb Q(\zeta_3)\). A global rank constant \(K\) specifies the layer the window argument uses. Lowering a cap for one component throughout \(K\) would change other definitions, including the outer search exponent \(J\). We retain the given \(K\). Instead we prove a statement uniform in the layer index. It automatically includes the layer selected by this particular \(K\). \href{https://github.com/nousentllc/riemann-hypothesis/blob/40e56bab026552e089d7afcfeff08d2defb9c005/src/Development/AutomorphicArithmetic/AutomorphicVarieties/StmtFields.lean#L58-L67}{[Original auxiliary set]}\quad\href{https://github.com/nousentllc/riemann-hypothesis/blob/40e56bab026552e089d7afcfeff08d2defb9c005/src/Development/AutomorphicArithmetic/AutomorphicVarieties/PrimeSetBoundDataOfPrimitiveChar.lean#L32-L42}{[Original even and odd obligations]}

\subsection{From class groups to the quantity in the original sum}
The earlier class-ascent study establishes the tower class condition in the reference realization \(L\) of \(\mathbb Q(\zeta_3)\). In its stated range the relevant class group has order prime to 3. Cubing on such a finite abelian group is surjective. Its quotient by cubes is therefore trivial. The rank estimates along the tower give a presentation of the projective module by a free module of rank zero. The target of this surjection is a subsingleton. Descent to each admissible finite layer identifies the invariant at a positive weight with a module length, which vanishes. We use its exact previously checked formulation
\[
  \operatorname{shaInv}(T,3,e)=0
   \qquad(e\ge1,\; T/L\text{ a layer of index }m).
\]
The assertion covers every positive weight, including any weight that a center-based inspection might designate as exceptional. This last quantifier is what permits a direct route. A bound known only away from a ball would fail to decide precisely the weights needed here. \href{https://github.com/nousentllc/riemann-hypothesis/blob/4b1c9fbbfb9bc50a3f791c2c1f3d7f2adb336c87/research/uriel/experiments/slot-three/RankZero.lean#L34-L70}{[Rank and descent]}\quad\href{https://github.com/nousentllc/riemann-hypothesis/blob/e674cb750ef96b2f5e7cd46586ea4edb9998715a/research/uriel/experiments/ascent-witness/TowerResult.lean}{[Class ascent and all positive weights]}

There is a small but essential transport step. The local predicate does not promise to use the reference field representation in which the tower theorem was proved. It quantifies over every realization \(F\) of the abelian field and every \(E/F\) with the prescribed layer property. Given such \(F,E,m\), first take an index-\(m\) layer \(T\) over the reference field. Uniqueness up to isomorphism of the base and the corresponding layer identification give isomorphisms \(L\simeq F\) and \(T\simeq E\). Invariance of \(\operatorname{shaInv}\) carries the zero value across. The original summation lemma can now be applied with per-weight bound zero. This also gives the finiteness clause in the source definition of \(\mathrm{VVat}\); a bound on a sum of conversions to natural numbers by itself would be insufficient because the extended infinite value also converts to a natural number.

More precisely, for every integer cutoff \(t\), every allowed local datum \(D\), and every real ceiling \(B\ge0\), the resulting lengths are finite at all \(1\le e<t\) and their sum is at most \(B\). In particular we may set \(B=0\). Nonnegativity forces equality to zero. This remains true at \(t=0\) and \(t=1\) where the positive range is empty. Since the original ceiling \(\exp(x^{4\varepsilon/5})\) is positive, the bound supplies the slot in exactly the format needed. The new kernel-checked statement \texttt{zeroOnModel} performs this transport and calls the original summation lemma. \href{https://github.com/nousentllc/riemann-hypothesis/blob/40e56bab026552e089d7afcfeff08d2defb9c005/src/Development/AutomorphicArithmetic/AutomorphicVarieties/LengthEqAddEigen.lean#L107-L110}{[Finiteness and sum in VVat]}\quad\href{https://github.com/nousentllc/riemann-hypothesis/blob/40e56bab026552e089d7afcfeff08d2defb9c005/src/Development/AutomorphicArithmetic/AutomorphicVarieties/LayeredIsomorphisms/Vvat.lean}{[Summation lemma]}

The new \texttt{eraseThree} equivalence deletes exactly the slot at 3 from the universally quantified odd-field predicate, whenever the character's model at 3 is the cyclotomic field just identified. The forward direction restricts a universal quantifier. To reverse it, take an arbitrary slot. If it is 3 use the zero theorem on the submitted field model and the submitted layer. Otherwise apply the reduced predicate on that very model and layer. The interval theorem \texttt{principalGap} repeats the same implication pointwise for any preassigned \(a,b\), the same cutoff function \(t\), and the same \(K\). Even if two selections pick completely different generators, this direct component never needs a handoff. Eliminating a component of a conjunction does require an independent avenue to the other components. One cannot erase a premise from an opaque implication merely because its conclusion has a redundant conjunct.

\subsection{A numerical label is not the invariant}
To see how much freedom the screening procedure has on this field, fix one generator and a surjective presentation in which the describing matrix is zero. Such data exist by the prior \texttt{same\_lift\_unconditional}. In a zero presentation the root estimate accepts every integral 3-adic center \(s\). The vanishing result also gives the away-from-center estimate at every positive weight. Finally, the single-center filter at any cutoff has cardinality at most one and satisfies the original first numerical bound. These facts give, within one kernel-checked declaration, the order
\[
 \exists(\gamma,\pi)\quad\forall s\in\mathbb Z_3\quad\forall Z\in\mathbb N .
\]
Here the omitted predicate consists of the generator and presentation identities, the root-ball and away estimates, the second bound, and the first bound with its \emph{specified} weight set. In particular \(s\) can be fixed once and subsequently used at every finite cutoff. This is stronger than picking a new integer label after each query. It does not identify the center hidden inside an arbitrary previously issued existential witness.

For this one integral center let \(M_Z=3^{1+\lfloor\log_3 Z\rfloor}\) for \(Z\ge1\). This is the first power of three strictly larger than \(Z\), including when \(Z\) is already a power of three. Strict proximity in the original norm gives the finite description
\[
 E_s(Z)=\{k:1\le k\le Z,\ k\equiv s\pmod {M_Z}\}.
\]
At most one representative can fit. At \(Z_1=144\), \(M_1=243\); at \(Z_2=256\), \(M_2=729\). Seeing 13 on the old list cannot decide whether the additional ternary digit is 0, 1, or 2. These possibilities produce a later label 13, a later label 256, or an empty later list, respectively. The first is active and too large throughout our toy gap. The second lies above all active cutoffs there. The third is empty because the representative 499 exceeds the later outer cutoff. All three continue the same observed old residue and match at all lower precisions. A finer query supplies strictly more information here despite the single-item bound at both scales.

\begin{center}\small\begin{tabularx}{\linewidth}{@{}Yrr@{}}
\toprule
Finite experiment on \([8,12]\) & Shared center & Independent centers\\\midrule
Resolution pairs examined & 729 & 177,147\\
At least one ordinate missed & 135 & 17,957\\
Both records fail throughout & 51 & 2,601\\
Earlier left and later right both pass & 484 & 112,424\\
Both endpoints pass and an interior hole remains & 0 & 186\\
Both endpoints pass and later record alone covers & 484 & 111,672\\\bottomrule
\end{tabularx}\end{center}
For this table \(C_t=2,\eta=1\). We subdivided using \(y=x^2\), on which every entry or exit of a forbidden band has an integer coordinate; testing open cells and individual boundary points is exhaustive for the real continuum. To see the number 51 without the program, a single label must satisfy \(k>12\) at the right end and \(k<64\) at the left end to obstruct the entire interval. The integers 13 through 63 give precisely 51 residues at the later precision. At least some obstruction is possible for 9 through 143. Every later representative in this range projects to itself at the old cutoff. These direct counts also explain why treating refinement as independent resampling would distort the conclusion. The denominators count a finite menu, with no probability law on arithmetic centers.

With the source exponent \(401\), \(\varepsilon=1/64\), and the choice \(\eta=\varepsilon/(2C_K)\), \(C_K\ge1\), the same interval has an almost constant allowable integer size. In fact \(1\le x^\eta<2\) everywhere on \([8,12]\). Here, for a fixed record, a passing right endpoint alone suffices for the size screen on the gap, since the cutoff at the right includes every active earlier weight. This finite simplification neither determines that endpoint nor bounds the global late thresholds. At \(X=12\) and 16 the strict moduli are \(3^{908}\) and \(3^{1013}\). Out of the \(3^{105}\) refinements compatible with the old residue 13, exactly one retains that integer; a particular retained refinement is enough to obstruct the screen at both scales. The explicit center \(13+3^{908}\) instead clears the old active label at the second cutoff. On the zero field both choices still give the identical direct length bound.

\subsection{The same prime row can cross a real ternary-digit boundary}
\label{sec:cross-digit}

Two distinct real inputs might still request the same 3-adic precision. To test the interaction directly, retain the previously certified \(q=3\), the same 1,400-bit \(d\), \(\varepsilon=1/72\), \(J=4\), and the previously prescribed interval \([e^{72},e^{76}]\). At log base \(u=\log Y\), the outer query is \(X=e^{4u}\). Its slot cutoff is \(Z=\lfloor e^{1604u}\rfloor\). Before the new calculation, we specified the ternary transition nearest to 36 and placed one input exactly \(10^{-12}\) on each side in this log coordinate. This choice did not depend on the digits of either candidate center.

The nearest transition is
\[
  m=52{,}561,\qquad u_0=\frac{m\log 3}{1604}
      =36.00010006526465928\ldots .
\]
At \(u_\pm=u_0\pm10^{-12}\), both genuine size bands still hold for the same primes. The strict upper limits on the small primes remain below two, and the congruences and divisibility conditions are unchanged. Both outer ranges contain the entire prescribed interval. The conservative integer guard for retaining an unobserved second prime fails between these two inputs. The actual certified prime has enough slack to survive. The strict moduli for the exceptional list are \(3^m\) and \(3^{m+1}\). To see why rounding the real cutoff cannot spoil this comparison, observe that each \(3^j\) is an integer. It lies at or below \(\lfloor e^{1604u}\rfloor\) exactly when \(j\log3\le1604u\). The directed calculations isolate adjacent exponents on the two sides without constructing the enormous cutoffs.

Now compare the two admissible integral centers \(s=100\) and \(s'=100+3^m\) on the same rank-one presentation at slot 3. Both old records consist of the single label 100. At every prescribed ordinate this label is active and too expensive, because \(101<e^{401\cdot72}\) and \(x^{\eta}\le e^{76/144}<1.696<100\) for \(0\le\eta\le1/144\). On refinement \(s\) keeps exactly the same bad label. The center \(s'\) returns \(100+3^m\) instead, which lies above every active cutoff on the gap because \(e^{401\cdot76}<3^m\). This later list is genuinely nonempty. Indeed its canonical representative is less than \(3^{m+1}\) and the excess of the later outer power over \(3^m\) satisfies
\[
 e^{1604u_+}-3^m
   =3^m(e^{1604\cdot10^{-12}}-1)
   >3^m\,1604\cdot10^{-12}>100 .
\]
Consequently the later screen for \(s'\) succeeds on every prescribed ordinate, while that for \(s\) continues to fail everywhere.

This control isolates a digit change while holding the character, prime pair, target and earlier exceptional record fixed. It distinguishes forgetting an old harmful weight from producing an empty later record. The original large-prime guards exclude \(q=3\) downstream, and this numerical \(J\) is the finite row parameter. The witness supplies no twist bounds, global host or identification of the odd table chosen by a particular \(K\). The two centers belong to one admissible rank-one presentation, and the direct length bound gives zero for both. In particular neither the agreement at the first reading nor the shared arithmetic row predicts the second reading. The missing input is precisely one ternary digit of the center, at a position identified in advance by the outer schedule.


\section{One center can resist every legal cutoff on a fixed sequence of gaps}
\label{sec:sparse}
A finite failure might disappear as soon as one requests a much larger outer scale. Whether one has that freedom depends on the coordinate being certified. For fixed \(x\), the ambient condition \(X^{1/J}\le x\le X\) sets an upper limit \(X\le x^J\). Increasing \(X\) indefinitely would eventually put \(x\) below the lower boundary. This observation makes possible a single explicit example resisting every legal refinement, indefinitely far out.

\subsection{Construction and order of quantifiers}
Set \(n_1=2\), \(n_{j+1}=n_j^2\), \(A_j=3^{n_j}\), and prescribe \(I_j=[A_j,A_j^{3/2}]\) for all \(j\). The intervals are fixed without consulting an outer query. In the ring of 3-adic integers choose
\[
      s=\sum_{i\ge1}3^{n_i},\qquad N_j=\sum_{i\le j}3^{n_i}.
\]
The series converges since the valuations of its terms increase without bound. After division by \(3^{n_{j+1}}\) its tail \(s-N_j\) is a 3-adic unit: its reduction modulo 3 is exactly 1. Therefore \(v_3(s-N_j)=n_{j+1}\). The ordinary size obeys \(A_j\le N_j<2A_j\), since the sum of all distinct smaller nonnegative powers of three is less than \(A_j/2\). This one pair of inequalities uses the same integer in two geometries. Archimedean size places it among the active but expensive weights, and 3-adic proximity makes the screen remember it.

Fix any finite \(J\ge1\). If the whole interval fits one envelope then \(A_j^{3/2}\le X\le A_j^J\). When \(J<3/2\) no such envelope exists, so the stronger pointwise version below is the informative claim. At the source exponent, \(Z=\lfloor X^{401}\rfloor\le3^{401 Jn_j}\). Once \(n_j>401J\), we have
\[
    \lfloor\log_3 Z\rfloor\le401 Jn_j<n_j^2=v_3(s-N_j).
\]
Also \(N_j\le Z\). Thus \(N_j\) is on the exceptional list of \emph{every} outer panel that could contain the prescribed interval. For every \(x\in I_j\) its active cutoff satisfies \(N_j+1\le x^{401}\), while \(x^\eta\le A_j^{3\eta/2}<A_j\le N_j\) whenever \(0\le\eta<2/3\). This proves failure at every point for this filter. Positivity of the allowance exponent is unnecessary for the particular comparison. The source rates just considered are well inside this range.

Even a separate panel for each interior point offers no general remedy for this same choice of center. For \(x\in I_j\), any panel on which \(x\) is admissible has
\[
 A_j\le x\le X\le x^J\le A_j^{3J/2}.
\]
Consequently the stronger sufficient condition \(n_j> (3/2)401J\) works for all such panels simultaneously. The lower bound on \(Z\) is still large enough to include \(N_j\). This variant follows by the same strict valuation test; it is an additional deduction after the original protocol. The sequence of centers does not change with \(J\); there is just one center for all \(J\). The onset index does depend on \(J\). No universal numerical onset is asserted.

\subsection{Checking the arithmetic without constructing astronomical integers}
The first three partial sums are 9, 90 and 43,046,811. Their first disagreement with a longer truncation occurs at the ternary digit of index 4, 16 and 256, respectively, numbering the units digit zero. Beyond these examples it is sufficient to compare exact integer exponents. For \(J=4\), \(n_4=256\) is still below \(401J=1604\); \(n_5=65{,}536\) exceeds it. The margin in valuations is \(n_5^2-401J n_5=4{,}189{,}847{,}552\). At \(J=333{,}632\) the first index satisfying the same sufficient inequality is six. These indices characterize our sufficient test. They do not estimate when any unknown global arithmetic threshold is attained. We include pointwise tests for six values of \(J\), and exact congruence and real-size checks on the smaller partial sums.

This admissible presentation has rank bound one and local constant one. The globally retained \(K\) may choose other ranks, fields of coefficients or tables. The sparse construction does not identify the odd table selected by that global witness. By contrast, the zero-length proof applies to the layer demanded by every such \(K\).

The example disproves a tempting form of compactness reasoning. For each \emph{fixed} integer \(k\) different from \(s\), sufficiently fine precision will eventually clear \(k\). Along our intervals the integer in question is \(N_j\); the precision needed to clear it grows faster than every legal cutoff budget of the form above. Pointwise eventual clearance of each fixed weight gives no uniform rate over the moving active range. Conversely, the center 0 has an empty list at all positive cutoffs since \(M_Z>Z\). The slot accepts both behaviors on the very same generator and presentation. Thus the example refutes a uniform guarantee for admissible records. It leaves open the existence of favorable choices in other slots, and in this slot a favorable choice is explicit. The cofinal obstruction and the semantic zero theorem are mathematically consistent at every one of the prescribed points.


\subsection{The conductor-one roster closes to two odd primes}

How many local restrictions would a conductor-one gap certificate have to reconcile? The term indexed by $3$ is compulsory because the source inserts $3$ directly. A second prime survives. Combining the source definitions with a short class-number argument gives
\[
 B_{\chi}(\mathrm{one})=\{3,7\},\qquad
 L_{\chi}(\mathrm{one},\ell)=\operatorname{AbF.cyc}(\ell).
\]
These are two different pieces of information. The first equation counts the indices that contribute. The second identifies a field descriptor for each index. Neither equation chooses the obstruction data over the corresponding field.

The primitive character called \texttt{one} has conductor $1$ and has order $1$. The first auxiliary prime $p_K$ is the least prime at least $7$ that is $3$ modulo $4$ and coprime to the conductor. Its lower endpoint $7$ qualifies. In the definition of the second auxiliary prime $p_1$, the odd part of the character order equals $1$, so its congruence condition is automatic. Coprimality with $2\cdot1\cdot7=14$ eliminates the prime $2$ and admits $3$. The remaining number in the roster formula is $h_*$, the class number of the splitting field of $X^2+7$ over $\mathbb{Q}$. A positive class number already ensures that $7$ divides the product of the five factors. The appearance of slot $7$ thus has a shorter proof than the assertion that there are exactly two slots.

The splitting field is $K=\mathbb{Q}(\sqrt{-7})$. To compute its class number, take $\omega=(1+\sqrt{-7})/2$. An algebraic integer $z$ has the form $(t+u\sqrt{-7})/2$ with integer trace $t$, integer norm $n$, and rational $u$. From $7u^2=4n-t^2$, the square of the reduced denominator of $u$ divides $7$, so $u$ is integral. Reducing $t^2+7u^2=4n$ modulo $4$ shows $t$ and $u$ have the same parity. Consequently $z=(t-u)/2+u\omega$. Conversely $\omega$ satisfies $\omega^2-\omega+2=0$. This establishes that the entire ring of integers is $\mathbb Z[\omega]$ and justifies using its lattice to test Euclidean division.

Given a complex number $w$, choose $k\in\mathbb Z$ so that $w-k\omega$ has imaginary part of absolute value at most $\sqrt{7}/4$. Then choose $j\in\mathbb Z$ to bring the real part within $1/2$ of zero. The squared error is at most
\[
 |w-j-k\omega|^2\leq\frac14+\frac7{16}=\frac{11}{16}<1.
\]
For $w=a/b$, with $a,b$ algebraic integers and $b\ne0$, multiplication by $b$ now gives a remainder of smaller absolute norm. An element of smallest positive norm in any nonzero ideal therefore generates that ideal. Hence $h_*=1$. The source takes odd prime factors of $1\cdot1\cdot7\cdot1\cdot3$ and reinserts $3$. The full cyclotomic descriptor at $7$ refers to $\mathbb Q(\zeta_7)$. The smaller field $K$ entered solely to calculate which indices occur.

Six local Lean declarations were compiled and checked. They establish the character order, both auxiliary minima, unconditional membership of $7$, the roster equality assuming $h_*=1$, and the formula for the field descriptor. Each axiom report contains only propositional extensionality, classical choice and quotient soundness. The argument computing $h_*$ above is a written proof. The independent rational calculation checks the rounding constant and the finite arithmetic. These distinctions matter for replay: the kernel theorem about set equality exposes the class-number premise explicitly. The remaining experimental question is correspondingly precise. One would have to certify intervals for obstruction data over \emph{both} cyclotomic descriptors, and prove that those interval certificates apply to the actual uses of the arithmetic selections at the two ends.


\section{What the genuine prime sum does on the prescribed gap}
\label{sec:gap-principal}
The fixed interval \([8,12]\) admits another test with a different kind of evidence. For the conductor-one principal character the previous experiment traced the cyclic level to the normalized zeta prime sum and showed that the fixed rescaling corresponds to one. We can therefore evaluate the actual scalar function on this very same interval, using no scale conversion. The numbers 8 and 12 were chosen before inspecting its values. Positive or negative outcomes here classify one finite scalar inequality. They carry no implication about which pairs, exceptional labels or fields the large-scale arithmetic construction would have returned at these ordinates. In particular the existential lower thresholds in the field argument have no numerical values in this example. \href{https://github.com/nousentllc/riemann-hypothesis/blob/d79bd956c6fb2b8a3de51b606765be81d24fb1fc/research/uriel/experiments/cofinal-endpoint-certificates/evidence/transport/extension/UnitCases.lean}{[Checked identification of the principal level]}

\subsection{A calculation on every real point}
For real \(x\ge1\) the original principal normalization is
\[
 G(x)=\sum_{n\le x}\Lambda(n)(x/n-1)-x\log x+(1+\gamma)x,
 \qquad H(x)=G(x)/\sqrt{x},
\]
with \(\gamma\) the Euler--Mascheroni constant. The summand of a newly encountered prime power vanishes at its entry, so \(G\), and hence \(H\), is continuous at each knot. Its derivative can change there. The nonzero von Mangoldt terms up to 12 occur at \(2,3,4,5,7,8,9,11\); every power \(p^e\) contributes \(\log p\). The second \(3\)-power and the prime 11 are the two interior knots. A calculation at integer nodes alone could miss an extremum strictly between 11 and 12. \href{https://github.com/nousentllc/riemann-hypothesis/blob/40e56bab026552e089d7afcfeff08d2defb9c005/src/Development/AutomorphicArithmetic/DedekindLFunctions/CharFamily.lean#L10-L23}{[Original ramp and principal correction]}

On an open knot interval let \(A=\sum_{n\le x}\Lambda(n)/n\) and \(B=\sum_{n\le x}\Lambda(n)\), which are constant on that interval. It suffices to study
\[
 H(x)=\sqrt{x}(A+1+\gamma-\log x)-B/\sqrt{x}.
\]
After differentiating once and clearing a positive factor, we obtain
\[
 D(x)=2x^{3/2}H'(x)=x(A-1+\gamma-\log x)+B,\qquad
 D'(x)=A-2+\gamma-\log x .
\]
Since \(D'\) decreases on an open cell, an upper bound for it at the cell's left edge proves strict decrease of \(D\) throughout that cell. The directed calculation gives upper bounds \(-1.929643250928\), \(-1.925358254510\), and \(-1.908038470627\) on the three successive cells. The left and right limiting values of \(D\) are both negative on the first two. At 11 its right value becomes positive, approximately \(0.2414862764\), and the value at the right edge 12 is below \(-1.7106887180\). There is exactly one interior stationary point. It is a strict local maximum. The jump in \(D\) at a prime power \(n=p^e\) is \(2\log p\), because both \(nA\) and \(B\) increase by \(\log p\). This explains how a strictly decreasing derivative numerator can restart above zero in one cell after having been negative in the preceding cell.

We isolated the zero by rational bisection, branching only after the directed evaluation of \(D\) enclosed a strictly positive or strictly negative value. The resulting enclosure for its location is
\[
              11.12618470653389\ \le x_\ast\le\ 11.12618470653390 .
\]
Interval evaluation of \(H\) over that whole root bracket encloses the value at the unknown true critical point. We compared that enclosure with the enclosures at every knot and both boundaries. The resulting bounds are as follows. The values in the second column are rounded to the displayed precision; the adjacent machine-readable bounds retain outward rounding.
\begin{center}\begin{tabular}{rll}
\toprule
Ordinate&\(G(x)/\sqrt{x}\), approximately&Role on the fixed target\\
\midrule
8&0.64680899314394&global maximum\\
9&0.61325384296753&prime-square knot\\
11&0.53719292425826&local minimum\\
\(x_\ast\)&0.53740075021814&interior local maximum\\
12&0.52812895767871&global minimum\\
\bottomrule
\end{tabular}\end{center}
In particular, the global maximum belongs to \([0.64680899314393,0.64680899314394]\), and the global minimum belongs to \([0.52812895767871,0.52812895767872]\). The actual sum satisfies \(G(x)\le(2/3)\sqrt{x}\) on the entire preassigned gap. It satisfies \(G(x)\le x^{1/2+\delta}\) there for every \(\delta\ge0\). This finite verification uses direct information about the normalized prime function at each analytic piece. No arithmetic witnesses were assembled to obtain it.

Comparing just the exterior readings would happen to return both global extrema. It would supply no proof that these are indeed global. The small local peak distinguishes these statements. As a post hoc control choose \(0.5373\), strictly between the value at 11 and the peak of the last cell. Both \(H(11)\) and \(H(12)\) pass this limit while \(H(x_\ast)\) exceeds it. Thus within the same fixed gap there is a directly exhibited subinterval where two apparently satisfactory scalar readings miss a violation. This contrast was selected after observing the peak. Its role is to show how continuity and an upward derivative jump can coexist with a missed excursion; it supplies no distribution of such failures.

\begin{figure}[htbp]\centering\bandfigure{.98\linewidth}{scalar}\caption{The genuine principal-character ramp \(G(x)/\sqrt{x}\); the shaded part of the left panel is enlarged on the right. The small local peak near \(11.12618\) crosses the post hoc ceiling \(0.5373\), while both ends of \([11,12]\) pass. The entire original target satisfies the stronger uniform bound \(G(x)/\sqrt{x}<2/3\). The plotted curve illustrates the certified enclosures in the text.}\end{figure}

\subsection{How the numerical claims were bounded}
The program uses rational arguments and outward arithmetic on a binary grid with denominator \(2^{160}\). Logarithms are reduced to \([1,2]\) and evaluated through the odd-power expansion in \((v-1)/(v+1)\), whose absolute value is then at most \(1/3\). A geometric tail bound closes the finite sum. Integer square roots rounded in both directions enclose every required square root. The fixed-point module is the one distributed in the preceding experiment, and the new code adds the stationary-point tests and the following two normalizations.

For an accurate value of \(\gamma\), take \(N=256\), compute \(H_N^{\rm harm}=\sum_{j=1}^N 1/j\) exactly, and set
\[
 C_N=\frac1{2N}-\frac1{12N^2}+\frac1{120N^4}
                  -\frac1{252N^6}+\frac1{240N^8}.
\]
Euler--Maclaurin summation for \(1/t\), taken to the eighth derivative on \([N,M]\) and then to \(M\to\infty\), gives
\[
 \left|\gamma-\bigl(H_N^{\rm harm}-\log N-C_N\bigr)\right|
               \le \frac{1}{240N^8}.
\]
Indeed the absolute periodic eighth Bernoulli polynomial is bounded by \(|B_8|=1/30\), directly from its cosine series at even order. The integral of \(8!/t^9\) on the tail is \(7!/N^8\). Dividing their product by \(8!\) gives the stated remainder. This procedure encloses \(\gamma\) between \(0.577215664901532860606\) and \(0.577215664901532860607\).

For an independent, wider check take \(N=10^6\) and sum outward on the binary grid. The correction \(H_N^{\rm harm}-\log N-\gamma\) is the sum over \(k\ge N\) of
\(\int_k^{k+1}(t-k)t^{-2}\,dt\), by integrating \(\log(1+1/k)-1/(k+1)\). Bounding each integrand above and below by replacing \(t^{-2}\) by the endpoint values bounds the total between \(1/[2(N+1)]\) and \(1/(2N)+1/(2N^2)\), after integral comparison on the series of inverse squares. The resulting second enclosure contains the entire first enclosure. This check isolates the harmonic correction and still shares the directed logarithm implementation. It does not independently test every rounded operation.

At each of the five integer nodes a further calculation reconstructs the von Mangoldt coefficients recursively from \(\log n=\sum_{d\mid n}\Lambda(d)\). Factoring \(n\) gives its integer coefficient vector on the basis of formal prime logarithms. Subtracting the vectors already assigned to its proper divisors recovers the required vector for \(n\). The rational coefficients of the resulting ramp agree term by term with the prime-power construction. The normalized node values computed using the independent harmonic enclosure include each of the sharper primary values. Agreement therefore tests both the divisor support and the smooth normalization, while the stationary-point calculation separately certifies the behavior between integers.

Three ablations illustrate those obligations. Omitting \(9\) leaves the global maximum at 8 unchanged but moves the global minimum at 12 down to approximately \(0.42241494091245\). It also removes the interior maximum. Looking only at the upper bound on the whole interval would miss that error. Retaining \(-x\log x\) but dropping the entire linear correction \((1+\gamma)x\) gives an apparent normalized maximum near \(-3.81423057503816\), which would pass the upper-bound test for a spurious reason. Conversely the ramp alone is strictly increasing in its normalized form on the measured cells. Its maximum is approximately \(3.67246276446990\), which fails even the unit ceiling. The sign and magnitude of the normalization matter independently of whether the prime-power support was enumerated accurately.


The scale of the peak also suggests what a useful error tolerance means here. The normalized values at 11 and at the local maximum differ by only about \(0.0002078260\). The larger change caused by dropping 9 can be evaluated at any \(x\ge9\) without repeating the search. The missing contribution to the normalized height is \(\log(3)(x/9-1)/\sqrt{x}\), reaching about \(0.1057140\) at 12. At \(11^+\) the missing contribution to the derivative numerator is \((1+11/9)\log3\). It is large enough to reverse the sign of \(D(11^+)\), explaining the disappearance of the interior maximum. Directly checking the scalar at an early point before that omitted term has entered, even with unlimited numerical precision, cannot expose this support error.

In contrast, changing the Euler constant by \(u\) changes every normalized ordinate by exactly \(u\sqrt{x}\). Thus over this fixed interval the entire plotted curve moves at each point by at most \(\sqrt{12}|u|\). Even the wider harmonic enclosure is comfortably below the precision needed to distinguish the intentionally chosen local threshold. The higher-order remainder permits tighter root and height enclosures while the monotonic derivative proof establishes completeness of the candidate list. Without the latter step, reporting additional decimal places at the integer nodes would improve precision for a different proposition. The original target requires every real point of the interval.

The calculation is reproduced by running \texttt{python3 -B extrema.py} in the scalar evidence directory and comparing the JSON with the archived receipt. Repeated runs under the two installed Python versions returned the same bytes. The critical-point isolation and Euler constant estimates are directed numerical and written analytic checks, respectively. They are not a Lean replay of the scalar computation. The fixed-gap positive result answers precisely the independent finite scalar question. It leaves open the identification of active exceptional sets in two actual successive arithmetic selections.


\section{What the two calls jointly certify}
\label{sec:conclusion}
We can now state a single implication in the notation of the original physical witness. Fix the principal character, admissible tolerance, constants and a page function \(e\) satisfying its published bounds. Use \(C_t=401\) in the arithmetic data, physical scale factor one, and \(J=J_\chi(K,401,\varepsilon,\mathbf 1)\). Assume the source prime supply \texttt{PairEx}, the pointwise height theorem \texttt{PointHV}, the field conversion \texttt{ThmBc 401}, and the coherent field supply at the stated rates. The last can be constructed from the fixed-field, moving-field and \texttt{LemmaK} inputs as shown earlier. All original input predicates retain their model and layer quantifiers.

There is a late threshold and a family of odd tables. Now prescribe \(a\le b\) and choose a late \(Y<a\) with \(b\le Y^J\). There is a positive rightward radius and one pair \((q,d)\), together with its moving-field table. For \emph{every} \(V\) in that right neighborhood, the pair supplies the original \(\mathrm{PairGood}\) record at \((V,V^J)\). In particular we can take \(V>Y\) and use genuinely distinct scales. Define \(U=Y^J, X=V^J\). On the moving table and on the remaining odd tables require the earlier record to pass at \(a\) using \(\lfloor U^{401}\rfloor\), and the later record to pass at \(b\) using \(\lfloor X^{401}\rfloor\). For each applicable allowance exponent \(\eta\), require \(b^\eta<\lfloor a^{401}\rfloor\). Then one has
\[
 \forall x\in[a,b],\qquad
 \mathrm{WinDatum}(\mathbf 1,\varepsilon,J,1,x,V,q,d,e(x)).
\]
This is the conclusion of the new combined Lean theorem \texttt{GapFirst.Bridge.twoSelections}. The same physical mark \(V\), the same primes, and the same page function apply to every \(x\); its page number \(e(x)\) may vary. The same finite set need not serve at the two cutoffs. The odd tests omit 3 from the outset because a direct proof supplies it during the assembly of the field estimates. The supplementary \texttt{checksAtSeven} proves the endpoint equivalence with just one remaining odd test under the quadratic class-number identity; the written proof supplies that identity. The classical roster calculation reduces the remaining odd tests to the single prime 7. Thus the outstanding inputs have concrete identities: the moving 2-field of this row, and the fixed slot at 7 over its cyclotomic base. For each, there is a specific old test and a specific new test. Full certificate availability on an arbitrary late target remains open in this study.

\subsection{Three different answers on the very same interval}
On the finite interval \([8,12]\), our directed computation proves a genuine scalar ceiling, \(G(x)/\sqrt x<2/3\). It attains its largest value at the left end. Every permissible layer over the base in slot 3 also satisfies its genuine local length bound there, in fact with ceiling zero. Meanwhile the equally admissible center 13 causes the source-shaped screen in that slot to fail at every point at the two tested source cutoffs. These assertions have the same interval and the same principal character. There is no inconsistency: one gives the value of an analytic observable, one gives a bound for a cohomological invariant, and one evaluates a sufficient screening procedure. Conflating these verdicts would mistake a failed method of certification for a mathematical counterexample.

The terminal field theorem only consumes the later \(\mathrm{PairGood}\) record. If a later pair is already known, one can define its own table at an earlier cutoff and test it there, even when that pair could never have passed the earlier prime-size bands. An older recorded answer on a different moving field supplies no such test. The joint pair theorem is valuable precisely when one wants to reuse information from two genuine outer requests rather than retrospectively request an earlier cutoff on the later field.

At late scales the source additionally makes a promise at each outer query: at least one inner interval will succeed. Suppose only its placement data are retained, so \(Y\le W\le Y^{J/2}\) and the marked window is \([W,W^2]\), with \(Y>1\) and \(J\ge2\). The intersection over all allowable placements is \([Y^{J/2},Y^2]\) if \(J\le4\), and is empty if \(J>4\). Indeed the last possible start and earliest possible finish supply the endpoints of the intersection. For the arithmetic construction the rank sum \(R\) is a nonnegative integer and \(J=2^{R+2}(2C_K\,401/\varepsilon)^R\) with \(C_K\ge1\). At rank zero this is 4, giving a single guaranteed placement point. With \(R\ge1\) and \(\varepsilon<1/52\) it exceeds 333,632. This coarse positional information by itself supplies no nondegenerate prescribed gap for this construction. We have used substantially more than the geometry of the displayed window.

\subsection{What is decided, and at what cost}
\begin{center}\small\begin{tabularx}{\linewidth}{@{}p{.26\linewidth}YY@{}}\toprule
Question & Verdict & Cost or remaining input\\\midrule
Can successive outer calls use identical primes? & Yes on a sufficiently small right neighborhood. At fixed power spacing the required size bands eventually separate. & Prospective witness with strengthened modulus and reserved horizon. No identity extracted from two independent old calls.\\
Can distinct exceptional sets support one bridge? & Yes. A later exceptional integer within the previously inspected range must have been marked earlier for coherent tables. & The later envelope must be valid throughout. Both endpoint tests and both slot-specific floor inequalities remain in the theorem.\\
Does a bound on list size force those endpoints to pass? & Even a coherent list of size at most one can block every eligible panel on cofinally many fixed gaps. & There is one admissible 3-adic choice for the obstruction. A different admissible center can clear the same slot.\\
Does the actual arithmetic necessarily fail where its screen fails? & For the conductor-one slot at 3 the length sum equals zero everywhere. & The direct proof uses the previously established class-ascent theorem and transports it across field models.\\
Is the principal prime sum controlled on the finite gap? & Yes, on the entire continuum \([8,12]\). & An independent enclosure, specific to this interval. Its two endpoint values alone cannot serve as a general interpolation rule.\\\bottomrule
\end{tabularx}\end{center}

\subsection{A discriminating next question}
The fixed slot at 7 now offers a sharper test than enlarging the enumeration of arbitrary lists. Its field descriptor is exactly \(\mathbb Q(\zeta_7)\). A formal computation of the base class group alone would not suffice for the result used at 3. The argument at 3 needed a statement along the tower, the passage to the controlling module, and a descent theorem valid at every positive weight. An analogous theorem at 7 would leave just the moving field. A weaker result could also suffice: an estimate for the two actual endpoint bands, uniform in the relevant layer index, with constants compatible with the very same \(K,J\). The falsifying test is clear: if only a base-layer computation is available or if the bound loses its uniformity when transported to the requisite layer, the direct discharge has not been obtained. The moving field presents a different question because its base depends on the auxiliary primes. There the first useful comparison would fix the pair and seek an explicit estimate on its table before comparing different pairs. No improvement to the published Riemann-hypothesis conclusion is inferred from either research direction.

\section{Reproducibility and verification scope}
The four primary files named \texttt{PROTOCOL.md} were recorded separately in a preliminary Git commit. Subsequent controls are labeled in their own files. These include the actual scalar evaluation on the already fixed gap, the deliberately chosen ceiling on \([11,12]\), the study of transient failure bands, the explicit odd roster, the paired digit-boundary experiment, and the stronger pointwise version of the sparse-center obstruction. A follow-up test can expose a failure mode without becoming a prespecified discovery statistic.

The companion directory contains a standard-library Python runner for the enumerations and the directed arithmetic, pinned source maps, and the supplementary Lean modules. The packet includes the text of the 18 research modules needed by the slot replay, with the path and Git object of each. It does not duplicate the entire original project. Replaying the formal checks requires the specified Lean release and a built checkout of the pinned Mathlib revision. The replay code checks original inputs against the recorded hashes; depending on the unit it can also authenticate cached compiled objects or rebuild missing project inputs from source. A fresh isolated directory is used for the combined bridge so that independent files carrying the same namespace cannot conceal a missing import. The checksums file and commands in the package README give the inventory and entry points.

Our strongest new conditional result is a kernel replay of the combined bridge on the source's types. Its hypotheses explicitly include the arithmetic supplies. The combined replay rebuilt all nine supplementary modules together against 1,150 authenticated original project modules and 18 pinned earlier research modules. A separate cold rebuild for the prime unit covered its 21 original project dependencies. The larger field and class-ascent units used authenticated objects and the pinned original source. The auxiliary roster calculation has a short handwritten norm-Euclidean argument for the class-number equality, with its formal boundary recorded in that section. The sparse infinite counterexample has an ordinary mathematical proof, a checked theorem admitting the center, and independent finite controls. The numerical enclosure of the zeta ramp uses directed rational arithmetic outside Lean. These scopes are individually reproducible. A clean build of the entire original 174,793-file Lean source tree was not run in this experiment.

\medskip\textbf{Primary revisions.} \href{https://github.com/nousentllc/riemann-hypothesis/tree/40e56bab026552e089d7afcfeff08d2defb9c005}{Original proof tree}\quad
\href{https://github.com/nousentllc/riemann-hypothesis/tree/08188ff0c6dd02dfec0bd65d30765862e9db56b7/research/uriel}{Research baseline and traces}\quad
\href{https://github.com/nousentllc/riemann-hypothesis/tree/4b1c9fbbfb9bc50a3f791c2c1f3d7f2adb336c87/research/uriel/experiments/slot-three}{Conductor-one slot at 3}\quad
\href{https://github.com/nousentllc/riemann-hypothesis/tree/e674cb750ef96b2f5e7cd46586ea4edb9998715a/research/uriel/experiments/ascent-witness}{Class ascent}\quad
\href{https://github.com/nousentllc/riemann-hypothesis/tree/d79bd956c6fb2b8a3de51b606765be81d24fb1fc/research/uriel/experiments/cofinal-endpoint-certificates}{Previous endpoint study}\quad
\href{https://github.com/nousentllc/riemann-hypothesis/commit/2ddb0b1a35431324d62739d775c65f1b3f4df46b}{Primary protocols for this experiment}
\end{document}

BEGIN FIGURE scalar
255044462d312e340a25acdc20abba0a312030206f626a0a3c3c202f54797065202f436174616c6f67202f506167657320322030205220
3e3e0a656e646f626a0a382030206f626a0a3c3c202f466f6e74203320302052202f584f626a656374203720302052202f457874475374
617465203420302052202f5061747465726e2035203020520a2f53686164696e67203620302052202f50726f63536574205b202f504446
202f54657874202f496d61676542202f496d61676543202f496d61676549205d203e3e0a656e646f626a0a31312030206f626a0a3c3c20
2f54797065202f50616765202f506172656e74203220302052202f5265736f7572636573203820302052202f4d65646961426f78205b20
30203020363132203234342e38205d0a2f436f6e74656e7473203920302052202f416e6e6f747320313020302052203e3e0a656e646f62
6a0a392030206f626a0a3c3c202f4c656e67746820313220302052202f46696c746572202f466c6174654465636f6465203e3e0a737472
65616d0a789ced5bcb921cc515ed757d8257155e490b4af9b8371f0b2fc036b2bd71008af002b3c0232184850089873fdfe7dcaaaecaec
a91e356339a2e5b008a19ed35959f79df791f3e80fcf7e7e71f3ecd3c71f8dbfff6c78b4fd74f366f0e337f8fb7c74e337f8fbcbe8c7c7
f8fb7c70f8e9db21f9807f5fdabf41642af8ecd64f5f0fc357c3a30fb1fc0d563d1e0695495219c54fc9273c1e4a9a9c04cde988bd6cb1
e0dc9423c1e5c10db0ad7f18fb0d43f49366c1a3218320af75c23faf9f8d7f1b5f8d8f3e0c46c754e3f1cf801f3415ac16f563f7c3ebe7
e0f2976ef5787ef5a78f870b57be7e3e0c212472516a682401910549bed446121bb631de3cdc49e3a3e193f1878dc990963f193fc496b2
58839fe99999845aa73cfe325cf8cca78fcfef3e9cec3e689e420daed1f88ac4b27009ca21e9e107fcdf8d1f385044a5060187a9161943
99420e41c79b6f878f9e0c8f3ef6631d9f7c6506f8e4e9f0f9f8e0501e8e5f8c4ffe32fcf1c9f0c93073e37d9c32e524dbbb1bec8eb77b
0f85b924d5d8bbe0ed75e7ed1932520da5e1bcc1ee7a7b924983ab3efaa8e992d7fb83bb4dc0ae896dd81d0484e027e782faece1859711
e0770828710abe408aadb7afd85d04e43c49111fbccba15c4640e8097025440d55189c26df38220244cc41a5d2f4cfad1bda7566ee17ed
d76951d2a26bd3a2d749351e5970dc44135802fd2b3f0f0e7f7d383ef966c8785b84ddc7ac33870f0eafed0bd8ae73be6af461fde6a97d
a313a25d81cd6e5fbcb02fc2a4ce4929a15898b66f5e1d1f71e2dd4cc6f2cd97cb37ead5c7943d0c75f9e6c7850032e48ad6b87ef3ecf8
4c84b9d486e67121a0a4859de591471f8723bfffb225a63286ae5f1fcac3f897257e25b5000619f89a5c8683c3878f81482d8a9b9a5ca0
42fb036e79af4258f8d31f4a471076fad955c6d753dad577c0be7f21bce13b8144634c777a963b4c073dc4d6bbae405525400cee445547
f0ca55b5d1bea8aaa1fbacaa72992a377abbaad295a9cafb3021789fe86a45af5c590df58bb65acacfaacbbb321577916bd56bd39784c9
a553df5ad16bd7d746fd515f0de5e7f5c56c245ee05fa94f34ae415fd9e3f95bfe7544af5d5f1bf5477d35949fd7572a93e40bfc2b1df4
3a124352ff01f940726d2c94327928aa490893e5757d86f427cb9090b5a7acb9ac795f9333fefdc1313543648a0e223c9767e914438278
4bded9e5e1b2242449d2e59ebfb36f2a8a369118b7d4afd9fff14265cd96b4374bfe3b64dedee5c1e1d1f28d8b01654e93d5fef6370b75
485b4b8a73263c4169a7b25eb351f8344a0e8d733e8d0741498657af544db63d9cfaabf1937b4701fa171b0c9f8f11ff7d016a9ffe0a4f
db8f05856e712b761fd13516dca76712189e68d3b32fa7f947062c150f4bc61bd2316ccc7d93ed89f1ec13c31c342e58697d93b577e013
7c2204273277144a0852520b4385110e2586d5889a3211c34e306480798a784d1c6f06f0eb50a5493134a156421cd6babe098e5ae19cb9
c796d7d8e347149f34888eed8e15beeb838bedcb33bc9dc12bf3f1954ea030ebe8fdd83004f0c8d0cafb86e171c606d02ddd52d69e58dc
eda9b0e35c6057fdeb51b939f844ec48d5db2ce92eefba23a566c74d9cdbdb5bc9af74363ada186a94b9a3f79bade7f5deda73d38f4225
86a0e445b2f5a9044212575b982082579179316290cf58e07542542ab518ca842152b6c4b1a42a2410b19b0fa1645bbdbe10db215a8bef
513dbe6fdea3c14bf145e3d8ed9c58bb80c19e0eea5053b22d1aa2f3e4c56508a2e52f6ffc35c2c82d15710a02eb28fdea30e5281a4e36
4631a555b23fa102a72e68f0b927d9efb2e7cf08c3ef8aaed9b911734346a79486e846850d8371571c71a3e3fdb7faa60d09c7ceaee058
b6ee24bc3ce5000bda6082080b4e6d6df5552b780d8836886bce423a422af2e89b81a82275d04a34a59a5db4a5c7774122200961af0797
37cd1b2c706202c7f6f3d8ec0a304b7121b604243f3115c9d9365889052c1541d08f0d5b0057b656116ce08d4926b8c4c8dfaed509cce4
d8efca3c09fbd6130a6482d959ded2502b7b7cc9be0c644f5eebaead64570a3a25acc4b6fa5ad96a35bb6305ff13e6dd34be212d080e89
c0dcfbc627d6721b8cc0383149535b8bc2408a3710892a8287e0f82dce1ac7785e27072d434a8093aa722a04707997b28490681dfb165c
de346f70843d7b181e6548b32bc01c111a4b438086c9bb5253b0e7575a01a34870d4cbc615c095ab55021bc80df0432988faed52bc21d7
18fb4d5103227b4592d01300f3471e106a47abdfe3caef4bc0ef496bddb595eb4a41af8395da565d47b65abdeed8c0cdad291d3b015663
5f5a977799ffb065fedb74ed744c8a94ff9299ebbd2b83633323cc8d824b9ce456f9d2e6f565cbeb11d814c2b5941145356dece5909092
40718183329d2aca551cbd2f07e4b4386d7d98e7672ec138c8564652927335988537e2b010ae38ed052f329ed4398444c03070ad5e0d66
bd18bcaf8499e19768300ef384533b02ae9055d68250e8614d38d311cb08d388222d8733a0e21d7b002f07640b5341944e622d297c2c85
13720e1b056993a5044ca21147b3e171270d218e6a0ea6b9a42160a22696763802987fc19fe7742109fe311cc5441567784079ecbc098c
3f2407228893049d65c0640a997a44dee479f4e10870b61e4457f88610470a05b6bcd10ffbc69b102d4692805408991171882ac2759967
059c50585d6c7f25fd65ce92706e38389ed109b5213713164d9e6710beb17d5896e0ad340ece656b644a071caa65cec4b2d725e25c021c
9eabea13734b872c2c079a1c711423198655c70a734ac83e4d6c156286f48106363f20678ed8e13e3123211ce12a48e420644313bbf028
d7c682f081e891684cb006bc31173c99e394020c7d1ed323f2f862c94062620835f6d3fbeddc050c57abd8321b5c9d432e6a709e5cb5bc
1470c964c7ee0590585e37200c632f6586c119ccd05bfdcaf023462054e1aad46c55b1309731025991d259adbe920a6fa5834108f05c8f
fc9eb518b3f464303cccc52a88ae081225d73acb046a488515db88b70468cddb2ba18518d50a3a4f63acd5e4da1c0c6b5c3e0b9f762e2f
ebdd6d67fbb67e38b7fe9ec35dda1ef2218ec839dcadece4c221b75e240cfba41bf93d87b83009483962b1a2527e70f8eef0f3e1d9e1e9
613cfc8405dfe3f36b7cfe07f09f0eaf0c1f0fe1f0a89fbf3584a438c74de7794f62ed00c0bc7c6aa9d9ed8c0afe167cae673647a23839
6f735d8ee17926986660dffa96cd1554174ee5cf6e7e3ca7fc7c4e1db763f935c5931d3f7ec84394a7094e101ad20388e725fe1b295508
edcde106ffbe80e06661be80f07e3461fe7cf8f2f0b22561b98bd4de6772cb7da62122d4a02695d0dc184188647db71d9447603b619bc7
4e6e27ddbfe7de3d33dcafe77ebafb805379728a54266efc35d81d3732229b36384161e4be5c782765e74a8808db4081316c25a0c1ee20
40589f8b979c034ed50b098021ee10a1de1232e6045bd6b46177dd4a722c7e7831c2b1f2bb9888b04344f53c5ae9491b111b7617114c57
8b1654532ee9e544c4eb98836c1a65d95a214b18149b43ffbf20f3f60b326ba48144cbbd478e03a3cbbd468e4da4436a80f54cc3b718b9
61d73a753cc30092dddbf06e1c84a89117fa3c2b80190ff260885ed99a7bdb259a43ee2fea5d994a91ea324f43b6b7a974c3de0395b60c
342aed78789b4a9989a2dcc7898402ea2295fa2b56292a11143d01c54ea3d3067c0f94dab1d068b5e7e26d6ab516bb328fcd6dba79975e
6f5df4b826bd0a98af36986bf4ba81ef835e5b165abd765cbc55af02bde618b363c3e512b5debaca784d6acd09794a6631dfa87503df07
b5b62cb46aedb878ab5a330bdbc45655bde47a2af42aef4eafd6419848e1e76344351a50ccdbfd0c6afac839e515b2d3c2c4973ff8addb
f94e5cf7fef45fc91ca4292d39f7542921d7a5e4d4e2aab470cc32b1cf5fe7c559e14ac150c18ecb38d4477136cf25ce9cd979c3d55576
6a915b6f2fc47630371079822eefe31e0d8e8fb9a4a263b73392f69225d69e0e769d83d4b8d0b1529de0229608b60ca68dc1461aa92523
f0d73d3cfba7edea00fa2164e9378611c414c329151eaeea035bb62dcd7e973f7f461a7e5776cdce8d9c37323aa53434372a6cf80bbbd2
081b19ef7e40b2d7d4d9edd8c0ed2e6a00fdc7eef94e26269d00a575ae8204365689d67617a773ef3a96801309152d714e5b96b144645b
3c3a4f3ca216654b9dfd688894320e8697c92594a7c17054dfd0bdcc638caa1295bd5a719cc855313c4f2c8f67182780e0e863858f1512
8a35b05158237b03f5c479eb307bbbf229e0a5ba62f7a1948377610b9978e5158fca2b1e1c0bf9e422c987f8716ee0f8179bfb04f1c1a6
2a021e137f856bb946e49c1d3c902cbd3b30e4a6c27ba69ca5134f7010a8626ea29620e28d1c71ec80170e6d32c7939c20192eace7938a
1d55a550de840b963b76d43d224385e338db5e7907198e24d63ace21cf9303e1b528a4a5753e1f218da0c62d8752256b5ece4de8c10651
92780fc4399df18af3c2daea92f20437d2b8e0ae246f3879c9eafc8cc3ff75de2773be9b65796fe06f8a04c3d966ad35cf74f2f70e6de6
241c7b22dae8cc57866bcffc16251135ce624036be68abb27d937398c50691443342bcd07ed5c1cf624e9e77640c2fa8ba6374b35ae0e2
d56636ea60919adca24581464d9e48e4278ecf4c8d71aa152c921e65f3382c97c73c8f67317ad4ba4b42a1d3aa5004c4a08673b6c4cccc
acd089169327cc6e72880c8bd5aa24b51998f282085c7c31f2a809ec118f719200012db33da101195ee0fed59c854e8447b3ed6303c060
ce42a72b78cae8a7e744356f810103149be1413b5388739f0c4ecddf4c34a753ce6d7c307ba0e360b9e95d13426a51e317a22ab1b8997e
d84946c642f970c2050588ad87926a0c748bca8698541ba529ccc13b4d103f47ff38e56c40a436f2af70d59187a5c6d968213b28d7c649
f0c4ca43c86450950103d4b0e7bf566c0a95e75a1900c8b4a49aeb79788eb98329004ee2ed5e3f194b255ae794d7732575f0cb0e86a4d4
2d70b34903f71dfb8bf2bc2e83da9e18eecc0c2fdebbebfa0608d791e6384fad1dc72778d433a33e3f8cf9fef0dde1cde1c7c378f81a9f
6e6ca074a616ba869eb047d40e2e31b0795e5c462c135ea97335ddc5655305e0b3078f5f1acfcb0ce0100e09ff2ffdef0274e38dee2cbf
6828f5e775c6f402927d6db3a86778ef3ff1e9957d7a3d53f0eeea124874359681064079ba5af0c5892dbad6a8fa75b366eeda67c079c2
fb5e08116af15e2b75c0594de2cd4f501e5a1c0ec6d3c1f9626811973961029a9c0ddf799604c6388f0c8f6715b290588f67217f616024
1aeaf2464ecb2bcee65374b685658f15870727e423f31eebce85d77b3871efe8c0d982d33c44dba3a1ba32f6e3f4183b0e79ed63e1b091
c78ade989c18d273e957233e7a1ca3a1df999d8a6a675c4f07cf82421be8a80ebc8a728bc306ede4d1e1abf49a9d1b493774747a69a86e
b4d870d8687ccf3e8ed9fbf06ff0333f510a656e6473747265616d0a656e646f626a0a31322030206f626a0a343032380a656e646f626a
0a31302030206f626a0a5b205d0a656e646f626a0a31372030206f626a0a3c3c202f4c656e677468312037363234202f4c656e67746820
35323530202f46696c746572202f466c6174654465636f6465203e3e0a73747265616d0a789cdd387b7c94c5b567bef3cdbedfd9cd3b9b
0d61799410424280f05c3010481082441a1e916c5e84471e2401a410082d16aa58451421a82481a255c008547944448d956781025a8bd8
5263b5ed42b95ea4b886d97bbe2f01d17b7b7fb77fdcdffdfdee373b8f3373e6cc3967ce3933b3c000c04e850cb6accc71e3210c220058
2cf58667e54e99e64dec3785e04194c3b3a63d38b6f78ed1ef024893687cd39469035217242efa82e00e82a71757f8abd92efd25001c41
f0ade225759e510707ea8838cd97f2caaae756cc92165c07d0aca3f15d73fdb5d5a0a5045a1a07d3dc85cbca3ee9f8d328827d00c675e5
a5fe127de53b73001cc7687c703975585ed5d2dc3002a1677945ddc3658b0ccb08ee49f0ec8555c5fec8b33133096e2038a5c2ff70b566
b6a63fc1ad047b2afd15a549a631d5049f25794475556d9d3418a3015c04436e754d69f50a7fcbe724ba8778781114dd98a0eb93084208
57fb946c8064180152e6f849796059e8afab8448d2217da110c0dd968abda0b4a61274ddf3188d496aada37aa58ab91e5e020e469a41b2
8576ab737787eaa96c84fff60b1dffb6fc973f466b72d2be191c2497449c287bc1083613e4a0a4c81b4efd8cfa6d2af70ae7a8f2cc296b
286b55c9f4a48f3b7a790a9e06a7aa971ff96bfc45b0c65f5351096b8a6afcf3604db1bfb296caf2d21a2a97d52c8435734baba83db7a6
7401ac29f757124e796911f52cf057fa61cd427f95472949bf6b2afc75e5b0a67281d25335d75f016b6a165712665d59e55c2acb15faf7
ecc13d92cacf486d0ac73c8d6f21d0dd55e3875026394828a30e51234b92dc2dc3dd2fb76c5c0978282dd1388593356a2bd8a7dfc1c1ee
1cab4a0eb08220a6c2321452ed24bd29faf2c040180f13a114e6c142a886252a87777afdd4bb002aa126140a7d1a3a1f7a2bf466686768
7ba825d41c6afa2e4777bf2ebb5cd10de9556a5d59e167607756241adf9d951d9bd89d959df3532e55771068fdaeac505d4079216533e5
4acae42f60a55c43790965072b6369ac0cdae02fa48a91d002418c27cb90a08c7a957a17cb83008d1711e66af91196477585bc03241a5f
259f56ec88a541112ca29657dec1dae0107c46b357b3c7f9043e53c156a55268dde4efb2eb3c43ca807cb9421e29ef9557cb7b0963b15c
26af86562a33a473f273f272f98cbc1cf215ced824252b7c4023cb6689d02835b24c16c532a5d3f036f15cc646b346369c9fe427e1225c
64b984b90b964a06f63efb920d60f96c2fcdba0937593c41e9523abbc63e278e37c339cce7066884279883a036384d7c7f065f422d3949
193cc12f4afdf8457817aec087d40f309f291e1387fdf9454ad7c9cfe79366ae30895fd438b5097299740b02ec27d24ee9164b64122507
8b276d3e84a7e542f97df967344ada6112a6613c8ea572b682c12fb246e2e28aa68c2d233c252da77502d2bbd20192f1085c26b9687569
b6b45c6a84cb6c0f3b441c033cc2f6c885da2239061a358d723e5c537403e7a4d3a48f5c551f8fc1639a817053d6c0759cc40ae597148d
8197bfcd802568b3350ed8c4b2b53f21490087c072b27080e3144adeee4a84a5d3c4c126b937be40bc4b52fd1dbdb165705acac022784e
4d1bd901d80807a0567112ecf58656c365941824796cad92776249ab6f6abee7d88c84fe49df033d36ada715725bcdcb3c0742a1dc7c39
86cf68e5b1ade8d5b5cadec42bff6cf04affa49cdc7ccf01d6675c6637d9718599d4392d9f9a0a44ddd43f2e531d53566de55efa4d2c6c
f514977b1eb53d9a38ec515be9b0feaa6552f4261f56bc9dce0b2ed3ae1ba08fcf857b64690f7f5c0b7bf4ba484d3442246746db85f640
7b6760200c08a4de085c08a48425d8b9dd9b60a7749de589ed6c16ab6479df049801df9bc03413be49173789721b39f049a2ac877e3ebb
7c5473443a0a8feb38d36238a0c1d61948bdd04e446f046e04525eb719738dd5466405612a5d25b5c9d6db35d2dcdb5bc42a7e515c167f
15976fafeda22b8d24ba08496fc0db12236ab2ad33a7d59a97d36acb9b957f9024f40d9da1329cc25eb7711fcfe5a8126e63ab1562c17e
8a0ef2439fcaef91e519c00bcffb7ea0d74054bc117e1b7146b3cd72d6ee39157f32765be209fb6613244660a4596f368e8c47b373782f
62bdfd462035d59e9191418b74dce8bc11b0fdf9dacd6b198e0cbb2323c5376940ef919e9109a37bdfefb93fa1c0539050e9a94c58e959
9950dd7bbd677dc2f39ee713767b7627bce97933c195ea4e89bfcfed8b9fe6ce8d2f7617c63fe26e88dfe87e32bec5dd14bfd7dd1a6f2b
6005ac87c6e50c4f4b1d3c9279ed09e91696d8a357faa09e0969e9837a25f6d068eda3589a47965aaa17cd9a5afa289b279e9db07ff59e
8f9895f538ffd39fd7fefac1da2fea281498d9ad49d999f76fa8e8bbf6f6ea9d6505275bde3b10fbe094e464668f8dfbbb1271429ff287
344e3acb6361a82f46d36c3a6b8566d766eb8998ed9167a370acc31769d29ba36d71caee913128d2073a54e953f64f7117ba2556e04d1d
dcc516b18cc4582ad86d400cf776d10620d42e5b565bbb7cf9f2ecfd8bcf3083b87966f1fe6cd1c8ca3e7fa9a9e9a55f6edbf64be96251
817843dca6f44641510bb1437baeec553aed55140cf145c3297641b69c325f3034dbe5e608daa268ed5833388747db3adb5389299bc292
b871cdf615b1352766550cb1c5ecc92cf18e1e87d849871e85b184d470397de2c6299b5e7c7153deb3bebcdd3f14e7c4cb6c3a1b90ff8a
3c527c9c9af2eaf3cfbf9a3a505c8a8f6743988bd29078c57a6693ae2c1ae584547515d10c672df6667e56b7d972826d47a70c66c9173d
d6385cd195622d8a3976dce808d83abed51553ec5cdd4f8597b4d470bc57793b59b9d89cbd77f119719319ced4ed6b517457b36c19b649
f95f075a8a67b3890c294d2ce83cae684fc95d76cd0db212db62e03e5f02b87ec7f4a77417f83613fb30729be38469736c8c4bd2b9cc90
2999adc36355fedac970bb14d7a16aee5a46cafe0171a3e314265d096ef65da5798077591e37741e363fb77ed1d5fa556295382b5e6539
ac07d3b191e2c9a585e53fb64969652b57de9729022903593a8ba0e83f4cbcb3b1ac7e7125a8271c480ed221d229ded767d15c867df276
8933335de775b6ce11015569141e7cfa5c7da1be5adfa097bf8d10bb58996238a2512efca6896e381f939d1045be8c281a60bdcfa9e312
e23ed8aed770266b74cc0c366357d8b12b7e3b22d01141ce5a18855172148fd20cd74ec489da9958a0c9d7ced156e15cbe0a97c84b34cb
b5ebf011792d7f4cbb0d9fe19b345bb507313a5c0ae7c374995216cfd64d970a7473a542a9942f95aaa5257c856eadb48e3fa67b4a7a96
6fd5b90a549659829e2530956bf61c6bba7d599a2c268b1cd1a87176bec6b692b9b7b0f3a27f970cdd5a71bf01db25856d8a6f5de18c14
c10b79356fe0dd8a5095a0717e1d806ee91db4f31190ee8b242532c325c759d709db3e0b93cc9065379badb6484503aae78e50ac518961
29fb0ba31aa2bacc31dd3ea857ef2e4f5174ec521c389c3bc466b3cd3521b9ba41d1f8d4d72bdf392ebd7c7b7a15dbbaa1323ab1f7ee2d
b73f920b7714155ceb8eabbc377161243efafbc234cd0e68369d706c8ed48fb54ec2b1aee10a0b9d77dcf45acaeb43c246476585e13d71
8e7c00ef581a95d2db4bebeb972e5eb16231a9709c382cfe203e110759162e7fa5b9f915253310c74480d231369439290d55ed4b4ce70f
11278a8f0ef7c57deba3272c9bd927b82f8efcd3a77aea3d11cdd6d1f17d37f5de5545772c0bbb8751c9a1e8e49e70b683c25b2db96acb
edfd1ac38e7b82190e51e29beaa5ea5e618078b353440bb75bf41a2de051e3d39613fa7d5a83c60c3a9b4339d4c2f2e83cd3858e0e9dd1
ae7a03b1d649e6ab1c33ecfbdb1481810105031e7d46e167dcbe7ac70ffae08070d76bbfb8dd29171ea82c45aed8c87c8aa445b46e6ff8
d8e7339b248b71b03bdecd355a9d9ecb86c1f1f16eafc1e88e975d6a94759e725d88dc6697b77929caf6711b8cf1315a9816f380c5a9cd
ed717f1fe524bc10e8206ed4a3b02bee7ea5c45d07b916dd4db436cb557b44467735a3c75e9d91b102df437abddea0371a4d46b3deca13
a34dd1e6684ba4354997ac4f36241b934dc9e6be9e0cdd70fd70c370e330d330738e3edb906dcc364d302f352d351fd21dd21f321c321e
321d327b2d1a8bd6a2b3e82d06b371887974df397df505704fb88f6772b8cb2927d0096457e318edd900d25dfaa0c1a43339a2f6833965
c539fed12cec88b825825557eb175ca99b377f62c5e8bf1fbdd159fc7b3a05aea7a4a4a5f74b36ea139b5ed9b73f3191d9060d1a969132
c0ac73b7fc62ef2eb7a2d718dacf1dfc057adc57f8a22d5c67c5663bdba76b0683cea897f46464368725cf79cf4525a7d5ae140ef5ca62
51ae2c23da3b47b4b73bd45b457be006c53f72cf14e68b1c0da3c394e7045a6d567bae948bb9ae42a910bb045582a1337c3835e856604f
b4a7d9a534b6483c316af60171fac26b7bf7f217c43b2110dec94342f0da0576896eb9a388e326f2538d5c4867ea705f8c95ee82a64bce
b3119b6d6c9f43234198d56ccfa2a8618bee3abfd4f3614047eab73123a6413957c3c8fa9c12e939a24bb9774fb45e4dac5cb2d8c3b328
6a28416aeaaf2ade39c9f64bbbaa6789abc96b97c624f6dab545eafb4d538b1a3718d0b1caa3881f0d9cf7f5c63899cb521c3d35940a25
0d68581cbdfcbc12c21fb886d3659bd393df7661ff369a2815b44764a817d46eabe357b5365d77e65767f460be15ca7bcc46c923a5483e
69ba542635484f4a4d526b773a4ae92ca53faae93aa5284e2f3d331ac817c2582c46cbbda027eb877de5c130886560869ca21b0fe3199d
1a72a1ee61cd5ab60ed7f2759a4df4c6d8825be46778a3e6257c9d1dc49eb44f616afc4f88a7f01fc1c2d95c912596ca859d41d47cd3d4
f52697666c19bafe45c31ceb88af205ea73e21cfadb45cba53ffa3a6738ff5b4be54fdefe5ee7f03f400ae10a4185bc23f6abe4eb59efe
4fefed087ab196a9cfd197291f271d6f80cff952689382d0c6cb209f1f825df211b59ecd2b205f3a02bbf803b04badf7525ed585835f12
debb301fcf430cf537f121ca8eddfdfac30ca8671a368d4eb7bf48a3a4b9d271e9cfd8172bb10d6fc93be4db7c1c6f57b98bc0c9d00f1e
a6d7ba4491d9a748c3af4a0954cbea7b7ea6f22e91f5245c8afa7f80d266104e50579b5e2f6c7c771bc1c32677b7e57bda1c225965775b
036ef653b80faaa01a96410dcc83b9500e75e0813e500c7da94e85144a69d42a220c0f8c259c3a7adbd5117629f8a10292a8772254127e
32b5c6c0424a1e78e02ead5a152aa5ba94e62ca1b284300dff835507df5d358f565a426bcda7399584adf0e1a739ffda8a99d49a4ff3a6
c362c228265cbf4aad549de15725f210954a2aab09a788e8ce233c0fcdafa2d5fdead8f7e94c53a9d4124755845ff24f463d77c7a7ab5c
d512ad2a75a554e22d0dd2bf33efceacfe776775fd43475f682559c17ff545a8364f81906eb77d695732611c8c872c98009360324c81a9
a49369f020d1ff21e4c38c316ba406d60b0420f3d2498bac274b8061c0594f08129408f476a4fb6a575f0f154f6923f3a8e3f170984a37
7186147c94d1588a96c862c04d65b4da13a596916a19a196e16ae9624e0aea9cb9544869230b53db0eb5b4320bd4d3b855859436d203cd
04eba9cfacf699e128c8ccc48ce4595c1d412a1ba8cfc80cd08bfa9411a4d2477d4a0f32bd3a53a7965af230a5546668f63e9bccc78431
8d2a17574b59c542552249ed616a09be503d8646a110d8f94d12ef14f84d1206057e7d2b8b7f5d8fb7b2f01f41bc29f02b813704fefb61
fc52e0bf09bc2ef0ef6ebc26f06ac0c0af0a0c1830e093fff65703ff5b2afed5807f09e2171bc2f917023f0fe29f83f819019f09ec10f8
a9c03f09bc22f08f02ff20f093205efe38925f2ec18f23f152939b5f2ac1df7fe4e5bf0fe2475efcdd392fff5d103ffcc0c93f0cc70f2e
daf8074ebc68c30be78dfc8207cf1bf1b784f1db209e23fae7bc78f669133f9b88677ee3e4677ae16f4e3bf86f9c78da81a768f8541c9e
74e289e387f90981c78f15f0e387f178837ccc177adfcb8f15e0319ffcbe177f2df0bd126c7fd2c6db05be1b8bef087c5be0d1b786f1a3
417c6b770c7f6b181e79339a1f49c537dbecfccd686c3b6ce56d763c7cc8c40f5bf190090fd26207051e10f8860b5f77e0af04ee17b84f
e0de087c2d0a5bc3f155a2f36a10f750b52788bb097f770ceea26a573dbe22f0e55ef84b812f097c51e04e81bf30e00e81db5b2c7cbbc0
160bb6f8e466525473109b684a931bb751b52d882f90f02fc4e2f3029fdb7a983f27706b6301df7a18b736c88d4f7879630136fae42d02
3793756c16f86c326ea2899bdcbe103e43539ff1e0d326dc485d1b73f029aa9e12b881f4b0211c9fb4e1135efcb9c0c705ae17f898c047
05fe4ce0bab55ebe4ee05a2ffe54e02302d7a4e24f36e18f05ae16d81085ab0cb85260bdc015029707f147415c2670e9929d7ca9c0253b
71715d0c5f1cc4ba18ac0d624d3d2e12585d95c4ab92b0328815415c18c40502e70b9c27b0bcd8c4cb5371aec0b2542c2d31f052812506
2cf1c9c545065e6cc22203fa0b5ddcbf090b999d17ba708e011f1258207036c1b305ce9a19c367099c49d0cc189c21303f883f14389d60
5f68bac00705e6b9719a131f981ac51f08e2541a981a85b953a2786e10a74cb6f3295138d98ef7bb71528e934f72614eb69de738317ba2
8567db71a2052704316bbc9367b970bc13c70531f33e0bcfb4e27d161c3bc6cbc706710cd11ce345df682bf7091c3dcac2475b71940547
8e30f391e138c28cc34b7098c00c270e1538240c07a747f3c15e4c1fe4e4e9d1987e541e6430f3414e1cd420a7a59a789a13d37c72aa09
07a6ece40305a610fd949d38c084c961d83f6918ef1fc4249797270dc37e25f88312ec2bb08f0b7b47d8796f37f6f2a0d78d3d134901fd
7aba31d18e3dc0cc7b0431c18a093ed9e3c47803badd18171bc5e3bc186b0de3b151187b8062c60639c68cd151393cba1ea368d1a81c8c
141861c7705a2d3c882eea7379d1598261767408b4136c17682b41abc5c6ad61683d2a5b6c686990cd34620ea229158d249a311c8d0db2
81ae613e592f5027502b50c30d5c23901b90fb6439885842f7413b9704452f3367760433b203ace491c759bfff1f1ffc5f33f0bff8c5c1
7f0002043c6b0a656e6473747265616d0a656e646f626a0a31362030206f626a0a3c3c202f4c656e677468203538202f46696c74657220
2f466c6174654465636f6465203e3e0a73747265616d0a789c6360a010301354c142b459ac60928d4417b0337030703270317043f93c40
cccbc00724f91904180481b410833083088328831800183000fb0a656e6473747265616d0a656e646f626a0a31392030206f626a0a3c3c
202f4c656e67746820333131202f46696c746572202f466c6174654465636f6465203e3e0a73747265616d0a789c5d524d6f83300cbde7
57f8d81d2a0a6d41480869ea2e1cf6a1b19daa1e203155a411a2400ffcfb2571c2a44582a7f7fc6c8c9de4d2bc344a2e907c9889b7b8c0
209530384f0fc3117abc4bc5d20c84e44b60fecdc74eb3c426b7ebbce0d8a861625505c9a70dce8b5961f72ca61e9f180024ef46a091ea
0ebbef4b4b52fbd0fa0747540b1c585d83c0c1967bedf45b3722243e79df081b97cbbab7697f8eaf5523649ea7d4129f04cebae3683a75
47561deca9a11aeca9192af12f5e52563f6cf6ccd909ae8437271f532f3bb8127af9947bd9c19590e492e432c825c93915c9cf24479a11
1c094ed1432954290f95f258a927990779a314454f8bf00f81e64354bda9a0cf16a1f548a98b82ba28ced17373c38b637273744bdf96c4
1fc6d8fdf89be117e35622156e97474fda65b9e717cfc9aa5c0a656e6473747265616d0a656e646f626a0a31342030206f626a0a3c3c20
2f54797065202f466f6e74202f53756274797065202f434944466f6e745479706532202f42617365466f6e74202f45564943414f2b4465
6a61567553616e732d426f6c640a2f43494453797374656d496e666f203c3c202f5265676973747279202841646f626529202f4f726465
72696e6720284964656e7469747929202f537570706c656d656e742030203e3e0a2f466f6e7444657363726970746f7220313320302052
202f5720313820302052202f434944546f4749444d617020313620302052203e3e0a656e646f626a0a31352030206f626a0a3c3c202f54
797065202f466f6e74202f53756274797065202f5479706530202f42617365466f6e74202f45564943414f2b44656a61567553616e732d
426f6c640a2f456e636f64696e67202f4964656e746974792d48202f44657363656e64616e74466f6e7473205b20313420302052205d20
2f546f556e69636f646520313920302052203e3e0a656e646f626a0a31332030206f626a0a3c3c202f54797065202f466f6e7444657363
726970746f72202f466f6e744e616d65202f45564943414f2b44656a61567553616e732d426f6c64202f466c6167732033320a2f466f6e
7442426f78205b202d31303730202d34313620313937362031313735205d202f417363656e7420393239202f44657363656e74202d3233
36202f43617048656967687420300a2f584865696768742030202f4974616c6963416e676c652030202f5374656d562030202f466f6e74
46696c653220313720302052202f4d6178576964746820373136203e3e0a656e646f626a0a31382030206f626a0a5b203332205b203334
38205d203439205b20363936205d203730205b20363833205d203733205b20333732205d203937205b2036373520373136203539332037
313620363738205d0a313035205b20333433205d20313037205b2036363520333433205d20313130205b20373132203638372037313620
5d20313134205b2034393320353935203437382037313220363532205d0a5d0a656e646f626a0a32342030206f626a0a3c3c202f4c656e
677468312033383630202f4c656e6774682032363236202f46696c746572202f466c6174654465636f6465203e3e0a73747265616d0a78
9cb5567b7054d519ffcefdceb9777377b3ec86252410c286cd066208e046f282d105f220042198109340209bec6e42c88bbc0306504229
a082a82b20166d91225abbb5b662d18e1615a6347f3840a9e338ea48c7e90c76ac833acd98b3fdee26a4cad499fa87e7ec39e7fb7de77b
9e7b1e0b0c00ecd471b015171416c164d000581271e38bcb569717f4aeec217c07616771f9daa5b34fdc758e7005e1d0eaf2f99ee6a4b6
0b84af13ae6c68f575c09f957d004a01e19686de6ee753c7c36984c324d312ec686cfdb4f95f6f9133c3fefe465f570779237ffc14614b
63cb40f0efc53da3845f06c0e6a680cfaf55fe290ec0e4a4f9ec2662589e13ef10ae229cdad4dadd3fed2b207ba6fb0827b5b437f8e0aa
e2217c98b0add5d7df812d6a0c61c3beb3cdd71a9897bdac9330c5acd83adabbba23cfc3c700fadf687e414767a0e3d023ca2c00b38962
b080b13616182b0a210435ca339a0e49900caca06865059868f5a84422e3b263f387e0317044e7b7fa3a7df530e4eb6c6d83a1fa4edf26
186af0b57551df14e8a47ea0b305861a03ed4437760636c35093af8d649a02f5c4d9ec6bf3c1508bafdd69f4dd64a1d5d7dd04436d9b0d
4e7ba3af15863a7bda48b23bd8d6487d9361ff9688c6e3e25676100480c81247209b251b63a410af4250a15556cc2a2272b3c2c7739828
65c1423f718aa15f7548073baab5b24fbe2383e32d299a3940012116c51ce6d0a8d2c8680d97437f34a662638cbc1e7978c2c6d83acf99
586d4306481ea299f48f37959cbad87c380317a9bd09a7e1bc721e9e861caaf7c01fd95e2593664ec21e784f28f0321c6379ccc1f268f6
92ea5007c46e718ae6d793ee72b2f21e3c4a360d4b676087b24d2983209c17c370946a7b94ff39fc81ed822b70182e2acbe12bd8851570
80ea51e8e220ae301da49201270c4f5401eaa9cdc04c71255a3f871db00d2ae0847a4675b04bd1a84fb273745afe41315fc2f5b885348e
c129b69bbbf829be1c0e8cc58b757040d9c58ef2ba68dd460bd207c7781d3bad3ae06d2356e29451a441788d5a1f0cb3456c37eea5c8b6
1911882b30acade0f3c7a2d2067121e503d45e8497201343a41fcd450dc2312548bebea24886b100d2c95a1780834e374cf9bd2a382a0c
e63a6d61c55de20f7bd754392f54a764cebd053a6d9a330c65e1d801e79948a4ac8a4f17d561911446b729ccddae8fbf6ff2e3ccb9a565
55cef06f0b0bc6ad16d61510afbc8a4803119bf8850599637bc2040aed1e6367ad8f7c2286697d5db0d3bb90092ef82c97920242089e0f
4ca498580df6a4d498a6041cdd0935f64ad39a5423132e92a7704772923936d576e3fae5b7f2ec7953f36e07dbb51bd76f5cb7c96ffe69
8f23d6020a40b389cfec53f38ca17a162b0dc75694862755acab7a1d52236f5000a954bdb9f69cdc6af7e4949c1454b9cb09f63be272ee
64591e3e559bc75cb3544d55a638e2b33cd97cc32af9f24c25eefeddf29de3f25379eecdd39be5930f32d6935ef3c5dee35f6c5c53fbd7
cbcdef7cc8f69c1c6d2f59cd7259fca7ec99fdf7ae2d2ebebaa16160478b7cfff0bb97328d139743a77798beaf1906bd0526a10a4da8f4
9d8840419d2ecc608e116056a8313aeae67c538d429bb54857394d9b634c1aad840aaa2266d0899f6eb1ddb876e31aa54a99c7d15acc5f
7ccdf3cdf59b4ba08dfd6ce36d62a035c9c294c959cc3539055dace6629ff241df45f9c1c5fda3aefd17c595d1638a7f2443393f9a4731
472ed2211ee67514d076ef729e2138531866280a336750467a468cce8cc1ac6b262d26c364d24406ad1dcf40ae696abeae71c64db05d28
3a056fec489305548bed7269d8425fc55ab1ae346c333a3b7d9f57418dbc915bfd17fbd4f15ca68e7f4ed37fb388d2dfc5d5b35e327106
b5decb0254a67121e2593c7788db209da529b3d12dd2d474537a4c3a45990fb92c9b2f140bd51c2d4f5f645e0125ac98178a42b5582bd1
57982ba192ad15956a95766f4c10822ca8043120826a634c500f9a7ba187f5f21ed1a3f6e9bde6bdb887df2f86d43dda4f62f6e8fbcc8f
43881dc627f82111529fd09f8567cdafe02bea8a181ea3a9fa349ec63278ba96a667b35c9eab65eb45ac849768457a15dbc49ab426bd8f
f7eb0ff2bdda6efd287b8a1fd142fa19ed8236b396d5b2970ec6b2da18c585cc15c3b226333e3c2ad9a24b9f49ff87ca6cf9fea5abac8f
d78dbe3ffa367b4d1628ab9438b9993d3e7ec72bd547cebd5e94b071d2e22f61a6297a41bfbbddfae1cdf1eb8f46b3628b4cfba32774e2
ada17ba755ce0088ddfdf547ff7e35b608fcf4dfe2db45e3c6fd65983f4d6d3f5dee59b05ee590c35f3376cc441984abec77d138345c05
19d0446f850236f01aff1cd0aac48ebfbf1ad4187703a7f79e2d88be4706cd209ed018ad8095158dd30869f42f668ce6dfa20524b0ade3
b40aa9ec102ca37bb30306a01336412379ef0627bd530d74573ac1030ba86611554f124e584a32dd7483769374007cd00a73895b026d24
3f8fa825d042d5496fd54d5b5d5114a031403abdd4fb4952ff3fbc664f78ad204fbde4ab9974da48da88c3473a3fcc630151cda457093d
24d140b2bea8b54054c317cdc84956daa8ef20997ab2bb89e49ca4df4ede7dd1b95bed9447ad74c1ea71f92dc40d7c8f8cf316a9ca6884
5d84dba35e3d1467162cfc8ef64dddcc5b74e91f44e44b6adb6957fcafa245f7a8425ffa36da53ebce283bbd91df480cbbf1d71e7c3184
bfb2e20b5d56f182079f9778da8dcf59f1941b7f19c29323f8ec089e90f88b7cfcb9c4673cf8f4f172f174088fdfbd441c2fc79f79f029
071e0be1933a1e9578240e0f0fe213673124f13192786c101f9578e891627168101f29c68307a68b83120f4cc787253e24f14189fb25ee
db9b2cf649dc9b8c3ff5e01e8943f1b84be20312ef97b853e20e89db250e96bac5a01fef93b8cd8e5b07ce8aad1207fa6bc5c0591cd8c9
fbfbdca2bf16fbbdbccf8dbd127b42d8edc72e2b766e718b4e3f6ee988135bdcd81187ed1456fb08b67923125b25b648dc1c8fcd9bf245
b31f37918f4df9d8b4ca2c9a12b03168158d1e0c5a31e0473fa9f943d820b1de6711f5127d16acdb9828eafcb871834d6c4cc40d36acd5
71fdba58b15ee2ba58ac218d9a1056575945f51cacb2e2bd2358b9f6aca894b8b6a256ac3d8b6b77f28a72b7a8a8c50a2f2f77e33d12d7
94cd136b2496cdc3d514c46a07ae32e3dd14d5dd4b70250d2b2596aeb08b5237aeb06389c4e5c576b15c62b11d8b24164a2c90b86ce9a0
582671e9202e91e81dc1bb46f0ce115c9cbd542c96b8e802e613955f8e79d2db81b9839843309b678aeca5b850e21d12b3f2d133820b2c
385f62a6c4b91233683ae376bccd86e96813e92e9c938cb3d3ac62b61fd3ace866ba707b30d592205207d125f2854be22c42b3ce620ac9
a74c47e74cb3704e42fa47f486f7289f69c6e4184cf6f219364c22f1a4104e0fe1b444b798e6c7c4843891e8c684389c1aef16539760bc
1ba74874489c3c8271f6441127d14e56ed8968933849a2952c5843184b0e6307d162b6084b029a2da84b34d19429842a89ab12056521f2
9113e29988367a5975a12420d391793924213bc3fcbb1f62193f6e811fd9fe0f2d33fe03b6b1484a0a656e6473747265616d0a656e646f
626a0a32332030206f626a0a3c3c202f4c656e677468203139202f46696c746572202f466c6174654465636f6465203e3e0a7374726561
6d0a789c6360185480858195c636b00100046900100a656e6473747265616d0a656e646f626a0a32362030206f626a0a3c3c202f4c656e
67746820323430202f46696c746572202f466c6174654465636f6465203e3e0a73747265616d0a789c5d50b16ec32010ddf98a1b932122
b12a39838554258b87a655dd4e55070c87855403c278f0dff7c0a9abf62440efde7bc7d3f14b7b6d9d4dc05fa2571d2630d6e988939fa3
42e871b08e9d2ad056a53b2ab71a65609cccdd32251c5b673c6b1ae0af444e292eb07bd4bec73d0300fe1c3546eb06d8bd5fbab5d5cd21
7ce1882ec19109011a0d8d7b92e12647045ecc8756136fd37220dbafe26d090855c1a73592f21aa7201546e90664cd914a4063a80443a7
fff1d5eaeacd267fa8494ecf59c0c71ff899d9fa5c607d67ebdca6a93ffefc41dec6965ecd3152f0b2b2923867b50eb7ad061fb22b9f6f
051e79aa0a656e6473747265616d0a656e646f626a0a32312030206f626a0a3c3c202f54797065202f466f6e74202f5375627479706520
2f434944466f6e745479706532202f42617365466f6e74202f4743575844562b44656a61567553616e732d4f626c697175650a2f434944
53797374656d496e666f203c3c202f5265676973747279202841646f626529202f4f72646572696e6720284964656e7469747929202f53
7570706c656d656e742030203e3e0a2f466f6e7444657363726970746f7220323020302052202f5720323520302052202f434944546f47
49444d617020323320302052203e3e0a656e646f626a0a32322030206f626a0a3c3c202f54797065202f466f6e74202f53756274797065
202f5479706530202f42617365466f6e74202f4743575844562b44656a61567553616e732d4f626c697175650a2f456e636f64696e6720
2f4964656e746974792d48202f44657363656e64616e74466f6e7473205b20323120302052205d202f546f556e69636f64652032362030
2052203e3e0a656e646f626a0a32302030206f626a0a3c3c202f54797065202f466f6e7444657363726970746f72202f466f6e744e616d
65202f4743575844562b44656a61567553616e732d4f626c69717565202f466c6167732039360a2f466f6e7442426f78205b202d313031
36202d33353120313636302031303638205d202f417363656e7420393239202f44657363656e74202d323336202f436170486569676874
20300a2f584865696768742030202f4974616c6963416e676c652030202f5374656d562030202f466f6e7446696c653220323420302052
202f4d6178576964746820373735203e3e0a656e646f626a0a32352030206f626a0a5b203731205b2037373520373532205d2031323020
5b20353932205d205d0a656e646f626a0a33312030206f626a0a3c3c202f4c656e67746831203130303536202f4c656e67746820363936
34202f46696c746572202f466c6174654465636f6465203e3e0a73747265616d0a789cd539795c14579adfabafaafaa4e96eba91a3691a
9a06f18280a87876f03e624824464c7440016f40f188210e9e205183c68889f1e8281a438c8bc635a0c6d548341965c719353bd924a389
187310e3644c9c4178ec57d56030b33bb3fbc7fef6b75dfdd53bea7bdffbeef75e153000b0d24d04d7c861c34740107403603da8377864
faa3133c5bbbd7537b18c1ea91139e488bab1a720640f8809e573ef270c628774ed623009845cf173e3a21216976ea260f11bb4ccf274e
9f975d38bca6c7e30012d163c6e98b17ba6056442a80dc8fda3caf70c6bcf9bd17cf06d0521bde9c915d54081aba40db486de38cb94bf3
8e96ecf99edab701ec3033373b473be5bd5280086a439f99d411b04bf32280b32bb56366ce5bf8cc9e378223a83d8eda3be7164ccf2ed8
317b1db59b147ee6653f53286e95e70344923ce0cacf9e971bf7dda013d4ce217e2e1716142d6c199ff51c4094c2cfbec205b98503347f
a26ad43592612628ba3282ff27500b215eed53400fbd602008c3468ccb00d3dcec85f910423aa55f5b1bc0fd9a8a3d2777413e68dbc731
7a26a8a596ca8b2ae635680119cc34e22df85ffdb57dfa0f9e0f53ef97da3e5540c1f797f79f67ffdde18ca408223d8483039c104ded10
30a8522bfa4368d70bfd241537a2fdd97ff6ebc097084f43bad291c681a81921004c10d86e8317e125b0a93678367b41f634589dbd605e
3eac9eb6207b16ac9e9e9d5f44f799b90be8be74c15c583d23b780ea3316e4ce81d533b3f3096766ee34ea99939d9f0dabe76617b8943b
d972f5bcec85336175fe1ca5a76046f63c58bd60513e612ecccb9f41f7990afd4ef6eeac0385af6bf015a4c34023683e53855949a26491
d857fc484abdf3afa3cd66b6d79ff8bb7af6e39a88eee27f8c0793fdb84a79bfde89c603fd933bb517fc5c17283eee55507de0fda15dc9
5f43d4a870c208c8a75250adc5e8b2deb73342ac82ab5b44770f8de9aadb40cf63a806ba4a804e78f1f7f1badec78bfb1b3c81681386a4
449055792a291c30d1c436d2dc20254baf10e74e7f89ff067902e10b0619512b0a824245b18fb74388f4bce139d4724587c9366e63db34
f3d8f5769c0ebefce00050fd6f35b598da16816c05c1a405240f7511ffa930081e855c980585b0383a4cf50b17f1d84bedcf26fc7c58a0
f4b75da7b8fa43db95b6fab6b7db0eb5bdd576e08fbaff3212fcf9676d7b4bafd2f403aaba0255738a64bd0828d712374033fa41d1d6a3
ed1044a0c46f2e818d40e17f164128413e41214138816277c5af9c00d1617e00374b815a384fd769a886ed6c1fb5f2a87f3ef5f884c3b0
061651cf19769e950b3da96f1fdc864b845906e7b15a04360692a917e063b2d71d96014788462ab3b1548d2c82385e3c223e2ed68a37c5
06e82b16890d629658c49271b73451da47908aef931d3f8448a86557a1088ee1d7988c27c461a209ae620356c30d9a45d1c179a8802a28
265e6cac004a8462e171ea392735c036ba0ae87903dbc92e1177c7d82ab8022fa3288c829dec0ac9751e7e8255982194906a93853ce2ff
1cd16aa0f1dba0885cfb0ad30317ba531f714f734d53ef11d853baa25eb7a18466ce802ab956b669dc348ba2b17dec0c6b9237830f2ee1
d3381f3f616b44b7b85f1c05157e0d60165410ed6dca18398f2d25d995ab58a12e2c11b358357c2d6669a611edf7158968ce23c2e32451
1e9c2058229b49a6016c0d9613a7cad30868d08c1113683c51d02c23a9010a30056653ad180ec261e889955041945479e5bed24f3472bb
f839c95cc136083f41030e23afca136f91ae1557a1087c47234b220a0c7ab8cc358267744e8df7b149ae0f32a37af6f845d365d6b86a20
bd2660a9abb6ad2d7d92182e65d6488e1af4686b448ffbf3ffeae1e73d7b8c4d9fe4aa691d3eac9deaf0ac61d43761125595167553fff0
61ea3365d21ac943ffd15935aee9335dcf9b9f77f77fde9cdbbfa73f4fd0ea4a11a8446b0dff412896adb46ef4f506ca2fc356538006d0
4a4b95de64fe746c4d50c6a43ad0b79dea9739b62650ad83b75f666352932535f52130b734253259b0dbac5ddcb1424a6f6b5fa1b874e5
aa35beca2d2f6d95ad5ff2c1376ff20137be6567af5d65f54d345f15cd57a0ce17e90dd428f3691818ac62901668be81777ea61b941c6c
b5db048dbb8f35a5b7504524b754fad6ac5a255b9bf8c0abd778ff6f6fb0f76fde64ef11d54f282d1d044e1259de8155020b851091683c
04094d897d93edee4f2e5de25cc9226d69c2615a67107a7a6d10c604268421609ab00b5688942c31a15e9dfe4e5322a9566396be532033
9a2533b7a0ad6efd4bb574e5aff3140d96b55d172bc8ab0dd005dcde20d967059f7193757d88ce11e844873d3c84e6bf431c981bef3499
6f25b268c162b62627592d66212e092c6670472b7761ddf61d3be8bf63c73da6e377efdde377994e4ae70dfc0241034d9ccc7ab3641f2f
e2a5bc8c17b10d6c297b966d50f2f0e794da2693346412af3d0d7da2e0935668c0a7d346ca0e844866305f6eb71a53acd654ef5749d29d
a6cbaa8024d891400c1485297da32c528a27d912658fe26c0c7f85e5fe868d69a9aa168b46d58e6abe524d04c8cfc53124b103767ae342
c3c231c4619144b0489298667ecdf25280cfb649a47c0766bdc0f48e2e669423cc2d636bec19636b82339e1a5b63cb788a3841c597ea2f
379d3a65b1a6b673d349ddd277acc661b6744925debc494f8813a5899a67c567a5c5e165a11aca86a16218858563212c96178515852f74
ac84d2d095612bc3573af6c3fe70cb1498e2212152fa40dfc12ca577ac3b5ad6a40c66c949a2dd266b64a0147cba651ca93139fb91d74b
7f75e999672f4ffa8ad9863f15caef5457572f619bfacfdb3a7a4965dad00b0f257df5ded37b0b23f8b724fd76b2771149df150abdbdc0
1ea42fd54596ba827cf6009f6eb3ecf0b936bb37c9ebed7be2831d4180b65047accbec405ba44e8e5794109cd121bf4e959f1440aede45
75b6a6c63b8d4de62f6f99d58bb492c8bcba1c677664b62b274a8429ccc9ec36312a3a362ec54982f421a9bab3147fe501f170c8a63dfc
22ff6aeab9d9191fcc3b79ae6eefc1a35b76ee7979c2c905451f667ec98c2fa027b27ee3673f783c671e4aaaac58bd65df92c2a2e298d8
232ed7ef0e3ff7a6e2d7b4f717abc8a7048ad015de081680018018900668d0f824862b74cca80787ac158d6a7e30906001aa604645b0cb
03eb9b922c8a5d1b2f0f6c4a225954c38a1f92713f544cdacd4047aa5190498beb12781e34c1ac3bc4b2eed8878d678f1a1f0d98c8f2d8
22f62cae6101644a1d8bc2640b05b0c56d894a41990b8ca7f02b573e6c9d2a795aae63434bf27eee635967c8423bc94239c479044cf5ba
c5308da5d41c11e6d3d87ce6f200c1072b02d66baa9c5d1c4c8f0ed09b65a7b98575b68bb9538e332bd1422632d7df5202588960320faf
f75b2788fccba2e81cec3678c02c8a353ec3d0565f8f493d9a590cbfccbf9f7a66e6e45373defacd6fde7aecb50ce94a357f313090dffa
e64ffc4797ebfc438947b76f3f1a134bdaae20ee2bd57c120393bc31413204941ac1172cfb1cc17bcd3e6379f426c77a8f315ae7087506
39302a32dc4309869ca8514d318d2d8d3fbb8fd7467b11d620346083785e3a2f93dc879dc214368545cb765bb09f5766efc5dcd1027608
e27629e9282a2958a85abb6bd75a02a61bf7eab80f2e050e383ce77326f1db5ff0567e8ba5b3f071afe28063bb5f3b7efcb5ddc784a5b5
31b1fc07fefd9353f8f7df7ec9bf5113d434b6d7a964a8fde44d33c926324cf78648160105b488942f24b2074ac8443ae868cc2d17ead5
ac9ff037695731d0a477e988e325627418261b59faf6cbf45a27094cc63029551a25cdc01aa89135e42d6418e66651fbf154eb1797186f
4d96ae4c6c5e2175577683eb48bfeb54fdba2101867a3d21a4dd38d9e7ece9b36e72ae8fdb9318628ce9e6b0c738027594bd2985074685
279a5bea9beed437a98aed8855b5954a41da49999e5e946b629293829524a386ab3b3a26a5779fa00e04f20c61ddc6bd7b376edcb797ef
5db909dafe78956f5af1e21e7ef7ee5d7eb76ad4a6552b376f5eb96a93f0feb6b2b26daf96966d9be83abcfced8b17df5e7ed8157db6e2
e3afbefab8e22ccb5eb872e54202f298152451194914a27a8c5b1319ca4a21d4a7df2bfaa03c38d267de14bcdea37138a282e8b419ed08
501d86d8ef5893bee43f76f84b707de87b61a7c24f394e45bce7ac8fd4545b4f58bfb622794c5fd5b7ad4126f21548e90dc97e2f898e65
1d62910e3e1fb77d2cf949ffc373aff17bccfc054366e187f88d71dbd9e0765f8a242f6101cc3af16916f8ed972c585dce76f1a79cc2d6
0e4f52b2cfc72c4bfa047793bd34e0f006c07641de2e6a2526825b6b6e19589f44bc5f56162f2506958bb6cccd5c01dcbdff763551286a
bb2e7ddc695dde62852dfe753934301943ed66655d6ee9bc2e2b0225ab2b729c5f38f58e719ff256869f7eca186ffb94f567cff0b5fc2c
7f9f97b1a5d2385ecb6ff02f792d1bc5c258381b55c59fe23b15795815f93e793f71524cf6e9493b553d9ded4ed08a1969e8a233c11b5d
e43a93c5551a79cc51e7aeb5acef62842e1812a0d31a22516b1b1e4bdc5db8dc9494e40f88fac63b2dc4e659d5ff2ca98aadf21323129d
899189aec4a8c4e82171de08afd31be97579a3bcd1e911e9cef4c874577a547a747a5c61dc9a883267596499ab2c6a4df4c6385fdced38
67c7d08e411d03b29c599159aeaca842676164a1ab306ab97379e472d7f2a890ce596310eb6b71a728ae104bbe9d1cd579fd09164e5e3d
b0a2e095badada2127d61e38df7a8f09af6fcd3a9a917b72f29f6f0bc979c5d38a3e3e123fae7545755ef6e9ddef9eb296acebd5ab3a2e
ae45b1fb31d255956c23ab39a09f3714eb8c81baba10fbfac0daf0ada160b58e0c31cadab011b49f20d5dc5137318d4a5c9ebd957834cb
b9dce97322f1d9e18fc42a53d31aedb588d7385a3e92f1c6eb2fbef8ba02ad2ff43f547c01dada2e141fea5f5727249cbf79f33c81f078
4e363fc1ff42d789ec9cfdc40d83f96dd7f126d930148678c3a194ad154da5016bf57516b1ae0b192f4c630d8051b6e161e616da11b7fb
14bf73cbfce3ad44af2130dc1cbe3c7c63b82f5c22e6d474dbce5d5fbb1a4ffe7c8b37c7ef487ffbecd9b7d3778c7f64ef9456fe11ebc9
e427768b2907ba77bfded070bd7bf7ea981812c8c4acacbf9bb4455c8993893fb35f5b617560b2d549daf5a65ab615bb88a015465aac86
e111ea0e3429e9beb6ea1fd09625b9dd98821a00c1ac53aec2ddb5b5fd0f3d77be0ddace3f77a8f51ce96dff7ed21d1e15a6feb5697f4e
361bc6b4740dcbe6f676f5b5f35542dab2d149b9d01b0376a62bd5ae95ec6f30a9cec88e87d4596b8deb1de176416bd7c258c11a38dca1
b258afee0415e5f917b13bfeac143f24a230c2177131e27684340486b021c210fb9070a98726419ba0eba12f8002562014d80bc27553e6
2b0a8e52976155b7aef6a0d6a84ad788252d878d0defcc3e376dfac539fc0e3fc7e25bbe609a5a61efda6d752661eae493e77af73ed8ad
07ebc7f42c880de59fd56f3d7270a7b282515a127348d741d09bfc520768627299c9526bdcaa678216c62b913bc2a6f8a5ea960395e5c2
62ed929a7838cbce142dbb2d7eb7a44ab2ba6b0816736a9f7b6ecb81babab4b7179d3e2b54b53e2decdcb5f364556b996c6bdd999bf3bd
1211a769f2a534af9209bb7b4df249f1109c1024a6156184920a9b14b336d231c96b30ebbcba745d96ae50477e1694ac4c64779faea59f
9875cf27dbbebe2f07d92608fa3c28c789ff5c8ec60e398e64d97f6b177e2989fd1f4822661d5404f1fbc4229a57c9c763281fd759a1ce
58abe4636be06368b50fffc539c9eb1e125a0cc57289a6445ba22bd197188a8d250125a692c0127389a5d8ea0bbd1d6a797027f3c071aa
e8a5036f6ed97ce0c0e6dbccca6fddfe13ff9e59f0eacd0f3fbcf9d507e7bedece3fe04dfc3b728054b2b38df553720f9f285611874a34
0df686774453ad693d7b174f4450248d5463aa53f6a12d6147407975fe88bae614d914cf7dd5b4a79e075252515dddcf9947e8d7918ff6
b71e94f5d59d720ffbb623a4fc76c331c49d05685b271b3460316099a9567742a397b5a01d6155b628aa2f501c5dbea004ce91f4a05d41
8ac5fc19e7677375c13191a37b6c7f9df838b626a897038f582de74fb61e2663e54d9794b7930594efced16c7170d33b30c028980c1322
9d5a9da0d14f888c74a6e90dce48d14e79b05cb495dacb43943ce8a13cd8d5a93744866be0f170ad49a3b5450fefaa7075b9a991d8494d
4ded488c3f2a8951f129759367fa8ece3f1af59e197d18e2285abcf31c7a87c161ec45e1ddc3d0c3384037403fc030c06870818bc5085d
f55d0ddd82126c09f66ec15d9d5d23e35df1513171a5fa5243a9b13440799bc70441d6cb063462009a3010cd188a61188e0e3142179710
3f24fe57f125f1cbe337c6fbe26fc787d09671fecf7939523d67c9eece1bfa04a66cecfa90ee70ddf8fd93cbcba7bd34a47eefdd3f4c3e
3337ef6cf6caf5b96f7adf7cf9dabfe61d11871cecda3523c33b3acad4ed95f2ed47ddee932929998f8d4df704c66c59b9f380ba1bee4b
81fd83b4936290b2b649d206e21b606127b4657a03e9987ccc6c35293138b05ed9eaa81b80a63bfe0315e5937ff2e713e5f0610b1ec0ec
ca2a4c193cd9c296b062be666cd1bbef5ed95d5626ede4ef55b4facac76fdbf57b21ab820d56b2c9418ac2496af4db6080d7f173fcafd7
b313b65a2345bfcd309ef2c008bb128ea97e8f6a4cba9f040aeca79424104439cc1f76f7578f58765049026fd5d60e3db4e8f407ecb7ec
98b0af357bd7ae935542f13ddf81bce9b771bf22fd20ca402562169d05ee79e3d46380606182a41428c820330b809c2620fc8b244b2830
49048df2fe42af9cc8403991294735e57d82725c03f518ddc5ff0621f5e76383b61dd4e383f78551c26ca15828114a85e5c226a14ad02a
13e950475e6c67611826c6d2e9331ee34597360552587fec2f266a47c008361a478b23a451b2573b1126b24ccc14d3b57990c766e12c71
863453ced22e8285ac188bc545d2b3f21a58c3cab15c2c974ae54aa8645b856df8b2f8b2b455de2fbd2ed7684f69af6adbb48395334ab2
4e799934e80c9bcaa69ee14f378b592d1978e09e4fd150dbbf49a1f232da3176f386e8755a8d2c19045100835ea7be60045a720de6cb17
14a9a9a8574dc328d9d3a94713e761e2d3ef16b51e4ce3a79785d2a614e465597f3d86453dd3a7b50e50def4b5b660bae645cac4c0dc29
e6206b50322099a0e88cefd83e1f6f995cd2daf20d6e655f08890c5bffcc4b5b6fb57ce71fc78a34e5ca1b756571a939734653fe5311a8
ef0e6d44b15ca5685708bae98067496145259bcffb2e6ef469cabf69d9cb9fe4369ecb86b35b422c267d03bfe024c54ca64eb606592d76
6143c964dee2db77cca730122458d952c1c45b5b2f7267cbac6f28b637a8e3ca95f7f9ca6a9488e92a2bead790cebc24db8928ba15a2c2
8633be8d177de7379728cc34b47ec2adfc387b9535b137f029ffb70c21f395dedfb9fbfc2a70e08f10a9553f24fceed7a6c68ef2ee472d
e34c993ae53b9ff6fe77071aa799c723004cfcee47cd8f9932ffe60b85476c50bf0180504db08eb27824d4105451fd13e13080d405ca08
3e27a824d84e9043b093a082603fc13a8215f211f8581e0745e24d28966c704ccc83f954ce179ba8b4c23121154e2ba57843ed3f8663a8
de1d0ad00d7da9ffa078020649a130889456e3bf840d74effceb07d361096c873f323b7b49d00ab9423d26e042fc4c1c29ce14df92864a
15d2bfcbb1f246f9827c4793a859a1b9a8d5688f6affac4bd19dd5fda8b7e89dfafefa89fa1cfd72bf6ea01b664077980946ca3d667845
d19c681782a954be436860b2f2ee59d41172a2faad46a93308a696bf2e90ab8f68af63a77eb1535d821036bebd2e838de5c15028804258
0a0b6016cca0d917aa5f98a6433c954990485732d5a611860bd208672114112c805cc88679d0837a47433ee1f7a2dac330972e173c7e9f
5691dacaa53297c62ca67b0e61eaff1bb3f6b93f6b06cdb498e652be36e413b6c287f2edeb7f36e330aacda671136111614c27dc6c955a
ae3a225b95c84554f2e95e4838d388ee2cc273d1f8029a3d5b7df64b3a13542a45c411edac610ef52ab316116e814a2989e64e86940746
758c11fc8ed4f66bf59be8dffebaa97ea17cb556be515bc0aaae447608a65d61089daec2d4afe2ca1cc360388c8091300ac6c238788ce4
9f004fd05c4fc224c884bd4c60ca5b2389c94ca359943f2b293935adbd1cda5e0e6b2f87b79723fce5c3e469b5c2726fdb3d8ecd36fcab
07ff9284772bf12713fec8f10ec73f7bf00713fea9126f7bf0fbe71f96bee778ab12bfabc4a666fcb619bfe1f8757ffc2a0d6f72fc3209
6f344e906e54622321364ec0eb5f2448d79bf18b04fc9ce3358e5793f08f36fcac123fe5f88915ff7d197e7c1cffc0f12342ff68195eb9
3c52bab20c2f8fc44bbf0f972e71fc7d38fe8ee3458ebfe5f8af1c1b2af1c279a77481e37927fe26093fe478768d453aebc0f783b19ee3
198eef713ccdf114c77fe17892e3bb1c4f703ccef19805eb4a3d521dc7da778e4bb51cdf393a457ae738beb35c3cfacf1ee9e8146f1b1e
f58affecc1231cdfaec4c31c0f71ace1f84f1c0fe6e05b263cf0a6473a90836f565ba5373d586dc53788e9379a713fc7d739eee3b8d78a
551cf7ec36497b9270b7095fcb411fa1f82a7117c79d3b8cb403c11d46dcfe6aa8b43d075fdd66965e0dc56d667c458f2f73dc5a19206d
e55819805b68d0964a7c69b3497aa92b6e36e18bcdb869e3716913c78d1553a48dc771e372b1e2058f5431052bbce20b1edcc071fdba5e
d27a8eeb7ae1f324e6f30f63f95a83546ec3b5b42da68eb21c2c254d957a708d0557735cb5d222ade2b8d2822b382ee758c2d1dbf6eb65
cba45f735cb60c9fcbc1e20cbb54ecc167392ee5f88c09971871b11e17715cd88c45cdb8a019e7376321c7028ef91ce746e11c8eb32d69
d2ec09388be3cc6538831a791c7339e6709cce711ac7ecfe98d58c538d3885e3531c2773cc9ca497329b71921e9f0c0e959e4cc2891c9f
a0999f48c30c3b4e60666942083e6ec3c7c604498f714c37e0a31cc73f6296c6737cc48ce3388ea52763398e196d96c604e1e8880069b4
194705e0488e232a7178250ee33854e8290d6dc6b4e3f8f058f4721cc271f020ab34d8868306064a83ac3870408034d0db16880302b03f
c7548efdfadaa47ecdd8b78f59ea6bc33e2906a98f19530cd8db89c90198f490414ae2f8900113130c5262002618b0574f9dd4cb8c3d75
d82309bb77f348dd73b05bbc55eae6c1782b768df3485d1fc6380fc67a0c526c207a0c18c3d1cd313a10a348ce282bba7230b2199d2482
33072302d0411a74700c6fc6b0340ca54628c7901cec429aeac23198060587a29da38d6310472b2158395a48564b1a9a9761600e9a3806
1883a5008e46c23606a381a3de8c3a8e5a42d372d4d850ce41911e8ae40176a45ee4b4f3304b424f6466048eac96e5acd9c0baff7ff8c1
ff35037ff717f11fb3563d330a656e6473747265616d0a656e646f626a0a33302030206f626a0a3c3c202f4c656e67746820313132202f
46696c746572202f466c6174654465636f6465203e3e0a73747265616d0a789cedcc370e02510c04d091084bce199678ff3b8210da82ea
4b3414ef35f6682c273f6a7de5763acdde4d955efa19649851c699649a59e69f7651f07d5970b3ca3a9b6cb37ba77d0e4d73cc29f56b9e
73c935b7dc0bbe01000000000000000000000000000000c0bf7b3c015ec502740a656e6473747265616d0a656e646f626a0a3333203020
6f626a0a3c3c202f4c656e67746820333432202f46696c746572202f466c6174654465636f6465203e3e0a73747265616d0a789c5d52cb
6e833010bcf3153ea68788008949258454a5170e7da8b4a72807622f115231962107febeb60713a9966034b38fc1ecc6a7eab552ddc4e2
4f33889a26d6764a1a1a87bb11c4ae74eb5494a44c76625a987f8bbed1516c8beb799ca8af543b4445c1e22f1b1c2733b3cd8b1caef414
31c6e20f23c974eac6363fa71a527dd7fa977a5213db4565c924b5b6dd5ba3df9b9e58ec8bb795b4f16e9ab7b6ec91f13d6b62a9e7093e
490c9246dd08328dba5154ecec2959d1da5346a4e4bf78b243d9b55df353970f38032f5e3e427e5ee495224a9e66210a9ab650d1314b00
292003ec01070007e4806368ea3d32092ae161d1cb7b98383803bdcce1c60f900385398739df871c94c08f2f9778504471a77cf92f0be5
6d507d528efe394752a0b0cb61971f428e2d49d3a4b11470065edca0c244dcccdc82ad0b21eec6d85df05be897c08dbf53b42eaa1eb4ab
72cf1feaccc1620a656e6473747265616d0a656e646f626a0a32382030206f626a0a3c3c202f54797065202f466f6e74202f5375627479
7065202f434944466f6e745479706532202f42617365466f6e74202f424d515144562b44656a61567553616e730a2f4349445379737465
6d496e666f203c3c202f5265676973747279202841646f626529202f4f72646572696e6720284964656e7469747929202f537570706c65
6d656e742030203e3e0a2f466f6e7444657363726970746f7220323720302052202f5720333220302052202f434944546f4749444d6170
20333020302052203e3e0a656e646f626a0a32392030206f626a0a3c3c202f54797065202f466f6e74202f53756274797065202f547970
6530202f42617365466f6e74202f424d515144562b44656a61567553616e730a2f456e636f64696e67202f4964656e746974792d48202f
44657363656e64616e74466f6e7473205b20323820302052205d202f546f556e69636f646520333320302052203e3e0a656e646f626a0a
32372030206f626a0a3c3c202f54797065202f466f6e7444657363726970746f72202f466f6e744e616d65202f424d515144562b44656a
61567553616e73202f466c6167732033320a2f466f6e7442426f78205b202d31303231202d34363320313739342031323333205d202f41
7363656e7420393239202f44657363656e74202d323336202f43617048656967687420300a2f584865696768742030202f4974616c6963
416e676c652030202f5374656d562030202f466f6e7446696c653220333120302052202f4d6178576964746820383338203e3e0a656e64
6f626a0a33322030206f626a0a5b203332205b20333138205d203430205b2033393020333930205d2034360a5b20333138203333372036
3336203633362036333620363336203633362036333620363336203633362036333620363336205d203631205b20383338205d20373920
5b20373837205d0a3937205b2036313320363335203535302036333520363135205d20313034205b2036333420323738205d2031313020
5b203633342036313220363335205d203131340a5b2034313120353231203339322036333420353932205d2038373330205b2036333720
5d205d0a656e646f626a0a332030206f626a0a3c3c202f463320313520302052202f463220323220302052202f46312032392030205220
3e3e0a656e646f626a0a342030206f626a0a3c3c202f4131203c3c202f54797065202f457874475374617465202f43412030202f636120
31203e3e0a2f4132203c3c202f54797065202f457874475374617465202f43412031202f63612031203e3e203e3e0a656e646f626a0a35
2030206f626a0a3c3c203e3e0a656e646f626a0a362030206f626a0a3c3c203e3e0a656e646f626a0a372030206f626a0a3c3c203e3e0a
656e646f626a0a322030206f626a0a3c3c202f54797065202f5061676573202f4b696473205b20313120302052205d202f436f756e7420
31203e3e0a656e646f626a0a33342030206f626a0a3c3c202f43726561746f722028526570726f64756369626c65207265736561726368
2066696775726573290a2f50726f647563657220284d6174706c6f746c696220706466206261636b656e642076332e31302e3129203e3e
0a656e646f626a0a787265660a302033350a303030303030303030302036353533352066200a3030303030303030313620303030303020
6e200a30303030303233353437203030303030206e200a30303030303233333331203030303030206e200a303030303032333338352030
30303030206e200a30303030303233343834203030303030206e200a30303030303233353035203030303030206e200a30303030303233
353236203030303030206e200a30303030303030303635203030303030206e200a30303030303030333332203030303030206e200a3030
3030303034343536203030303030206e200a30303030303030323038203030303030206e200a3030303030303434333520303030303020
6e200a30303030303130363938203030303030206e200a30303030303130333238203030303030206e200a303030303031303534362030
30303030206e200a30303030303039383134203030303030206e200a30303030303034343736203030303030206e200a30303030303130
393237203030303030206e200a30303030303039393434203030303030206e200a30303030303134353838203030303030206e200a3030
3030303134323132203030303030206e200a30303030303134343333203030303030206e200a3030303030313338303820303030303020
6e200a30303030303131303934203030303030206e200a30303030303134383230203030303030206e200a303030303031333839392030
30303030206e200a30303030303232383830203030303030206e200a30303030303232353230203030303030206e200a30303030303232
373333203030303030206e200a30303030303231393230203030303030206e200a30303030303134383637203030303030206e200a3030
3030303233313034203030303030206e200a30303030303232313035203030303030206e200a3030303030323336303720303030303020
6e200a747261696c65720a3c3c202f53697a65203335202f526f6f74203120302052202f496e666f20333420302052203e3e0a73746172
74787265660a32333731330a2525454f460a
END FIGURE

BEGIN FIGURE handoff
255044462d312e340a25acdc20abba0a312030206f626a0a3c3c202f54797065202f436174616c6f67202f506167657320322030205220
3e3e0a656e646f626a0a382030206f626a0a3c3c202f466f6e74203320302052202f584f626a656374203720302052202f457874475374
617465203420302052202f5061747465726e2035203020520a2f53686164696e67203620302052202f50726f63536574205b202f504446
202f54657874202f496d61676542202f496d61676543202f496d61676549205d203e3e0a656e646f626a0a31312030206f626a0a3c3c20
2f54797065202f50616765202f506172656e74203220302052202f5265736f7572636573203820302052202f4d65646961426f78205b20
3020302036313220323730205d0a2f436f6e74656e7473203920302052202f416e6e6f747320313020302052203e3e0a656e646f626a0a
392030206f626a0a3c3c202f4c656e67746820313220302052202f46696c746572202f466c6174654465636f6465203e3e0a7374726561
6d0a789ced5b4b6f1cc711def3fc8a398a803deaf7e3e08395d80a8c0086650239383e0814254ba0655b922d0741fe7bbeee9de9fe6a77
67b9146d910122814b4e4d4f4d3dbfaaea997df8d7cbdf5e5e5c3e79fc68fccbb7c3c37e74f176d0e32bfcbc18d5f80a3fef473d3ec6cf
8b41e1e8c7216883df57f5b789e52f35fffe61189e0f0f3fc7d2b758f178302a4d2a8dce4e06d7f9ac263f1f5d2d47c6f929e2705eba1c
564ebf8cccc026337950cae79bcbf11fe3ebf1e1e7a6de6aca49d96cb40fb91c681b8d775999414d293bad6370518f6f5e408957f879bf
7ec528ae78f2783871e59b17c360559e7c0e3ea5516b3d59139c87da444e7e72c9db64a121784c2906a3f21ab93359251fbc65b5dea3e1
9bf154fb99f1ab6a1935055f8d93b2d239a8985cd1bfe89e8289bee8dfd47ff278944163749ab20f31c6c5e14459fc3a7c5be56a8e3361
fe177160dbbf72d0eebb384e81c3fbe1c46b20deeaca6187fb71d19daaf10a8b2272865ff0a9c64f1524323ac0b8d9a71c7d49853c251f
8c1f2f7e1c1e9d0f0fbfd4631ecf9fd7a4397f367c373ed8a4b3f1fbf1fcabe18bf3e19b8fee1f8a955949a2dc6fff1c157dd53f568529
9aec94d23504aff74fbe4bff40cec9d86073539228f7db3f47455ff74f76935730a8d5cad953fca337ea4e3dd4417851b353eeb9878e89
beea2154a5493b9dacf5b614e7533ca4efd2433ee9c9eb802b5bd3d129f7db4347455ff5908fd029e9e06cf6f1c41c32ec2105df18eb01
93a5ab9b7495c7795d24b773b353945f5b37f0baaaf049fc040e583da1f74135cd4937d955b9baea02c19b220f365f9f8de7af8688db00
5c9c2d1151563cd8bca927ec1481f7d9dbc2693ef3ac9e81a1b44154857ee2653d816054caa564526d6ceb99d7cb25ca69b515633ef374
3ee3b5d736448d0a349f79370be04d4a2af96cdb99cbe51a1bb5ca24f3380b90c2acce7cc9c32fcda2efeff3b5e8ed34eeb65c4b165963
f260f3d3725f15ac8b016a5cafdf3ab775fd7e58e57663cd1f6c9ecfb741e8381f95ed4e5977d78a918edc7f35263e44fff598580fca17
6b02ac8abcbd7fcdda82966a81a353216cbf633f329e012455b9c989f0d007abbe7e585bff6130614a6940f140671e00784ac1ffba011d
f003e2ef60ddd79bd79bcbcdb8b9c0e7ebcd3b7cbec1d1bbcdaf9b9f3757c5d08c83a796811321fd18f8eff23ca4266aa4412027444e3c
0c89664c5b8d1f6cbe58e23faaa0bc0b08c04ff5b44df906a01e13cd767958a232795fe34f2f31e6664601ad41c47093dcb8c7a7ddf5b3
ba384d3126154be2e585cdbf67364ec5321f25e7963371f5cc27f58c9bac8f798b22cb99b47acd7faee1a6cd146ccce834917b1d5a3fd4
6efa1abbdd8945afb3ce0d1bb192c966cadb5c8e0a1316f2b4c4b6ab19aad0d1e3efa3e3bf81d9add3c5e4b201ebf4fd2eacdf6bdcbd97
c497a3325dc347b61d80f2e0fc327e188d7485f5bc8b3e500395706587958a2b579ba70d4dde025d80f2a6dc4c632c05bf5450bc42ce08
a0798a256f71f41647979bdfea65ffdabcdffc50fedab970acf0fffd708fbab3198f80b8a80e1efd38c66a7b1c76cfcf40054acf4a6990
de6f7e82fa579b9730c33bfc7c528df13b8c7381e311c6780a733dc3aae7a501b8bfb8ac5d9a0c721781a360c3955eb525f4df67508bd6
fb5cbb99dd326fe00cdc937a96e7730792b371b654bab0db81586b62288d46deed5a22f24c3993c309c091af39e3a67c63c03d24c1ba6c
4b77e43212b57641bbad9e0dda47638096bb7d13069e128dd4512d5da84dd128d138feb176d3ab67d49f07c5a64cc22545b7e11cb687e0
3ff979131a61720c944b0263945739caed19a23750fe732b05dfdca34047645c929582e807e6f5668a71dd14a2669c6abc93790b602806
82bb5d0404a0b8473d459394b12e4393b53ab22d239725d931f846cc4d513b9302c8ef64d9d899d68f96c9a3e7ebd38c93cba30b19f9af
6c02eaa1cfd601437ad00ef51b4638aad61ba0fc0b20fa713dee4b5d53087fe713402164b75ed76e52d87eac4745e967387ad9468fdf60
8cabb381cbfbfd2a6e294d0872e5ad0dcefe4fd7b6f5a6f83a8cfe7f6dbbe9997b50db90c426a604f8148f8688feb16a5bdf55cfe06a91
005196b64edeaf6c6149d092ac86915a77ab88cad6af18d6afa8a870226f31213954a65aa8eb849431c4469533e658a8bf8a0e618b0ecb
0ea30bc012858f3d742883440e8e22a9654606c46a13b2df8bf2a4632c7bff7637336e729f253731f12ba34aaaede50cca1cee956ccbcd
9f97fbc4f2acdeeb9eb5cbee9f01dd66806193faed8c900ed48c192f85ebcfac4bbd7ecdbad4cb0e30b021ebe0956d18b0ee9f756edf2d
5e0880757437b9a1509b2512cc1363b6123a21b4c1f258a06b0f50762f593df1cfb38e34f7a774e658fa3dccc6a124890f13d0b9a05272
86dba4dd5ee2cbb3319ac9d47da012fae821aee68ee167f414dbae61dc7c3677101738fb2bce6e57f0ffbf6175b9f627f422fbd75e9e0d
3a4e7e5b4cea46391f1257de76b8f9633bc05bc523dd1ed455ac246b2e8863f2e97b477d13faf8e6d0647780fa437a5f2145a9218b1488
dc29e40c3464ba291b37ba607aa5224bbc8f959a2d34c65a749028f1c1220c86b2bbe4828ede543a2a522c0f54418d8bde30614e003e49
6dd6d8f268740be7383d0abe768ab8be6c1eb3146eb2f8b095411719e4e400086545574fd7e7bd553d7649a35e5423e1e2fac49b5723a8
ca130e27398709d522015e851820a3374f8002961954bdab1dd3d8124cef7663beddc62c047b844426ff917ae4eb839171f1a1ef596d93
a50c276e79b92c96f02c0fb54b314e471f6cdf6186b0149c211e2a1b20607991a8d30b15909fcae66ab7b04fa5274601dbf10706e809c5
1881c8de2b54bbe8dd3c2da81417823ec790e0dba24d48d1439345ee71cceaf598172ea10cf1096017adf372318c1b34c24d308e3096ad
ef51b210a02a8d96cb0b89639ae2ae6e4c633b30bd5b8df9760b930cc21d5d60765ed38dfd7c282a6e991d2825dbb73c901008cc251495
13674edff16a8972f2eed19c327ffcde9194514f6191919098e884da858ace4f1bcb10af83991c1ab5d29473a12974f81d96e39202ffa2
8df7a50e0b6ab3952834e81927dc1a357c149c3da61c03e0927204a81d33dc226a4da167635534a3d0300269b71ab23d1a55d41ab1ba01
b7e0dc605ecad18b024bdd0b086bc854b607d3bbf59873b734cbc17e61a9bb1759c3eef183f171eb844a3c2beaba5986a8f4db096ed930
3bba5b7b20a1ca6be2fba97224019e6c5f25977c29090a3d47009650de29fc99b53186cde7ca437298394769ec42c7f462a34802a7009b
cb1d9b1b05959c2ee82d4404e7164e420e0a3e21750b55a1610b6b610f4a029b7399f89211498065e5bd2f1f451214aaf118ceb294c396
19059e308ea52e54b7a7a1a0923d04bd594f706e961672905f84d4cd8b4243f2f8a1f8b87512f8f6367f1950d49212654039b2d7b4f6c8
e2e3b65f2c85cbc58d5b2938a03a9d83cf973df188515e842aa8211be3f34eca788d25d15a9932a062e8f749a64cb7864c198ffe4abb68
64ca7837956d9d245306d4e08df27127653c8232466d65caf8b207336b485e695499329dcec1d73973a8b21c1cd85d6a4e82aea1a0923d
04bd594f706e961672905f84d4cd8bac217bfc507cdc3e65dcf22a6ee980e8fb2f86bf29b33eb0dc61c6b0149c31d4d1139dba7fb2358d
0ac2333458901f6904219f0b2a4588a0b778129c5bec0939285285d42dae59c39e03c22b9431bdaf178bdb08c08cdbbc20a5e8c305cbdc
0711d68fa96c0da677db31e76e671243388564261736f5d8db8762e3f65d16bfb72e77c0c2d21f9dfe30e3be8e2daabcc336cbc84d7da7
f300002acc9fcb8c49e342a900984c53dc195b404fc629b913a6515be06fab93a4365bc9b14579649076752f8c387b20af8d726a51a124
4a4c7e676a29df1e4a36052da6161581395b05d91c8d2aa7165eddfb7fe6dca7052107cd165d681a43483fa6b23598de6dc78cbb9d590c
f60a0bdd7dc80a767f1f8c8e5ba753df211b1f0fab5f683b94331f3c9908be14e7b6ec8fc0025ec4b9d57a525e830ddba8503df8a6282d
8a86779a5f6e25fb17aa5deed87c25a8e459416f7120382f2123c4a0f81242b768140ab6c815e6a038b7f067f2aee227af4e65cf2b0611
e7984eca4bdb21e51d39ca8519fd3dc779fddedebe7e4c656b30bddb8e19773bb318ec1516bafb9015ecfe3e181db78ef348df59a64904
25ac3f853af569f3479e4ae88bcc197acc42b85c87b682079dec20a64598a18fc055cac45ce66110930f1a462b5f41f056b93a4de4fab8
b7d46fb0aaaf43607cc965f4dc6a0c47656581594c6c66d83258c870060656ed47e66aa75ce563011c7a02935c9dde495884a18d75d4ef
6a816817b5ba1f1af1a25a4697776978298ce18cb6087f629a1069989db5bc3fa8c807836ea78b0a9adad38988ac3f91bbad8869b76abf
3f3b804425672d3ab1570f44c0ad3342712335effeda5cbf773fa7873ff028f1e09b16773876704a504fdec9d4be772b53afcf1ea1c1a0
3b8f6688ee6626524830b9850f735de28c05a09064615bf4925a3dced90f9412bd13e7a5ad6727a6adbf17f7efb34017b54f0da4131159
7f22775b11d36ed57e7fe1802e2a396bd189bd7a2002969418fe0b9f6e8bde0a656e6473747265616d0a656e646f626a0a31322030206f
626a0a333338340a656e646f626a0a31302030206f626a0a5b205d0a656e646f626a0a31372030206f626a0a3c3c202f4c656e67746831
2039393630202f4c656e6774682036363736202f46696c746572202f466c6174654465636f6465203e3e0a73747265616d0a789cdd7979
7c5445b6ffa97b6eddbef7f6dee9ce9e5e084d4020895908619116d9832c1231089104921096844840641089e3a8332e8ff8902da21050
044589094f115011191584014650111975c4d11903f29b1f83da86ca3bf72620f21ce7bd3fdee7f7f9fc6ea59653eba9ef3975ea541a18
00b82991c1357cc8d061702de401b064aa4d1e3e6eec8450d7b4e344f7a5f8d8f009370f5eb174f30a00e920b59f1f3b21236b76eaed7f
03c03d444f9c5e555ac3b6699fd264f1447f37fd8ef9c1eb5eb95605e02e1a53585133a36ab234fb1c80a58cdab7ce28adad010b0550a9
1d6c33e62caa18bfe9dbe544d3fa31f32bcb4bcbb4eabd530192eaa9bd4f255538b659a89cf40ed15d2babe6dff9e4397d14d1ad44df35
67eef4d2fc9481c46f7225d191aad23b6b9432a537d19b890e56975695f7b25d5f43348d67a2666eedfc339ffed909e037e6bbb3665e79
cd5da51bbe04084488c767c0c0c6061d9f44012181ea12c14f65a34d87741800d29061a30bc131a7747e35c4138ef4b5b7035c2e193dd9
ecf279d5a01a2533ca348391ab94ffd9ec799e30b700f1d2fe1efcc2273ea624437c27ce8973267dae3d60441a17f8a571262ff5bf38f3
2faedb3ee9175b97fc4223a37d715028b5420cc4114612514660a6f425aab712b63114646a8fa3def1141413271fa506f232d502d5199f
c54412400343065630a4640747a74c9e84ade03565f2abd279a5d3a0be745e5535d44f9b573a13eaa79756d7525a593e8fd245f3e640fd
8cf2b9549e31af7c36d4579656539fcaf2695433bbb4ba14eae794ce0d1a29c9f63755a5f32ba1be7ab6513377466915d4cf5b504d3de7
5754cfa0b4d298ff9fc8dfc461cecc19a53fd101193a7480411f33e7b40f2f244317e84e3a86b43307a498b99d6a25a29228b593061a98
310a94ca2ba4dd06363c9bafa1e9fc1d397e00159287d4d6aa222ab224c99deb5efec6550c2d8320853b15aff0b2064b15fbfc277db033
269ba8d2f9228a99b40c25947bc165ca2508bd602c14422954c04ca8813bcdddff583b1d2a610ecc6b6f6fffbcfdc3f697da9bdb37b53f
ddfe54fb86f6c69f7274f9eb38737775529a395b4734f8e9d5198d1d8ded8c866e14764643474a3b23d9165a1f8833309137ecc2ccce15
e650245b40a802cca37827c51856c1b25905ec86bf12140361034431609efd0aaa35f2adac105aa97d1af5bc47be8f15525e253f0512b5
2f950fd11c12cb8669703b95c2f2536c37ec842f68f43dec113e82df6af4367765cc7581bfc9cef17c291f8ae42a79a0dc2cdf2337538f
0572857c0f34519a2f1d95d7ca8be5c3f26228323863a38d68f0010d6c144b8506a9810d61096c887408de209e2bd820d6c0faf377f9bb
701c8eb371d4732b2c9474f636fb3bcb6045ac99465d800b2c4054ae94cbceb22f89e3d570148bb80e0db08c7988da0d8788ef2fe0ef50
2bd3acb08c1f977af2e3f0267c061f503dc02c66d8b014eccd8f5338079b611621f31993f871c56b09c915d277d0caee953649dfb15426
51f0b000a1791b1e924be4b7e5df512ba1c324ccc6000ea6748ad1831f670dc4c5674a055b44fd8cb098d66995de9476d01e5f8353b42f
5a5d9a222d961ae0147b81ed248e01ee632fc8259669721234280d72119c35b081a3d221c2639c89c743f090722d5c90153887a35989bc
d9400cc2fc0d3a4a21cb28c5032bd928cbbdb413c03c584c1a0eb09f017fa323502f55498195721a3e49bc4bd2924bb8b1457048cac769
b0d60ccbd90e580e3ba0d63824d8ed658bc2659418f40aba9aa4f0c8b2a6c8f8a2e03b9342bd7b5d45065d9660138c6bb22f0aee686f1f
572427f1494d3cb909c36a931c4efdec9f357ed6bb57c1b8a2e00ed67de890ce6987960ca1ca0945543428aaa6faa143cc3663d5261ea6
bf91254dc1e995c1075d0fa6f67bd055deafb7a999742bd119364e7b953825e5b37be97c25446cb805362b2893d86314575bebb183d702
6599312137a6c61c6595df1e2f63f78ac3e221b6c098a5a2fd73a582f4c24a163d2b12a3bce081176c6b3dabe3b5dece3cecedeb194fa3
db5a6992d3e75b5d6733235648888b712564240c4ae0ac9875917c5e4f7696273747c220b85d103253ac68dabbb7e9c5bd7b5f64956cb5
a814abc51a3183ad914f88b6d6af451b93bf6e65328b136562855829cad85a368bcd666b693764c5f82cd2610b64471c8ab4071e91595f
8c85be5c751d2b68d20b0b9a1c85938b5ea1ed46fa4e3279cb38df9a95c99a231a2ba67d9a41e725c22d960a173f2e4efd30466ea6bd16
b57f2eff9ef45487303c11b94653202160853fc61d56d6398eb8830703ef26af4b3de05e6d83d4388cb76b76ebc000dabdfdbb1104fb68
892c777e7e3ead76fa7c1b61f197b317cee67bf2dd9efcccc8e88cb481c181a1416937066f0c15078b43d5c1ead0ddc1bb4335690f071f
0e3d117c22f47cf0f9d0abc15743be2c7f66e0067f2430c13f2e30dd5f12b8cf5f1758eeaf0f6cf0af0f34fb9b02ae620356c5e78dcdce
ea339085dda15c074bedd22d37a76b283b37a75b6a17c5e2be8e6507656943cded93c7973fc8668a5523b6dff3c209e6645ddebbffdf6a
dfbab9f6abf96438ececbbd1a386dcf868558f072edeb3a9a2f8dd0dbfdf917cf3d8f474e64e4ef986b03610c9254412202f920807d931
d971d07e4c6f74cb8d710444a265b01dbcfd135d6dfbb22ea980387fd6f58fb399dba7262d4d928853773a31d7c96d9e9b38edd082ac58
3977e4f2b12b9f796665e1aa48e1f3b788a3e259369165143d270f141f67656e7be2896d59d78a938100cb633e0a7901431fa7b47fce1d
8a716b2543df48525c231c71b81bf91175b5e300db885e19ec522471b0b57f0a09c59009c9a3f5f4f9d3adaed3c4d4587f89df642ad489
9ac14b76562c66f5e9408ef8c44d8cd47154f382c3e202d30fcf6fd950bb6851edbc458b70b754f47deb86e953d84886144616b7eddfbc
7efd662376680fd765c3de24c10d9110f83e64da41f5185f67631fc4aff31cb0ad4e4ef249aacf0e4324bbb37fb2c9df3e528f0ee04e9b
c89dcdcfdc9e913228c560d217f2b39f821604de215faeb7edb2af7df8f6334b96920e1f11db5801ebc2543650d42f2ca9fcb54bcaaeb8
fbee1b8688d6cc6b592e8b238bdc4fec5d5eb16441b5e11ba491657c84308c81f19164176736b55161eb60a34369d1a518f2e6346eb70f
755abdae8f0b9aec97cf93d5384f035adb06ecdbe7e9d0f3ac363af2594464b208afb3d4a912714dfca5b08e5396ea2685ec938d8fbc35
7a00cb127f100dcdcd07df57bc5fe70d19d30e6debb184c19897b7193c6da553ed219ee22037126fa7cb4a3fe939e23be0dae860921d86
bbed76a7cb303559e78df33ca095246a9cb6cced250975091d22cd75e7744bebd03673753a0359b1dc2356db5dbe11e93575a281558c7f
a97aef7ee9d98b13e7b235cbab1353d39e5f73f184e2bdb8795af159f33e06c9e002c9baf420eb720a5ae48d1267c4904b75b50d6835d5
e97c6b66441ba79568355a9d26ff6853b6b20a630dd12097fcb09efcb18f3bf6d539a3ff65d828313bb864976993cc497809afe175bc73
127302c5fb7dab81c8b3e4f86c277d0ac3cc488f249f47932d10542cf1ce93c123a978c0df924ca2f27a6caa5d19e1b57b460492ec292e
3245054dce4292568a21adb6018671ce37952c63c080d3e7079c6e252b4521ce909933909699362ead26ad2ead3e6d5b9aa5989998f94c
f02e999290bbe364183bf499edf2ce21bb6ade38205633366c4cc55c49ac8edc34a386c8cac1cfcd98df8c9b2aabce7e7e71a234c29e9c
b870f6e675173f9246ec9cbde5898b27e492a7a696d45c9637edee6a79b7fcbcbc4fff6b79fbfe85bc1f7ff492bc890b43dc9d67368db8
306eb7de74bb357aa0d176c0b8dd063b47e3605fffab6eb797f26206250c8fc12b2c30c174e5c526bdb170c992850beeba6b010bb1a162
97f844fc49bcc286e3e2e71a1b9f332203f18e68a5f00eebcbbc14fa1a9c6c1513f96dc48961d7fa47527eb46b071cabd99fb025856c5a
c4b46e2ec3b26599b8b84e9fbedab4852f43617043d742cc158c4a1e039351db171c66bab87078c1f6a76a172fae25f3b6e1e276452750
c4cbe22285978b316fcbba755b4ccb66ca0a5b893737dd02b16e87a65800f7581f731cd05a2cba62a787bfc750bb1843edd4f63d7d27ed
33cf09b1d6b6cf6d5e80ec6a31c5616b4671c6832b0c7e86b62cf15cd31d33627d2f3e7db14d2ed9515d8edcd0915974fb4ca375d3e0e3
48c46e931cd63efe809f2b1655e3b2de2710f08775ab3f20fbcc9bc97bd0772c7e9d5b5e17a69ba9bb5fb706922c3021e92687d732aecb
8ddd8d3bfa58eb69e2c6bca43beeaa7f18779587ce02f9581697e30c1d8bce6c529766d5ca5871e4364dd374cd6ab559ed9a93a726da12
ed898e78672f355d4bd7d3ade9b6747b8f60beda5febaff7b7f6b3f5b31768a3f451d651b611f685b685f69dea4e6da7bed3bad3b6d31e
76280e8b4375680edd6ecdb30fea31b587560c575c910126c7fabc72a84bb734b769fb496619849d6144e9d8c5d5be3fb5627a41e92016
f39af84e44e79e5932fbb3f933678dac1af4cd9ef36dd33fa29bf35c6666766ecf74ab96bafeb996eda9a9cc9593d32f3f33c3aefa373c
ddbcd56fe09a44f27c8a3f49f6bf2a92e8e0aa131bddac456d045db56a92464ae6f2380abd1da6a4a0c95538b9a0c96d241ef31e7098f7
c0be1fef817dade7071866c53029f18360508cf96f16a7cbe91e278dc371be12a9043b366a98496f6c7f2a90bf625c0f6e299bdd2e965d
3765873874ecc5e666fea4d8db0e223c26af1d5e3cc64e92b77e1d71bc9ecea92297901fd23f92e4249fd676d27b246eb58bb578147264
9d76f770b21aaec48e3bbfc3dc9dcefad16624d519be480c699f97fcd1d8b80e702f7b01ddd6b34ac9e18e1d4e56c330c1e3ffa36aefbb
6cbbb4b566b23893fec0c2a4d46e5bd7483d7e58bfc1b41b0cc815e109c48f02ef45d23045e6b294424f262343490185a5d00b362c217c
c2154e8f062e83c5756cfb3a1a2815ef8b33d86bcdead43a7ec6e2523b233f33a90b8bdc65bc2b5d148252a61491264a15529d542fad97
9a3ac31e0a47287c6a8673141238bd58eda8d3598861c998287783aeac27f690fb400ecbc77c39531d06c3d8481c2997a8772a0fb0dfe2
03fcb7ca4a7a2badc135f20adea06cc697d82bd895e41413d2c8788502f40c8d63b16c86182e16ca256d51547e586fec7e04dd4c29b47b
1bfb7564949262512d9c2050558b9e63d5090b2d8790b092c32fe9399a8e2819a8d8a83ba1620beb5c428590b1aa36abaea91daf2aab85
94aed37bb71aaa6633550d0c553b68e2451a9615f7b390fd1c847ff1e9419ddec53c531fae47284cd5e7e84b75133e4e412708f951fd88
4e00f26ff4b858c94bbea387bb14174de3d1bc7aacb51beb4a8fe51eca356a4fadab356ccb70e4422ecb91faf37e4a3f4b9edad736c831
5c1a89c3e4117cb83ac2360927ca13d5c95aa175926daaa3422a91a7f112a5c452a296e925d605528d5ccb6b941a4b8d3a5fafb1920894
dfa8bfd3eeb73e645be6784c7ddc56efd8246dc14df216feb4ba45db647dd64602515e525fb1bdc5f6e10179bf725c3a811fc81ff12fd4
bf6a7fb1fecd36d914561233fe58c8ca4223d894ddafb29e14a7881362d1abbb05bdbedbda50be284b6d3fac47a94d746aaf4ef2b3b245
91617a8a15ac5a8aaa8191e964dc1535c56251e414437553485646261bef48d27106618d8427cbec13dd629618b730abed92ec9c9d66a2
a0c9fba3cfb8efd83ef72519fef745f8119910cd85012d88195a360ed06ec0d1da4d38459b8673b45afc95760fd6e123da326d15d4633d
ae52d66aebb4f5b84d6b4223bcae19610feec1c3da61ed081ec14fb44fb44ff153fc46fb46fb16cee1b74abb96458a29db54d463a46439
598dd1037612badc53edaa93a0e5fe6a8e9e691f2215c805ea107daefd7ea89396c9bf5396a975fa2a582135c86b940675a5fe8cd224bd
2ebfae1a7a75583eaceed13f8023d227f287ca27ea11fd4bf854fa46fe4af946fd54ff1e2e28fd486c9a2934cd941a7b5bcc618bbff88a
2da67cb3b8afed7b719f34504a15dbd9e88b9f5c7c834d134f76faed8ae1b7a7c1b0480c2469f658478f78efd264c51e1a1a482dec6e38
a8e4300d3878f060db007ad747dcf1c9c1c6904b5a674fdbe8b3845a12fde44991eb407270d171ca349c758be1ade7e6e4f5c90e5da2dc
9782e9bfe77538f07e3f6ba8a8789f79cca2e9cc8b0ac3a1a74af2e89392878c19f3f22d4566a1c3b7bfe4dfdf5244f6ac27e9dc60d30f
4e80dbc8d34ad0643bd8ad6ca953b3c70cf5b90b132f714facbb8d6b251213ab491e47a3744ae7eb64d8e8d55b5c36837feae73ad3f1ec
88b8cf25ed496a22f35e9754935492144952c93523d6f1ca67c88f0f9113c65378d45b05fd24215689d962555393f11e111fb3aec69be4
872bb8ded689b8c9b50f2211370755473b5b4a1eb76b285d90b19d1cb775f2eb8a5165a7b55193e929e5d65aec92c9ece51792c9c9d56f
2369b1f436717296de4709e24bf37d64f8ffbcc7d52f24c9e0458d102fd7403a4c8c2438127ada335db1bd96825fb3277a9706147bd7a1
5dba15665c81e379e3b56628822f3190da987eaa2ba9424ffb35a40b5d5be25342c68bc0d006531dfe953ea4fe54257ad01dfe53b52848
49e94015afd00c798a81f8d5fa7109e71f35c47c2b193b34d0ee013746e2e2ddf69e2ec75248a1ddc52ef52bf6d4a1a170e13557eeaeb5
cdd89b27d1dfa531d5434ade63639c25b5c5951cbcb4b19f5773530eff52d13be5429429989fd375434e3fabeee67fc4a5496b2e388b5e
9aea1cf00f0874fcd471f46ec7c94bf9b7f3da5e701ed2cacddf0c2e7d34ce5225e89a7485be9df77d96f3d07ff9df7a403e64fe1f1c24
7a1f4afbc993d3a14aa13abe00745e0145f26b50c477c2145e054548be336f86add26b1d519d08cf1a345f6af6d98a7f87adf29b300bdf
8324aa5fcff32060298311bc07041463ec59e849314d6da578d690cfe5cf0f33e03196422e428ba44a37492ba4d7a493380a1fc13fca41
790a577929dfa27427dfe0b43a587d42fd5c2bd4769abb09e018b20677828d74da051163f7fc8c14a25c36ffd77fabf13f4b99ae239669
fe56609419c412d551964065c33acb084136a6b32c5f51e610cfaa3bcb0af8d9fd7003cc851a5804f36026715e09f32108dd613ae95a10
b22093423695a6518f200ca63ef3a196e23c288752a8825e543b12aaa97f3a95ae8739148270d3e5b96a4daa9cf2721a7307a565d453ff
6facdae7f2aac6af2c77d05ab3684c35f536f828a531ffb315875069168d9b080ba887f1ab4da9395bb939a2d4dc519066a9a6b486fa4c
a3796752bf208d9f4bab979a6d57cf33c19ca596389a4bfdcbfe496bf072fb4493ab5a9a6baeb95216f1960db93f19776954efcba3c0fc
3d95bef6bb490b7eee0b9867c4f825d14973e6d17cc36038f9a323a10046c3581807e30995097033ad700b14c1249a673293600fbcc190
c99605d533870dceccecccb33bf39ceb7f23d5b16e200059985eb9c8bab210f403ceba4294a85488a5b44b675d17b39f51461634db03b0
8b523fed00c9c5355a93e996439664fcfec712cd9a04338d37d338338d35531ff3d2838a339f491965643166d963a64ee68025d4ee3429
a38cccce6cf030d5d9cd3a3bed4e663666a5dd72b30529ada33a2bd3a11bd5192d486984ea8c1a649a395235530b9d4423354628cdabd2
f9f5314c31f7c5cd54367ba1b923c9ac61660a91f625d87e1d0a816d3ff4e26d027fe8855181df7f379c7fbf04bf1b8edf46f182c07f08
3c2ff0ffeec2bf0bfc3f02cf09fcc68f67059e69d5f91981ad3ab646e4afffa6f3afb3f06f3afe358a5f3d1acbbf12f86514ff12c52f88
f842e069819f0bfcb3c0cf047e2af013817f8ae2a98fe3f9a932fc381e4faef7f39365f8d18930ff288a27c2f8e1d130ff308a1fbcefe5
1fc4e2fbc75dfc7d2f1e77e1b1f7acfc5810dfb3e21fa9c71fa37894e63f1ac6238fd9f891543cfc072f3fdc0dff70c8c3ffe0c5431e3c
48cd0753f05d2f1ed8bf8b1f10b8ff9d62be7f17eeaf93df89b4bf1de6ef14e33b11f9ed30be25f0f765b8afdec5f7097c3319f70a7c43
e09ed7fbf13d517cfdf924fe7a3f7cedd544fe5a16bebadbcd5f4dc4ddbb9c7cb71b77edb4f15d4edc69c35768b15704ee10f8b20f5ff2
e07f08dc2eb04560731cbe98804db1b88de6d916c517287b218acf53ffe793702b655b97e073029fed865b046e16f88cc04d029fd6f129
811b3738f846811b1cb821223712508d515c4f43d6fb711d65eba2f8246dfec9647c42e0dac777f1b5021f6f28e68fefc2c7ebe4866561
de508c0d11798dc0d5a41dab05ae4ac7953470a53fd28e2b68e88a203e66c3e554b5bc00ff9db27f17f828e1f0682cd6bb705918ff4de0
23021f16f890c00705fe4ee06f1f08f3df0a7c208cf70bbc4fe06fb2f0de95f86b81f708ac4bc0a53ade2d7089c0bb042e8ee2afa2b848
e0c23b36f18502efd8840be627f105519c9f84b5519cb7046f175833b7179fdb0baba35815c539519c2d7096c099022ba7db786516ce10
589185e5653a2f1758a66359449e3e4de7d36d384dc7d2121f2f5d49d7b79b97f870aa8eb7092c163885e8290227df9ac4270bbc95a85b
937092c0a228de227022d191f689026f1658e8c7095ebc697c02bf298ae3a9617c028e1b9bc0c74571ec18371f9b8063dc78a31f471778
f9681f168c72f3022f8e1ae9e0a3dc38d28123a2387c98970ff7e1302f0e8de2901b1c7c88136f70e0e0ebc37c7014afa739af0f636490
9347040ebacec10739f13a070e1c60e7036371801dfb97613f81f95eec2b302f06fbe426f23e61cccdf1f2dc44ccdd23e7e8769ee3c59c
3a393bcbc6b3bd981d91b36c786de6267eadc04c9a3f731366d8303d067bf7eac77b47b1972fcc7bf5c39e65784d19f610d8dd8769716e
9ee6c76e410cfbb16b2a01d0b3ab1f53ddd805ecbc4b14434e0c45e4a017033afafd98929cc053c298ec8ce1c90998bc836cc6a372921d
13130a78e2124ca045130a305e609c1b6369b5d828faa8ce17466f19c6b8d123d04db45ba0ab0c9d0e1777c6a0738fec70a1a34eb6538b
3d8ab62cb4d2d6acb168ad93753bea115913a80ab40854b8ce15815c471e91e52862194a344a1264bdec9cb911ecc876b0b2fb1e613dff
fff8e0ff3503ff8b5f0afc2742a789d30a656e6473747265616d0a656e646f626a0a31362030206f626a0a3c3c202f4c656e6774682036
36202f46696c746572202f466c6174654465636f6465203e3e0a73747265616d0a789c6360a01030e31067a1d460306085b3d870aa6107
620e064e062e066e208b8781172ccac7c0cf20c020c82004640b33883088328831883348304802001d8601430a656e6473747265616d0a
656e646f626a0a31392030206f626a0a3c3c202f4c656e67746820333136202f46696c746572202f466c6174654465636f6465203e3e0a
73747265616d0a789c5d52cb6ac33010bceb2bf6d81e821337b1133086925e7ce883ba3d991c1c691504b52c64e5e0bfafa4955da8c01e
6676462cbbcacecd4ba39583ecc38ebc4507526961711aef96235cf1a634dbe520147789c53f1f7ac3321f6ee7c9e1d06839b2aa82ecd3
17276767787816e3151f190064ef56a055fa060fdfe796a4f66ecc0f0ea81d6c595d8340e9af7bedcd5b3f206431bc6984af2b376f7cec
cff1351b843cf21db5c4478193e939da5edf90555b7f6aa8a43f35432dfed54f94bacad59e073b41477889322799279993bc97510ed011
46f9b08f72808e30cac52eca013a42929f482e929c28dd511c9622798f444fc9bb52aa5297656a3ed1421020815c3c3152e6448f145928
f550520f25f5501604e592b884892eb30bc30d2f61dd1cbf5beb97169f4bdc56d893d2b8be28339a900adf2f216eaf510a656e64737472
65616d0a656e646f626a0a31342030206f626a0a3c3c202f54797065202f466f6e74202f53756274797065202f434944466f6e74547970
6532202f42617365466f6e74202f45564943414f2b44656a61567553616e732d426f6c640a2f43494453797374656d496e666f203c3c20
2f5265676973747279202841646f626529202f4f72646572696e6720284964656e7469747929202f537570706c656d656e742030203e3e
0a2f466f6e7444657363726970746f7220313320302052202f5720313820302052202f434944546f4749444d617020313620302052203e
3e0a656e646f626a0a31352030206f626a0a3c3c202f54797065202f466f6e74202f53756274797065202f5479706530202f4261736546
6f6e74202f45564943414f2b44656a61567553616e732d426f6c640a2f456e636f64696e67202f4964656e746974792d48202f44657363
656e64616e74466f6e7473205b20313420302052205d202f546f556e69636f646520313920302052203e3e0a656e646f626a0a31332030
206f626a0a3c3c202f54797065202f466f6e7444657363726970746f72202f466f6e744e616d65202f45564943414f2b44656a61567553
616e732d426f6c64202f466c6167732033320a2f466f6e7442426f78205b202d31303730202d34313620313937362031313735205d202f
417363656e7420393239202f44657363656e74202d323336202f43617048656967687420300a2f584865696768742030202f4974616c69
63416e676c652030202f5374656d562030202f466f6e7446696c653220313720302052202f4d617857696474682031303432203e3e0a65
6e646f626a0a31382030206f626a0a5b203332205b20333438205d203434205b20333830205d203739205b20383530205d203834205b20
363832205d203937205b20363735205d2039390a5b20353933203731362036373820343335205d20313034205b2037313220333433205d
20313038205b203334332031303432203731322036383720373136205d203131340a5b2034393320353935203437382037313220363532
2039323420363435205d205d0a656e646f626a0a32342030206f626a0a3c3c202f4c656e677468312033353638202f4c656e6774682032
333735202f46696c746572202f466c6174654465636f6465203e3e0a73747265616d0a789cb5557b7054d519ffcefdceb977f7ee66dd84
1812086143d8405cc223482081d605f22408c13c4c02916cb2bbd9906413f20e3488259402521ee22a2916daa2456aedd63a35364e6b6b
1de960fee820e358ff5047663a9d818e9df1d1669ab3fdeedd9851a6ced43f3c67cfe3f79deff13b67bf7b0e300048a48e83b3b4a8b804
e6800ac0d2499a525ab9bdaa68606b3fe17b09bb4aab6a362db974df6b84ab0947b657adc8db332f7c95f02dc2b52d9dbe6e78533900a0
1411ee6819e8733d75219a4d384a3a1dc1eed6cebfedf9e79f2898e1ff78abafb71b34aac02f13b6b7760c070bb71fd209bf48365b4301
9f5fabfd6312182a00f92112d89f15af135e477871a8b36f68ce248408ef229cdad1d5e263d95046b88fb0a3d337d48dedaa95f0a8c13f
eceb0c2ccfdfdc43f812f9777677f5f6c5fe0ab466a578b0bcbb27d07de6b4b288f04de26007e36cec102f0a210461ca8ca6433a64002b
2ad95a0d163a3d2ab1d88c6e7cfd0c9c8564737d9fafc7d70ca3be9ece308c36f7f8da60b4c517eea53e14e8a17eb8a703465b035d346f
ed09b4c368c817269d50a09924edbeb00f463b7c5d2ea3ef230f9dbebe108c86db0d4957abaf13467bfac3a4d9170cb7521f32fcdfc168
861777b053b40710abc539c86719c6182bc6b721a8d0292b361511b94de1337b982d95c1623f498a61484d96c96c4ceb641f7e4907675a
bab973802242ccc41c96d2a8d2c8e80cc983c9c91c63bf8f9d98f5113fe7a5b3a76de8c49bc16668a6117796c556c0385ca3f607b8026f
286fc045584bf501f81d3baae4d2ca337004de110abc08e759014b6605b47a5d4d5687c5617199d677916d197979071e239f86a7717858
d9af544210de10933046b5cb947f04bf6587e0063c09d79432f8140e61359ca43a06bd1cc40da683543c70c9884415a099da02cc1537cc
fa113c0cfba11a2ea9e36a32bb6eb27e86bd465fcbdf89f375dc857bc9e23c5c66877916bfcccbe0649c2f36c149e5101be34d66dd4f07
3208e77913bba226c3eb0657925412d320bc426d1026d97a76188f12b3fd0603710326b52d7c459c9536826b683f40ed797801723142f6
e65ed4209c578214eb53623289459043de7a0192e9eb86bb7fa30a8e0a83652e67547197fba3de1d75aeabf599b9cbee802ea7e68a4265
3461d8351e8b55d6f1f9a23e2ad2a3e8b644b93beb83af5afc20775945659d2bfaebe2a219afc54d4524abaaa3a981484cf2e2a2dc784e
5840a1ec31326b2da5824ebbb441bd779545d8c0a609b0716a4829622b541bb8d200259606deab5b6cba5533b7623172d036dfeefcf8e6
9b89730b920a56c18a5bb73ebee5fcc74a8aa939c56d2dfea359fd2296a965e24c630d18cd9c4eab55fe9c39bd36a4fc2bf33fdf1237ae
c8922bd31f5ea1948c5da3349de44d606307bc65dc233853187a1485d93c1453f75875660c365db368568fc5a2098faaa9dc835cd3d442
5de38c5be0805074ab4533cedc6207d5ee7cab226aafae883aaa7756449d469758bdb3ee655063afaeab373640f437dccc9b5b10676f31
79c79be5f69db87ed10b16cea0d1fb960095695c881496c293c53d90c3b29525e816d96a8e25c79a432c0b611dcbe76bc41a75ad56a0af
b76d817256ca8b45b15aaa95eb5b6cb550cb6a44ad5aa73d680d42900595200644506db506f5a06d00fad900ef17fdeaa03e603b8a47f8
2362543da27dcf7a443f667b1c22ec497c829f1111f509fd6978daf612bea46eb172fa8ff4793c9b79788e96ade7b3757c9d96af97b072
5eae95e875ac8d85b4903ec887f447f951edb03ec69ee2e7b4883eae5dd51636b246f6c2a904d66855b2906559d9ea398c4f4e4bb6fefa
6de97f4f5922dfbdfe361be44dd3ef4ebfce5e9145ca362549b6b3c7cdfb47a93fb744ac1fdb7dd7864f60a1c5bc80fe72c0f1dee7e367
ef4faf4e28b11c373370f62ea5efaa532e004838fcd9fbff7e39a104fcf4767eb1a8dcf83e0df79422ca71ca5807ace5af18b9325bcae8
bd8adfa32a6e030fbd65768ae204aff1e4a1434998795b346830f29ed35bc6569a77ad31679042283e57c0c14a66e608d9f442c7e7fc0b
7301a96cdfcc5c85c5ec0c6ca63ba11b86a107daa095a2f7818beee016ba075c90072ba9aea6593369b86013e9f4d1edd047da01f04127
2c23693984497f39cd3642075517ddc39ffbea355180c600d90c50ef274dfdff889a3f1bb59a220d50ac3d6413266d83878f6cbe5ec422
9aed21bb5ae8278d16d2f599de02a685cfdc918bbc84a9ef269d66f2db467a2eb2efa2e83e73ed4e3f55a6975ed83ea3bf97a481afd071
dda1556b32ec25dc6546cd239eab61cd97ac3fb7cdbdc3965ec7d827d40e5056fcafa29a39a5d03fbd1476028c2b07bdb15f498cbaf197
79f87c047fe1c0e77a1de2b93cfcb9c42b6e7cd68197ddf8b3083e33854f4fe125893f2dc49f48fc711e5ebc50252e46f0c2fd1bc5852a
fc511e3e958ce723f8431dc7249e4bc22747f089098c483c4b1a6747f03189674e978a332378ba144f9d9c2f4e493c391f7f20f184c447
251e9778ec68863826f168067e3f0f8f481c4dc14312bf2bf1118907253e2cf180c4910ab718f1e37724ee4fc47dc313629fc4e1a14631
3c81c307f9d0a05b0c35e290970fba7140627f04fbfcd8ebc09ebd6ed1e3c7bddd4962af1bbb93b08b68754d61d81b93d829b143627b0a
ee692b147bfcd84631da0a31b4cd2642a9d81a7488d63c0c3a30e0473f99f923d822b1d96717cd127d766cda9d269afcb8fb21a7d89d86
0f39b151c75d3b13c42e893b13b0812c1a22585fe710f54bb1ce810f4e616dcd84a8955853dd286a26b0e620afae728bea46acf6f22a37
3e207147e572b14362e572dc4e24b627e3361bde4faceedf885b69d82ab1624ba2a870e396442c9758569a28ca2496266289c462894512
376f1a119b256e1ac18d12bd5378df147e7b0a37e46f121b24aebf8a85342bacc202e9edc67523b896603ecf15f99b708dc47b25ae2ec4
bc295c69c7151273252e93e8a165cf2abcc78939e8143959b8340397643bc4123f663bd0cd74e1cec3c5f654b17804b344a1c892b888d0
a209cc24fdccf9e85a6813aebb905efb57bd637ca10d33ac98e1e50b9c984eeae9119c1fc179696e31cf8f69a94922cd8da9493837c52d
e66ec41437de2d3159e29c294c4a4c13491213c96b621a3a25de25d1411e1c114ca080092368b7d9853d156d76d4255a68c9124195d455
898276210a9113e2b9884e7a5375a1a422d3917939a4231b67fec32798e79b2df00dfbffba65c17f019aa3ef870a656e6473747265616d
0a656e646f626a0a32332030206f626a0a3c3c202f4c656e677468203137202f46696c746572202f466c6174654465636f6465203e3e0a
73747265616d0a789c6360183480852eb6b002000293000a0a656e6473747265616d0a656e646f626a0a32362030206f626a0a3c3c202f
4c656e67746820323338202f46696c746572202f466c6174654465636f6465203e3e0a73747265616d0a789c5d503d6bc33010ddf52b6e
4c86a0c4a4b483319474f1d0a6c4ed1432c8d2c908e29390e5c1ff3efa685de881f4b877f78ec7e3a7f6ad2513807f7a2b3b0ca00d298f
939dbd44e87130c40e152823c34f977f390ac7781477cb14706c495b56d7c02f713805bfc0e655d91eb70c00f8d92bf48606d87c9fba42
75b373771c9102ec59d380421dcfbd0bf72146049ec5bb56c5b909cb2ecafe36be168750e5fe502c49ab707242a2173420abf7b11aa875
ac8621a97ff3aaa87abdae1f9fe27a816bc15ba29f5f329de05af096cefd0ad3e514c36a5bcede47c739ab6c359934846b9ccebaa44aef
010791783d0a656e6473747265616d0a656e646f626a0a32312030206f626a0a3c3c202f54797065202f466f6e74202f53756274797065
202f434944466f6e745479706532202f42617365466f6e74202f4743575844562b44656a61567553616e732d4f626c697175650a2f4349
4453797374656d496e666f203c3c202f5265676973747279202841646f626529202f4f72646572696e6720284964656e7469747929202f
537570706c656d656e742030203e3e0a2f466f6e7444657363726970746f7220323020302052202f5720323520302052202f434944546f
4749444d617020323320302052203e3e0a656e646f626a0a32322030206f626a0a3c3c202f54797065202f466f6e74202f537562747970
65202f5479706530202f42617365466f6e74202f4743575844562b44656a61567553616e732d4f626c697175650a2f456e636f64696e67
202f4964656e746974792d48202f44657363656e64616e74466f6e7473205b20323120302052205d202f546f556e69636f646520323620
302052203e3e0a656e646f626a0a32302030206f626a0a3c3c202f54797065202f466f6e7444657363726970746f72202f466f6e744e61
6d65202f4743575844562b44656a61567553616e732d4f626c69717565202f466c6167732039360a2f466f6e7442426f78205b202d3130
3136202d33353120313636302031303638205d202f417363656e7420393239202f44657363656e74202d323336202f4361704865696768
7420300a2f584865696768742030202f4974616c6963416e676c652030202f5374656d562030202f466f6e7446696c6532203234203020
52202f4d6178576964746820363332203e3e0a656e646f626a0a32352030206f626a0a5b203639205b20363332205d20313230205b2035
3932205d205d0a656e646f626a0a33312030206f626a0a3c3c202f4c656e67746831203134303838202f4c656e6774682039353234202f
46696c746572202f466c6174654465636f6465203e3e0a73747265616d0a789cd57a09745445d670d5bb6fe97d4b77b6ced249a713c296
d8210981002d4b0c8b1805228b68024958040204548c4c0087c4289820101611331a940062404613880c485c10187500e7e31307171497
88cc7c191748aaff5baf3b1074e69b9973fe73fef3f74bbdaa57efd6adbbdfaaca23941062c39b485cb78d1c9543c6934584d0bed81b75
5bde1d13dc7fed9f85cfb984807cdb8449c3eb7aedfc1a81f3f07dc7edb74eccbd3d75db5a42a49108d371c78414ef9c9f4b1f23441988
eff367ce2f5c08a743dec3e75a7c7f64e6034b5c644e34e2d3e278ca4a16ce9abf68c0037309d1737cbb6715962d240a5e44df89cf8659
f39695fc79c3b609d8d410e23a32bbb8b04833fd8d4a42fabe83ef33666387f159652721fdf09124cc9ebfe4a19085c9d5f89c80cf97e7
95ce2cec18f7ed6142fadbf1f9e1f9850f2d141b65e4af3fd2435c0b0ae717277d37a4159f9b909e330b4bcb96b4d55e5c4248ea6ff139
72e1e2e2858395bf22e82df5c8e36cc2656520819f8017900ceccb2483b1cddfe9487f924d849139e32612d3bcc2250b4838ca157f7e3f
21d75b1c92de5fbc7801d1f0965a44c4c06b0d11042f8714f2844a62248ec064fe91eabd56bddfafde5f0bd6b53d21fefd9fffb5607dbe
47dff95fd4ff21ce7f63d6f3ff1e04bfff6bd8ff78f6a33ddaa7fde7fda76fccc79f7bbc3dfe7f7bee5f5072fe17d41cfde7b0fff247d1
4eccc442acc446c2d0e2a2490c719104e221892409df5ab08faad6dafd138974bd2d737f532d91ffb468c37ab4718a1823f04ed591e8fc
ea180eabc1ab1b8ae0cca6eb7387103b21413fd84ef610bbea070f172e2e9c416a0b17cf5f406a672c2e9c436a67162e28c3fbece2c578
5fb6781ea99d555c8aed598b8bef27b5b30b1720ccece219d8737fe18242523bafb0d4c5efe84fbf9d5fb86436a95d703fef299d55389f
d42e5eba002197942c9885f7d91cff3ff1395556f3e6cc2abcc9ef4412f03baafab206b934201f51249ef442bf06945d1419a4d64e9285
b2882203f1ee44af1754b952f5ce79fe847c45f248b681281ff3a98455387d01c68db341a117dcacb6ee673a3bd89ef46fa8da84781ff8
d770646a0096d7d7db3d70dcd43fb5c7f3e21b6d0163f8b51a6c675f1fda0b25e3428951b4acf1e401b29f5c0eca3004df25abd2300564
4eb8cdf02cd24bbb14ef7df07d2fed5a7cdf1b5b445b47480fb894eb70fdafc3f553e116f58093885785e334a6a9709b11ee16b5f5f84d
7da9ead8476f1a9bf10fc6a6ff83b1037e3556205c4abd24ee0726fe56e2d19e8a265acbfd484a93b6a0e4620235fc999408361ca49701
34a22004ecabc72faf645411f1a11c97cb7666a75b95f9f4f39b602058a208b72a424ee213559f45d2a46a81eb4083f76494b01729ce22
43480ec925e3c81d641a292473c83cb2902c23e564b9ea0501c85484ccf805e44c845c40160720fd9ffb3ff29ff37fe83feb3fe3ffc07f
dcff8eff4dff317fb3ff15ff7e7f93ff65ff5eff4bfe3d3773f3bffe02b9b223f814a25212289cbfe460e112ea1b2c3c2ea562e19a96b1
0c08162d16aec1aca04c86040bd74c4eb098b1e4068b05cbb860b162b923585033c877a0707a0a838547ae9958e660e139775eb0701d2c
c0b2104b34166e3fcb823c94070b5a04591e2cc9349d349313781d25bbc836fa023e9560ff22eca917f693d56429f61ca32768b5d00ffb
5e2057c86984ac222760173ad618b4c413087f0e6dac834e2407104716b5d32c4546371b2f1e10ef129bc54be22992299689a7c402b18c
a6c17352bef402962c78136def388925cdf402292307e16b48835671a4682217e014ec225fe02c5cde27480d6940ea9b899d96920aa15c
b80b7bde964e91ad7895e2fb53743b3d8dd41da48f92b3643388422ed94ecf225f27c80fe451982854a01ad38412a4ff6dc4750ac76f25
659826ce521d61421fec43ea71ae19ea3d1afa4967d5eb0aa9c099279206b959b62b6e9c854bec057a8cb6cbeb493d390df7c022f888ae
16dde24e3197d404240005a406716fe563e412ba0c79e75739c72e3c2816d05de46bb1409981b8dfe41ce19c0784bb90a312d28ae541d9
823c0da6aba11a292d57f5794a1923a6e078c4a02c47ae0929857432175be5642f46b77e50476a1093caaf9c29fd8023b7899f22cf3574
adf003390523d1824bc4cb286b6e4218d55e536449048192be2e4b93e0195dd4e4bb73b2eb9d2971fdfafee2d165515c4d24afc9b8ccd5
ecf7e74d169dd2942629aa093c9a26d1e3fef49fbdfcb45fdfb179935d4d5da34606b18e2a18897d132663933f6137f68f1aa9bee39336
491efc1b5dd0e49a39dbf5b8e571f7a0c72dc583d435b3c0579c1823788469607f134a651be6f4589f59d94c3699144af4363144434c96
f3d91dedd6acac5b88a5b33d35242dd4e6b00b8a3bc3963e4068a8afdbb8b1ae7ef5a38fcab676967de11336e8db2fe89b972ed13710eb
76bf8d1e230c3d3bc26780ede45119441a41c2654474e664005f669a03dc21574e37acbc8bed6147a80fe9aaf27f2ed6a0ade8714de3f6
85c8f536526f58675b13ae8d32c74094c3198e233bda11c1c58e76cbe5541a2f582db634afcd6a1192bcc46a21ee787e179ed8f6cc33f8
f7cc33d7a896fd78ed1afb916aa53c768a9dc4728aa6e13580a6d5b33256c9aa58195d4b97d187e95a1e913fc5e03415f3b78ef87c8ee1
502f0af5d24a85d46b35b172149058aab79c19db649e38b905817d03a7b4b775224129edde8ef633eda9a89129f1f48019cca2303d33ce
2aa57bd2ac718e3846c7b02db4f85d3aa6b361975896db9c7bf5ec2e4480d6238e418ea3c8765f5244a413c2a3ac9248ac92240eb7fcce
bac1586f5f27621421169d407551611690a32d9d639b1c13c736854e9c36b6c93e711a5202fe2303a7b49d693f72c46acb0a52d3a152a3
58a4ef14e93bda1465b18665216d3eef24315fca571e161f961e70564528186322c44834b6a825e401796964997349d42a5219b12a7295
7355d44eb2d3699d4ea67b9089f40c923994a60f4874c7cb4afa509ae6151d7659910906b6a39de3508c6985b7bf5879dfe9871e3e33f9
2b6a1f352d8275ecdab5eb41ba6ed0fc4da31fac1b3ee2e42ddeafdeb867c7c268f62deabb08b96f40590b687f2b7dd1d4084602601c4e
40afd44b14566aa94147a2648d68406b1cdba447ae8d2ac306cef099ecb676af95f37bf14c76bbd7961560583c8e4c1fe7acf6d6e3d223
974cc110ff20799c28a1b40f49a47d20838ea77718ee30e6d312ba943e0caba91159d4d23848b3a639dc56b7352e1d64265096cece9e3d
de75afe4e9fc1c4e75a6ed64f5b4e018d25b83965aa75a6a0299ec4b089189b1d240ea43e5faa8d01d967a4375fcbaa8351e43bc362a22
26240ae2629d1e345dd4c945d5782f765eb45ce617aa2bd567c7dc414f09a7e09478423a212b0d31fb6384e9743a8d971df6d0346f068a
9c3afa5377bc005ceac461276e1737f4386fa8d0f0d8b3cf3e86856ac73d3dee9dd3e6c1fbefff944aecca67ac8b5da679d439ee69187c
f0b9df1d3af4bbe70e0acb9a1312d9dfd8f7774f67df7ffb25fb4635fd1974470cb77ddc558bb3511f3299e90b97ac02086015d112a57a
b21224a02225b262e93cd9a64685941e16860505ce6d72f2ebb861f02132dcd4a392ac9903a7f86c93052a43a49425e54ab3a08934c90a
ca3b24ce41dd346e271ce9faec34655d69d2d9fcab2ba53e7ca5f004caf70955be6e5c298ef079c251ba49727d4cbf7adbba983549cfa7
861b127a473912a2cc5a8c0b181ccc71ce544b675b7b475bbb2a58356eb5ab71c272392b2bb5a7303dfdd18a13d2bca1dc7c63a8c32eba
e313d20764847403a0810b4fd4eed8515bfbc20eb663d53ae2ffcb05b66ee553cfb31f7ffc91fdd890bbeed155ebd7af7a749df0e6d6aa
aaad4f57566dcd77ed5ff1cafbefbfb262bf2bfead9a735f7d75aee62d5ab864d5aa2558d06256224755c851b86a316e2536825692887a
dd0eb19e5487c6d65bd685aef1285151712131243e3ecaa81a0c92df1dedbe647fefb697d0b68837228f388f441d897e23a62d56d9656b
b57d6d03b498ccf418b40f5b88096d85a40f2069012b894fa4dd6ca10c3e1db76d2cdac9a0fdf33e61d7a8e5330ad4caf6b12fc66da343
83b6148b56428dd4967f0f357ffb250d5503e5b36c5a8cb0a9db92781e39470ba48fe039d49742a27c46b24d90b7891a898ac4adb17466
b77991f6333c2ca2aad50b973857192ff0dcce2b3cfee19e5b74a3c56949b2cf266f14858d64a5f812a2500097be3a1481f74c1b0f686a
bcdf6fd653caed268ec754abfb384ced5a28e47535bd2b9dddc57277756592004e6930e2d493be3ebb66a3f0924856ea6444280dd452c4
6ae058f10fd1665fc446eafe3c23478b81da9ac67d1f11bf2bfce5dd77bbe2116fd736a1e86a1fe1edaeac206e712ee29688c7675089a5
513088883cb3f13470b13dd5a74d55f29415b0421403481d7188eecf9df7496779cc174899ff73e95c8f2cb7d1463606b25c84390d221c
169ee53a7b6639aec43435bf250514aade21e93ceba270fe3ca5cc7f9e0ea20fb1c7d85bec4d56459749e35833fb827dc99a692e8da44e
9adbc0a6b1ed5c87b401fd1d3d1eb969c2ec3f59b6a3f606fbcc62a3bc5f6824fb50f830826878bc35f02c4778961bdb64c3d81ba2c65e
c2636fe749af9ac73b54e5ba316cc635d19d57ae30c457f373670df2598e16df0fd76a3ae221ad98dd62f5615a13690c935b4c565765ec
c1a81677b3754d98818441b851abd1c782c63e2a11719e3cd3eef506424cdbc58e4e14c25baa475bb3b8f52f488d4e8d498d4d75a5c6a5
c60f4bf245fb627cb13e972fce179f179d1793179be7ca8bcb8bcf4b5a98b43aba2aa62ab6ca5515b73abe36a93ee94a524cf7d0ee41dd
030a620a620b5c05710b6316c62e742d8c5b11b32276856b455c78cf383c84665adde9dcb912315aa4c5f54c85a1c2e10b7b56966e6969
6e1ed6fad89e135dd7a8f0e2a6825727161f9efa3f5784b492f21965e70e248feb5ab9aba4f0e873af1fb1553cd1bfffaea4a44e6e578b
fc9fc32594550419e673924afa9868aa343ea66bb18a2d6128a448c56624b9f6519196ce8bde6ecb601d972d7fbf9cead39b9d16e70a67
adb3de2921b16aa208129ce95023412053c0a5f1cfe4bdf2d65bafe43d33fef61dd3bbd887b41f95273d27a6efe9d3e7f353a73eefd367
5742021d4a4dd44607b9518348953815edc3c24f407c11912dc4646f91346b4ccd741384e17e54b8cd6ad38f8a56edc0cb1d0a2316c6df
b6cba9af16c4ac88a98f0124883b952a344135e350da23cac273cdcd83f63d72c24ffc271ed9d7f5f68b4f3db573e7534fbd08af0af7fe
dcbeb3a8908ea41abc461632c7894b974e6009d25581d2b2132759e84b200eaaadd43c26391aa9d462a087c25b6ccd8635514e87a07168
c858c1661e15a592d8a6ae8eb8f002e9b723104f9387452f8cae8f7e3ffa4ab4348c0ca3c384618e614ea9af92a249d1f6d59592525a2a
943a4a9ddae98bb880e378900dc8d615744d4515ba225674ee379c7a6deedb3366be7f3feb606fd3e4cecfa8d22cec786c6b8b49b877ea
e1b7070cd8dbbb2f1d487534848e601fb76d3ab0773bb780140ca83fa1ac43c8145f9464a1064da34cab70412eb7ea84104ca65a496334
ebc7d9f9f24fc71d531f704c93dae63e99dd8631b7cda63a0e0f6c96cbb826c2ccf7aacf91e7a8776082702091d1548d9da893346ec4c2
4f4d336fa729ec8396a6a6bdafcbf62d79b367d674a6c00735e30fede6b266f9e25494b51ef7bc637cee0843b4d6561912da6286964477
7352abb6c5fc7a64746204d1186e936d36d7a864350b07cca1ed62c020d8592ee92cb48ade2b7ad7f7e656d19d8d50926116015354929a
bbd0c782a682ebfad0b074dce8eed8b861c78e0d1b77343376b570cf9d776ebfebf707b2f63ff2c7cece3f3eb23fab5918f2cef9f3efbc
7dfefcb7ec33f67574cc2b7d7bbffe876933676040e42b96413366eee2f2c5842516a9f21de08b002d011395ab4cd666c3261d1534643c
8f4039769e1af80623259b2f24acb6b0acd4fd050ecaadd86d0d908c8d34abeaf16251f3238f6cdcd3d232fc95a547df121abaee11b63f
bbfd704357956cefda5e5cf43dcf914771f265382fcf917d7c26f9b0b88fb40a12d58824872749350f5dc42ce4d35bb43e6d9eb640bb50
2ba999834755f7d166fc8905d7ea65fbd79c8f1bf8e25f239b04aa2139a225b00549f5192d924fca930aa485d215490e204104b2fde7f6
eb32405d86908c9b65d0fa8f6570b15b06070a1cef39845f4ac1f12fa42016ece54208f8eb52d586c2d08642e4161b693134f38c6733df
0936c7a85feceb7cee6111e5a45cae502a3415da0a5d85bedc5061ac3055982b2c15d6725b7dc49508ebcdebe39bb67f651bf6ecdeb87e
cf9ef557a88d5dbef257f63db5c2854bc78f5ffaea9db7bfdec6de61edec3b74ce2cf4413be567bb07d1ca1b90421ee986fa9cdd91aed9
b486be0eadd118e56e53e35d0e8f755e6f80d68bddc1cea70d44bb4f6230e57bae8b0669113055f434755ad6d232685ff949e2f79f2cdf
270cc478f7222f3bbbf6caba5d4585ac95fd84576b21fdb63bdc05f40663903a2bc1cd82ac5788550f55a6666daba29335449363e32ea7
da11c6b8332779503b9017f26c08d758201bdc5057188c891ddd77db8b48c7c1d521fda3e080cd7ae270d77e5456c94c49c2d94a3117bd
8db325914bbe6ca34130e927c4c668b482a29b101b1b335ca78f89151d98a3aa457ba5a33a9ce7280fe6a85e313a7dac53217739352645
638f1fd58b5375a6fd22f7fcacaceea4f5779eb4b84da95b07d377b82b55d4fb94f8fd24093dcd373f4a17a58f32f4c7d0db57dfd73058
3b5837583fd8a07711174d107ae97ae97b87a4d8531cbd437bc5f48a4d7625c7252455ea2af595864a233feba38220eb643d18c0082630
83052220129c1025466b9352928725df975c91bc22b936b93ef94a72386e4416ddc899b1946f0764f7f57084a69542f976210365074f8c
df39b5ba7ac686616d3b7efcafa9c7e695bc55b86a4df16edfeecd9ffcb1e480386c6faf5e1327fa46c7997a6fa9def6aadb7d383d7dca
9d63f33ce6848dabb6ef51f758b84415fe266d471fc48c6a923466682456daaaa9d2e951c66863169b89fba01accbd6a2c6fef086c7431
16bd1c88453c82db4307f3789e98ce23b9953e48cbd9eab165afbf7ef6b9aa2a693b7ba3a6abbe7afcd667ff2414d4d0a13c12ed452f9c
ac7abf1dd77b5137fc7f8d8eb6da9b0de8fd76fd788c03390eee8e59018bbae8bd1e044a1d47781008c1f81770bbeb993d91eee541e0a5
e6e611fb961e7d87be470f0a2f74153efbece106a1fc5afd9e9299576027e77e0846a00ab1007798d77c49eae652b05241e215083291a9
951079b800e40f922c8140259128fcbc45d763256a9fc8cf3faeaf44dbc202271e593736a39a605137a5be277385b942b95021540a2b84
754283a0e11369418b56eca09110292692449a0cc9a24b934ed2e9201824a66a72480e1d0da3c51c2957f669f2493e9d0253c43c4d0929
a173608e384b9a2d1768969225b41ccac5a5d2c3f26ab29a5643b5582d55ca75a48e6e12b6c26671b3b449de29bd2837698e682e68fc9a
a17ce79ba6a569d43de418bd97de7b8cdd73552ce89c087baed5a384f25100e9282103fdd6375a9a244b5a9d569ca4d3c2249d5e274ca2
825eaf93ad8a4691507a1a8dc205674068149c61b84e124046e1e93506bd4eab091c4b62bc30066438b6c9cacf8c6cfc16c26ffa8018b9
5c295feeabc72b639b2c37a47b52952e66296fd83f14f0f53ad0e4f2f64b42a8102ac5ebd275a385d1528ecea79b264c9326e9f2740b84
0552896e19ea6299542155095b84cdd2065dabd02afd51781bde93a225410bb2a897741abd162b83438880503152726a9c5abbde61f010
0f750b4910277aa478395ef16892b409ba38bddb900519628626cb906aca11722147f489c3259fec537c9a91da91ba917a9fc967e25acc
17f2c43ba5bbe4bb943ccd04ed44dd24fd4c52448b85b9508c3bbbb9f25c6581b6503fcb506a5a4a96d265c27278485c8edaad901f562a
948734cbb415da72dd03fae5862ae131a9c6b4896ca21b84f5b04d7c5ada2c6f56b6687c297586674d2f90176883d000bbc5dd52a3dca8
ecd634185e36fd5ed807af8b87a466ed1f4c6dc2313829be2b2d532dc249f91f75eba93bbff9cb2fce7df94533fbe8dc5fff760e6da30e
e6f272ad1eea3ae7a28d0c462f5a8636a2a7237c3992555664d10aa2c22b49a40205ab806ab722a4ceaad5515ee97568325a2b1acc709d
225251831e26045be810866e033107f5af9a0a1a81dced73bc1d38656cb376dbc43f33895ffbe0669d28ea2245872e513744bc453749bc
5b99ac2bd13d401f161f5096e8d68aab745bc467c54dca53ba5add0bb4517c59dca13cafabd745e94094d007f491e0901cda487d32244a
1e6d6fbdcb38886641a63440c9d066e9538da321471aa51da3f719a7705f15a6c0dd52be3c45c9d7e46ba7e8f38ca5c6876885f169ba41
d94d1b9426e37bc60b46bf31859ffb096e2dc53f744ab188dd4f779d6307d9c173f415b6f81c4da6c96241d785aea3b499e50a638450b6
88d6a8910c570c3c9299e913be118a46d05a89998b9910b3c96a2666a3d56024bc3219d1710d5674dbe146bdd642f45215bc6ed2b75a4c
46834e8bdeaa318b66bda55b011a55ecfa1e6257d7f66d01a95b301cb70777c2ffcc1ba5efc2bc5ce657642269642d18437561468bd16d
4c378ed6dda11b6f9caa9daa9babab32ae30ae37da740489404fd39bf4e630ea102ca2450ad3d9f57643a429d29c441230e7ba449794ac
e9a5f5e812f4098624636f536fb3cb9a89b1325d481553a581ba0c7d8661a031cb94654eb5de4a7cd427f8c027fa821e385c3b4a779b71
b469b4d9679d48eea4770a93204fcc43fd4c42fddcadbd1bbd7092618a698a39cf5a424b84d9ba39a639e6026bb9e621d343e66af2b876
b57eb5a1da586daa366fd16ed46f346c356d3537e81b0cbb4dbbcd4dd6f7ac17ac7e6b31ea5232d13e341d73d330caf59926ac1fbfe191
f5f3c64d4c8b638303e176f63b0f6fcdad9c288eefdc00f3509394b0fd6218aea813c883be011166875ef668232d8e68bde48a03a26fd4
92467a44eb680cd9e73168755242680489d64921829db82246e8cc9287ab4e3d2a91b8d6ba13372e7bf8cac76a0bfe07a58d2f6d318f62
f7f597a9743fa69ee912cfe343709d18a01d5b1983837b35f580329a6f99c5b0a13f37ce7b72c8909afb1b7f1e3a726dfeb405a553f3d7
1eaeddf0f1f79b96d4946dbcf2f1fa9ac96b7f7ae6c908e793db7e5a3bf966deaa7c19b11e67a8c1ac71ea1d11665172018968749246f7
11a7b9d1bacf131511e930535c80443adc369144c63a46a0b18a9eebe957e47618dc6f8a014ebdc86a90c97fc8abda19e4b79bd7208399
fde9758e4339833738ee8374e7ac9934b5b474eaa43539c37e7e71de9343873e39efc59f871dceaff969db93ce88279ff9e9c9bb6bd67f
7c656359cd924ddf7fbc817b64535727e4294fa1ef109a6ec1c497660bb1591dc2da8aa9acb3fe8583f55d9ddf408860c3a06e625d5def
b398ce39dfe06a71ad3aae9aff4f9aef6f5221ef9852fd4399fa454213b3abef38ce3407220537472aac3d565ffb7efd89f5154af5379d
a7ba3e623676883e4ddb69234cbb41cb7cfebf6deae6cb355a764ca867eeae4eeda26381f7b40cf102217c8fd6744c9d92f498918f54d7
79b48c7edc55a054ff38fe18ff5c820acf777d234d46dc4615f780cc8c34048b5764dce01f2c5d769086b04bc523b3babe51965e58b7bd
a9e4ee2dfc641fe999ae5492701c131a161a81032d49899ea4f434576646044db35a1459d830a877f1c10563d9b5214328dcbee060712f
e1a86bc87a36a2eb2bc1405fbbef56911e1a3593be2218babe62b99bb3392dcc8eb454aab458ed61a12a2949b82685e959238b6904fbae
b9bc54a9fc69c3d6fc92a6ed359f2a45010e3a909a359c7b24240de93ed8d5a17e89b1b3eb1b7c13e0cd3a2029d11d971ee70d0b755885
e7111fbb44430e2e2b45d6966fb91bf1adbbf0d313017c9d2a159cbbc4a444959fb0d0b83487222ab2ca6b6606cccbee85dcdd4e61c810
766d2c72d75b381a97bd99360be15d1decf699a30e51f1d6fbd8e8ae0e219cbebe7e08a70675313dc81dae791559252733235d78fe6069
7933fb8e46a0a491bd6d9fd6a0a4f3b706ce71a45daac7e5fa4212d5631b435cb831466335c459ece33ceab9380f11966c7e5ad3760bf1
59b5466ba34d88ac22e19be4585babde9c92fda5d7cbb22f7b2d97bdded49b8e6d6e1cdda80ea3f0177c7722edea3ec7618a7a94b3b769
665222fdf9a6339dee739d2dbd7acd9e1938df49e1ff81447a1d2482dce5b33a73489826d46c17351a08d5c9e3226fd0cbb271d7ebb369
70c362a932851f0edd67daa425ad12e5d45e66eaffb3bca93ebddf59efac75ae705ad4e3c85f938c14a32fb9c531014a5f7aa58553fe73
4b0b3fe5e8a6f1b59739d174ffd7c1b3b1208d993e73580ed5883a834663116da671a19cbe00799c3aa319e3b558453659b5ad468113c6
54aaa82aba5f1e7e613e272decb51be75ffca0442affc50918ce2e77e2ecc9a41f99ea0b4fc909eba3e96d713a3491bdb52456d624c468
e313c7f5bf2128352cf2a31d2428cc19eb6e4cb0d22ad2ef70ef7d16b2295449688d888e4bc9cebee8f5f26c6e69f7e25f40cb416d6666
645e1755b7ce7b1cd949283d7e6cc7d57b7754d2f84751d5b70bfbb93883ba079426aa3da0e5495121495c98ddc2ed664de896adca5d22
19ef0bed9563d45842c3ed1a8b96ffdb35cea98d758f4beac199ca986a06e151aec638ab506548dce450e25acd913101963ab27fcd0f46
a89bc5ff8bf3c7809df6d445908feb3ceceea997ebbae15f250953b6bcf774d58cfbccd97f27b1814f1a3ff88de96277fde3879de34c53
b4fc3bcfeeef1dd5efae94f92c9a1013fbf1c3ab779aa6fceafba6c1e229f51b1e22ecc2824106624903d6dba5305285e5532c75588ab0
d460d989e5092c2be503e49cb8951c9745725c5c8af5385286639bc44ba45c2c218b243b5924b69345c2072485b7251b39286491a3bcf0
b6f8850a7310c660bb0f290537c9c4febd622b19223f41f2b11eccdb521825bc60d26a1289b03650df2832119e9709d919a87bb691d647
716e5da028ad68012937f1de87dc46aac9113a8c6e1314619cf07b88051fcc85cd7010fc627fb14d324bd3a43765abbc54de279f937f54
d294df2aa7353a4d8be607ed40ed71dd47fa670cab0c578c25c655c65ae31e63abf1b8f1bf8d1d269769b029d734d354636a3453f304f3
53166a2908e88364c3449c793631a0c759c816ae2dd1218462cdbf5d52c854febd8aa845e054f55b32dea624149f026d8168684eb00d3d
fac51e6d8984d3f1c1b64cecb4848c20a5ea57738bc91c320b675f425ca4179989beee225e928a571ab66620848b0c479825a40ccb6252
4c0ac97cd2177b47930508df1f5bb7927978b9c85dd77195a94fc55817e39807f05e8490ba7f63d68cebb34ec4991ec0b9f8174a0b109a
d3518863feb31947626b2e8ecb274b116226c216aad88ad511852a472ec4b200ef0b116606e29d83702e1c5f8ab317aaef7e8967828aa5
0c292ac5eb7eece5b396216ca98ac98b73a791f49b46758f1102c6e6ff8dfaede9af7fd9aa5df02f980dc4ac9e59394818e6a848e22451
ead78d7dd06ebd642862cf418bcd454d8c2163d5af1cef44194c209370bebbc9643205679846ee21f79202be37a62295a84c15aaa15aaa
a37a6aa0466aa26672841ca5166a55962e98e34dcb1a15ac7302f5ada9c1fa9660ed0dd669c13a33580f0cd659c1fad6603d3c588f08d6
23799d333c353558a705eb0184340b2b7cfe6b0caedae1670ffce4851febe00713fc9d410783fff1c0df4cf0d73ab8e281ef1fbf55fa9e
c1e53af8ae0edaafc2b757e11b065f0f82af86c325065f7ae18b8b13a42feae022025e9c009f7f96227d7e153e4b814f197cc2e08217fe
62878febe03c838f6cf0dfcbe1dc21f82f061f22f887cbe1ec99dba4b3cbe1cc6d70fa4f4ee934833f39e10306ef33788fc11f199caa83
932762a4930c4ec4c0bb5e38cee0add556e9ad28783314da181c63f00683a30c8e30f80383c30c5e67d0cae01083835668a9f4482d0c9a
5f3b24353378edd5e9d26b87e0b515e2abbff748af4ef7f9e1559ff87b0f1c60f04a1dec67b08f4113839719ec2d82974cb067b747da53
04bb77d9a4dd1ed865834624baf12aec64f022831718ecb0410383e79f3349cf7be13913fcae08ea11a4be0e9e65b0fd1983b49dc13306
d8f67484b4ad089ede6a919e8e80ad16d8a283cd0c36d519a54d0cea8cb011076dac830deb4dd2865eb0de044f5d8575b587a4750c6a6b
a64bb587a0768558f3a447aa990e353ef1490fac65b0e689fed21a064ff487c791cdc76f85eac7f452b51d1ed3431576541541254aaad2
03abadf05b068faeb24a8f32586585950c5630a860e0f3ff66f972e9370c962f87478aa07ca2432af7c0c30c963178c8040f1ae0011d2c
65b0e42a945d85c55761d15558c8a094c10206f3e2e07e0673adc3a5b913600e83d9cb61163e9430286650c4602683190c0a0741c155b8
d700d3194c633095c194c93a69ca5598ac83bb4323a4bbbd90cf6012ce3c69384c74c0046a912684c35d76b8734c887427833c3ddcc160
fced16693c83db2d308ec1587c3396c198d116694c088c8e364aa32d906b84db18e4d4c1a83a18c96084d04f1a7115861f825bc7828fc1
30064387d8a4a17618926d9686d8207bb051caf6f9cd30d8088318643118986997065e85cc0c8b9469878c74bd946181743d0c88813423
786fd14b5e06b7e82135452fa51a21450ffdfb69a5fe16e8a785be5ee8d3db23f52982dec936a9b707926dd02bc923f5ba15923c90e8d1
4b8966f0e82181819b41bc19e290cf381bb88a20f62ac4200b3145106d842894601403e755881c0e11f810c120bc08703725853108c541
a111e060606710c2c08600360656e4d53a1c2ccbc15c0426064643a864646040684328e819e82ca065a041300d03c50e721188f852440b
7000f602c39dac4512fa01b50061409b69d1eab5b4cfff0f3ff2ff9a80fff517fd7f00cdb4aec50a656e6473747265616d0a656e646f62
6a0a33302030206f626a0a3c3c202f4c656e67746820313031202f46696c746572202f466c6174654465636f6465203e3e0a7374726561
6d0a789ca58dd90ac2500c440f74515bede2d6d5dac52efeff0f76b8dc07a1f8d4c0c94c2684c0ce7236896bba67279f0347e94904849c
b9d84db4b98c45627d6afaf5cfd7db8fbf8b074f32720a4a2a93d6e245c35bdad2d133f061646256b288ef0a9dc0035b0a656e64737472
65616d0a656e646f626a0a33332030206f626a0a3c3c202f4c656e67746820333938202f46696c746572202f466c6174654465636f6465
203e3e0a73747265616d0a789c5d933f6fc23010c5f77c0a8fed800221d8444248155d18fa47a59da20ec13ea348c5894c18f8f6b5fdce
546aa4f0f47ebe73eeccb9dced9ff7ae9f44f9ee077da049d8de194f97e1ea3589239d7a572c2a617a3db14bbffadc8d4519920fb7cb44
e7bdb343b1d988f2232c5e267f130f4f6638d263218428dfbc21dfbb9378f8da1d800ed771fca133b949cc8bed5618b261bb976e7cedce
24ca943cdb9bb0de4fb75948fb8bf8bc8d24aae41728490f862e63a7c977ee44c5661e9eadd8d8f06c0b72e6dffaa241dad1dee3ab180f
69a1df0937c00de386b106d68c35f0129b2c2be06c179922a886ad39a8662c81f94bd92ac83a2f22d6c01a8e35c03572a2b45060244769
a1c0e8a1e61e6aeea1b6c096b1055e1d138ed2421396e82d4a0b055e027333d9a273b982a05a890ee53a6760031427b938c9c5494a58f1
bfc456da4c5390aa60f9e3d9a206851a146a50a841a106b5ce89d8072d2b6e5971cb0a87aff8f0553cfc30627996e2b4c5ab711f657df5
3e4c71ba3f697ce3e0f68eee576c1cc69815df5f37f7eab70a656e6473747265616d0a656e646f626a0a32382030206f626a0a3c3c202f
54797065202f466f6e74202f53756274797065202f434944466f6e745479706532202f42617365466f6e74202f424d515144562b44656a
61567553616e730a2f43494453797374656d496e666f203c3c202f5265676973747279202841646f626529202f4f72646572696e672028
4964656e7469747929202f537570706c656d656e742030203e3e0a2f466f6e7444657363726970746f7220323720302052202f57203332
20302052202f434944546f4749444d617020333020302052203e3e0a656e646f626a0a32392030206f626a0a3c3c202f54797065202f46
6f6e74202f53756274797065202f5479706530202f42617365466f6e74202f424d515144562b44656a61567553616e730a2f456e636f64
696e67202f4964656e746974792d48202f44657363656e64616e74466f6e7473205b20323820302052205d202f546f556e69636f646520
333320302052203e3e0a656e646f626a0a32372030206f626a0a3c3c202f54797065202f466f6e7444657363726970746f72202f466f6e
744e616d65202f424d515144562b44656a61567553616e73202f466c6167732033320a2f466f6e7442426f78205b202d31303231202d34
363320313739342031323333205d202f417363656e7420393239202f44657363656e74202d323336202f43617048656967687420300a2f
584865696768742030202f4974616c6963416e676c652030202f5374656d562030202f466f6e7446696c653220333120302052202f4d61
78576964746820383338203e3e0a656e646f626a0a33322030206f626a0a5b203332205b20333138205d203431205b20333930205d2034
34205b20333138205d203438205b203633362036333620363336205d203532205b20363336205d2035340a5b2036333620363336203633
3620363336205d203631205b20383338205d203730205b20353735205d203732205b20373532205d203736205b20353537205d20373920
5b20373837205d0a3931205b20333930205d203937205b20363133205d203939205b203535302036333520363135203335322036333520
36333420323738205d20313038205b20323738205d203131300a5b203633342036313220363335205d20313134205b2034313120353231
203339322036333420353932203831382035393220353932205d20313233205b20363336205d203132350a5b20363336205d205d0a656e
646f626a0a332030206f626a0a3c3c202f463320313520302052202f463220323220302052202f463120323920302052203e3e0a656e64
6f626a0a342030206f626a0a3c3c202f4131203c3c202f54797065202f457874475374617465202f43412030202f63612031203e3e0a2f
4132203c3c202f54797065202f457874475374617465202f43412031202f63612031203e3e203e3e0a656e646f626a0a352030206f626a
0a3c3c203e3e0a656e646f626a0a362030206f626a0a3c3c203e3e0a656e646f626a0a372030206f626a0a3c3c203e3e0a656e646f626a
0a322030206f626a0a3c3c202f54797065202f5061676573202f4b696473205b20313120302052205d202f436f756e742031203e3e0a65
6e646f626a0a33342030206f626a0a3c3c202f43726561746f722028526570726f64756369626c65207265736561726368206669677572
6573290a2f50726f647563657220284d6174706c6f746c696220706466206261636b656e642076332e31302e3129203e3e0a656e646f62
6a0a787265660a302033350a303030303030303030302036353533352066200a30303030303030303136203030303030206e200a303030
30303236373836203030303030206e200a30303030303236353730203030303030206e200a30303030303236363234203030303030206e
200a30303030303236373233203030303030206e200a30303030303236373434203030303030206e200a30303030303236373635203030
303030206e200a30303030303030303635203030303030206e200a30303030303030333330203030303030206e200a3030303030303338
3130203030303030206e200a30303030303030323038203030303030206e200a30303030303033373839203030303030206e200a303030
30303131343931203030303030206e200a30303030303131313231203030303030206e200a30303030303131333339203030303030206e
200a30303030303130353934203030303030206e200a30303030303033383330203030303030206e200a30303030303131373231203030
303030206e200a30303030303130373332203030303030206e200a30303030303135313339203030303030206e200a3030303030313437
3633203030303030206e200a30303030303134393834203030303030206e200a30303030303134333633203030303030206e200a303030
30303131393030203030303030206e200a30303030303135333731203030303030206e200a30303030303134343532203030303030206e
200a30303030303236303332203030303030206e200a30303030303235363732203030303030206e200a30303030303235383835203030
303030206e200a30303030303235303237203030303030206e200a30303030303135343134203030303030206e200a3030303030323632
3536203030303030206e200a30303030303235323031203030303030206e200a30303030303236383436203030303030206e200a747261
696c65720a3c3c202f53697a65203335202f526f6f74203120302052202f496e666f20333420302052203e3e0a7374617274787265660a
32363935320a2525454f460a
END FIGURE
